% archon:covers MR0555258CompactifyingPicard/Basic.lean

\chapter{Base-change theory (\S1)}

% This chapter formalizes Section 1 ("Some Base-Change Theory") of
% Altman--Kleiman, "Compactifying the Picard Scheme", Adv. Math. 35 (1980),
% pp. 54--63. Throughout, $f : X \to S$ is a finitely presented morphism of
% schemes (proper where stated); $I$ and $F$ are locally finitely presented
% $\Ox$-modules; $F$ is $S$-flat where stated. We write $\sExt^q_X(I,F)$ for the
% relative (local) Ext sheaves, $X(s)$ for the fibre over $s \in S$, $k(s)$ for
% the residue field, and $(1 \times g) : X_T \to X$ for the base change of $X$
% along $g : T \to S$.

%==============================================================================
% External dependency anchors (leaves of the dependency graph).
% These are results the paper CITES; they are not proved here.
%==============================================================================

\begin{theorem}[Existence of the representing module $H$ (EGA)]
  \label{thm:ega_H_existence}
  \lean{MR0555258CompactifyingPicard.External.H_existence}
  % SOURCE: [Altman--Kleiman], §1, (1.1), p. 54 (read from references/MR0555258-compactifying-picard.pdf)
  % SOURCE QUOTE: "(The representability results from [EGA III_2, 7.7.8, 7.7.9]
  % in case f is locally projective, and its compatibility with locally
  % Noetherian base-changes is proved in [EGA III_2, 7.7.9]. See [ASDS, (12)]
  % for a proof that the formation of Q(F) = H(O_X, F) commutes with arbitrary
  % base-change; the proof for H(I,F) is analogous.)"
  \textit{Source: as invoked in Altman--Kleiman, §1, ref [EGA III$_2$, 7.7.8].}
  Let $f : X \to S$ be a finitely presented, proper morphism of schemes, and let
  $I$, $F$ be locally finitely presented $\Ox$-modules with $F$ flat over $S$.
  Then over a Noetherian base the functor $M \mapsto \Hom_X(I, F \tensor_S M)$ on
  quasi-coherent $\Os$-modules is representable by a locally finitely presented
  $\Os$-module, and the representing object's formation commutes with locally
  Noetherian base change.
  \begin{proof}
    Established in [EGA III$_2$, 7.7.8/7.7.9] (and [ASDS, (12)] for the
    base-change compatibility); used here as an external input.
  \end{proof}
\end{theorem}

\begin{theorem}[Acyclicity of quasi-coherent sheaves on affines (EGA)]
  \label{thm:ega_ext_affine_acyclic}
  \lean{MR0555258CompactifyingPicard.External.ext_affine_acyclic}
  % SOURCE: [Altman--Kleiman], §1, (1.2), p. 55 (read from references/MR0555258-compactifying-picard.pdf)
  % SOURCE QUOTE: "Since g^{-1}U is affine and F is quasi-coherent, the
  % right-hand side is equal to zero for q > 0."
  \textit{Source: as invoked in Altman--Kleiman, §1, in the proof of (1.2).}
  For an affine scheme $U$ and a quasi-coherent $\mathcal{O}_U$-module $F$ one has
  $H^q(U, F) = 0$ for all $q > 0$.
  \begin{proof}
    Established in [Serre / EGA] (vanishing of higher cohomology of
    quasi-coherent sheaves on affines); used here as an external input.
  \end{proof}
\end{theorem}

\begin{theorem}[Descent of finitely presented data to a Noetherian base (EGA)]
  \label{thm:ega_descent_noetherian}
  \lean{MR0555258CompactifyingPicard.External.descent_noetherian}
  % SOURCE: [Altman--Kleiman], §1, (1.1) p. 54 and (1.4) proof p. 57 (read from references/MR0555258-compactifying-picard.pdf)
  % SOURCE QUOTE: "By [EGA IV_3, Sect. 8], there exists a Noetherian affine
  % scheme such that all the data descend to S_0."
  \textit{Source: as invoked in Altman--Kleiman, §1, ref [EGA IV$_3$, Sect. 8].}
  % NOTE (iter-016): re-signed to the DATA-descent form actually consumed by the
  % 1.4 free-resolution step. The earlier Lean signature recorded only the
  % reduction-to-Noetherian conclusion (a base S₀ + a morphism S → S₀); the
  % source quote ("all the data descend to S₀") covers the full data, so this
  % strengthening is FAITHFUL, not a new axiom. To keep the affine bridge directly
  % applicable, the descended base is delivered as a literal Spec of a Noetherian
  % ring A₀ (EGA descent produces a finitely generated ℤ-algebra), so X₀ = Spec A₀
  % and no scheme→Spec transport is needed downstream.
  Let $S$ be an affine scheme, $f : X \to S$ a finitely presented morphism with
  $X$ affine, and $I$ an $S$-flat, finitely presented $\Ox$-module. Then the data
  descend to a Noetherian base: there exist a Noetherian commutative ring $A_0$,
  a morphism $p : X \to \Spec A_0$, and a finitely presented
  $\mathcal{O}_{\Spec A_0}$-module $I_0$ with $I \cong p^* I_0$ (pullback of
  $I_0$ along $p$). In particular one may reduce statements about $(X,I)$ to the
  Noetherian affine $X_0 = \Spec A_0$, where $I_0$ is finitely presented.
  \begin{proof}
    Established in [EGA IV$_3$, §8] (and 8.2, 8.5.2, 8.10.5, 11.2.6 for the
    individual descent statements): the morphism $f$, the module $I$, and the
    flatness/finite-presentation data all descend to a finitely generated
    $\mathbf{Z}$-subalgebra $A_0$ of the global sections of $S$, which is
    Noetherian; $X_0 = \Spec A_0$ carries the descended fp module $I_0$ and
    $X = X \times_{S} S$ is the base change recovering $I = p^* I_0$. Used here
    as an external input.
  \end{proof}
\end{theorem}

\begin{theorem}[Flat pullback preserves short exactness (standard)]
  \label{thm:flat_pullback_exact}
  \lean{MR0555258CompactifyingPicard.External.flat_pullback_exact}
  % NOTE (iter-016): standard commutative algebra / sheaf theory — pullback
  % (inverse image with the tensor) of modules is right exact always, and is
  % LEFT exact (hence preserves short-exact sequences) when the term being killed
  % is flat over the base. Mathlib HAS this at the module level
  % (Module.Flat.lTensor_shortComplex_exact / lTensor_exact) but NOT for
  % Scheme.Modules.pullback of a ShortComplex of sheaves of modules
  % (verified absent via leansearch/loogle, iter-015). Anchored as the genuine
  % EGA/standard input used in the last step of the 1.4 proof ("since I is flat,
  % the pullback of the sequence ... is the desired sequence"); flagged buildable
  % from the module-level lemma + stalkwise flatness, to build/upstream later.
  % SOURCE: [Altman--Kleiman], §1, (1.4) proof p. 57 (read from references/MR0555258-compactifying-picard.pdf)
  % SOURCE QUOTE: "Since I is flat, the pullback of the sequence on X_0 is the
  % desired sequence on X."
  \textit{Source: as invoked in Altman--Kleiman, §1, (1.4) proof p. 57.}
  Let $p : X \to X_0$ be a morphism of schemes and let
  $0 \to K_0 \to J_0 \to I_0 \to 0$ be a short exact sequence of
  $\mathcal{O}_{X_0}$-modules such that the pullback $p^* I_0$ is flat over the
  relevant base. Then the pulled-back sequence
  $0 \to p^* K_0 \to p^* J_0 \to p^* I_0 \to 0$ is again short exact.
  \begin{proof}
    Pullback is right exact (it is a left adjoint composed with the tensor with
    $\mathcal{O}_X$), so the surjection and the middle exactness are preserved;
    flatness of the right-hand term forces the left-hand map to remain a
    monomorphism. Established in standard commutative algebra / EGA (the module
    case is \texttt{Module.Flat.lTensor\_shortComplex\_exact}); used here as an
    external input at the sheaf level.
  \end{proof}
\end{theorem}

\begin{theorem}[Flat base-change tensor preserves short exactness (standard)]
  \label{thm:flat_tensor_exact}
  \lean{MR0555258CompactifyingPicard.External.flat_tensor_exact}
  \uses{def:tensorBaseChangeFunctor, def:isSFlat}
  % NOTE (iter-024): standard commutative algebra / EGA flatness input, the
  % SES-preservation implicit in Altman--Kleiman's assertion that the functor
  % T(M) = Ext^1_X(I, F ⊗_S M~) is half-exact. F being S-flat means F_x is flat
  % over O_{S,f(x)} stalkwise, so the base-change tensor F ⊗_S (-) =
  % (Scheme.Modules.pullback f) ⋙ (tensorLeft F) ⋙ sheafification carries a short
  % exact sequence of O_S-modules to a short exact sequence of O_X-modules. Mathlib
  % HAS the module-level fact (Module.Flat.lTensor_shortComplex_exact) but NOT for
  % the sheaf-of-modules base-change tensor through sheafification (verified absent
  % via leansearch/loogle, iter-023; a genuine from-scratch stalkwise-flatness gap
  % parallel to thm:flat_pullback_exact). Anchored as the genuine standard input
  % that discharges the hypothesis `hSE` of lem:ext_tensor_halfexact, making the
  % half-exactness unconditional; flagged buildable from the module-level lemma +
  % the stalk/IsLocalizedModule layer (as in lem:tilde_exact), to build/upstream later.
  % SOURCE: [Altman--Kleiman], §1, (1.3) proof, p. 56 (read from references/MR0555258-compactifying-picard.pdf)
  % SOURCE QUOTE: "Consider the functor, T(M) = Ext^1_X(I, F (x)_S M~). from the
  % category of finitely generated A-modules M to itself. ... Since T is half-exact
  % and since T(k(s)) is equal to zero, T(M) is equal to zero for every finitely
  % generated A-module M [OB, 2.1] or [EGA III_2, 7.5.3])."
  \textit{Source: the SES-preservation implicit in Altman--Kleiman, §1, (1.3) proof
  p. 56 (``$T$ is half-exact''); standard flatness.}
  Let $f : X \to S$ and let $F$ be an $\Ox$-module that is $S$-flat
  (\cref{def:isSFlat}). Then for every short exact sequence
  $0 \to N' \to N \to N'' \to 0$ of $\Os$-modules, the image under the base-change
  tensor functor $F \tensor_S (-) = \mathtt{tensorBaseChangeFunctor}\,f\,F$
  (\cref{def:tensorBaseChangeFunctor}) is again short exact:
  $0 \to F\tensor_S N' \to F\tensor_S N \to F\tensor_S N'' \to 0$.
  \begin{proof}
    Flatness is checked stalkwise: at $x \in X$ the stalk $F_x$ is flat over
    $\mathcal{O}_{S,f(x)}$, so $M \mapsto F_x \tensor_{\mathcal{O}_{S,f(x)}} M$ is
    exact (the module-level fact \texttt{Module.Flat.lTensor\_shortComplex\_exact}).
    Reflecting this exactness through the stalk / \texttt{IsLocalizedModule} layer
    and sheafification (as in \cref{lem:tilde_exact}) gives the short exactness of
    the base-change tensor at the sheaf level. Used here as an external standard
    input.
  \end{proof}
\end{theorem}

\begin{theorem}[Flat restriction to the fibre preserves short exactness (standard)]
  \label{thm:ega_flat_fibre_restriction}
  \lean{MR0555258CompactifyingPicard.External.flat_fibre_restriction_exact}
  \uses{def:isSFlat, def:fiberExtVanishes}
  % NOTE (iter-029): the correct flatness mechanism for the (1.3) fibre diagram,
  % replacing the iter-028 named gap `hpullflat`. Altman--Kleiman restrict the
  % acyclic sequence 0→K→J→I→0 to the closed fibre X(s); the right-hand term I is
  % flat over the BASE S (not over the fibre), which by base change along S → k(s)
  % (equivalently Tor_1^{O_S}(I, k(s)) = 0) keeps the restricted sequence short
  % exact. This is a DIFFERENT, correct, anchor from thm:flat_pullback_exact, whose
  % flatness slot is on the pulled-back term: the prover (iter-028) showed deriving
  % `IsSFlat (j≫f) (j^*I)` from `IsSFlat f I` is FALSE in general (j a fibre
  % preimmersion is not flat), so this anchor states exactly the input AK use.
  % Standard commutative algebra / EGA; absent from Mathlib at the sheaf level
  % (parallel to thm:flat_pullback_exact).
  % SOURCE: [Altman--Kleiman], §1, (1.3) proof p. 56 (read from references/MR0555258-compactifying-picard.pdf)
  % SOURCE QUOTE: "Consider an exact sequence, 0 -> K -> J -> I -> 0, in which J is
  % as specified in (1.2). Since I is S-flat, the sequence remains exact when
  % restricted to X(s)."
  \textit{Source: Altman--Kleiman, §1, (1.3) proof p. 56 (``Since $I$ is $S$-flat,
  the sequence remains exact when restricted to $X(s)$'').}
  Let $f : X \to S$, $s \in S$, and let
  $0 \to K \to J \to I \to 0$ be a short exact sequence of $\Ox$-modules whose
  right-hand term $I = \mathtt{sc.X_3}$ is $S$-flat (\cref{def:isSFlat}). Then its
  restriction to the closed fibre, i.e.\ its image under
  $j^* = (\mathtt{pullback}\,(f.\mathtt{fiber}\iota\,s))$, is again short exact:
  $0 \to j^*K \to j^*J \to j^*I \to 0$.
  \begin{proof}
    Restriction to the fibre is right exact; flatness of $I$ over the base forces
    the left-hand map to remain a monomorphism after $\tensor_{\Os} k(s)$
    (Tor-vanishing). Standard; used here as an external sheaf-level input.
  \end{proof}
\end{theorem}

\begin{theorem}[Fibre pushforward of a base-change is the residue tensor (standard)]
  \label{thm:ega_fibre_tensor_residue}
  \lean{MR0555258CompactifyingPicard.External.fibre_tensor_residue_iso}
  \uses{def:tensorBaseChangeFunctor, def:tildeMod}
  % NOTE (iter-029): the fibre identification (brick 5) that turns the homological
  % assembly lem:ext1_X_pushforward_fibre_vanishes (Ext^1_X(I, j_*j^*F) = 0) into
  % the residue value T(k(s)) = 0 of the functor T. For A Noetherian local with
  % closed point s and residue field k = k(s), and f : X → Spec A, the projection /
  % base-change formula gives j_*j^*F ≅ F ⊗_S k(s)~ (= the second argument of
  % (functorT A f I F).obj k). General EGA base change along the closed immersion of
  % the closed fibre; NOT §1 content. AK use it implicitly ("the latter Ext is just
  % T(k(s))"). Absent from Mathlib at this generality.
  % SOURCE: [Altman--Kleiman], §1, (1.3) proof p. 56 (read from references/MR0555258-compactifying-picard.pdf)
  % SOURCE QUOTE: "Hence Ext^1_X(I, j_*F(s)) is equal to zero. However, the latter
  % Ext is just T(k(s))."
  \textit{Source: the identification implicit in Altman--Kleiman, §1, (1.3) proof
  p. 56 (``the latter Ext is just $T(k(s))$''); standard EGA base change.}
  Let $A$ be a Noetherian local ring with closed point $s$ and residue field
  $k = k(s)$, let $f : X \to \Spec A$ be a morphism, and $F$ an $\Ox$-module.
  Writing $j = f.\mathtt{fiber}\iota\,s$ for the inclusion of the closed fibre,
  there is an isomorphism of $\Ox$-modules
  \[
    j_* j^* F \;\cong\; F \tensor_S \widetilde{k(s)}
      \;=\; (\mathtt{tensorBaseChangeFunctor}\,f\,F).\mathtt{obj}\,(\widetilde{k}),
  \]
  where $\widetilde{k}$ is the tilde of the residue field viewed as an
  $A$-module (\cref{def:tildeMod}).
  \begin{proof}
    The closed fibre $X(s)$ is the base change $X \times_{\Spec A} \Spec k(s)$;
    the projection formula identifies $j_* j^* F$ with $F \tensor_{\Os}
    s_*\widetilde{k(s)}$, i.e.\ $F$ tensored with the residue field pulled to $X$.
    Standard EGA base change; used here as an external input.
  \end{proof}
\end{theorem}

\begin{theorem}[Fibre pushforward of a quasi-coherent pullback is quasi-coherent (standard)]
  \label{thm:ega_fibre_pushforward_qcoh}
  \lean{MR0555258CompactifyingPicard.External.fibre_pushforward_qcoh}
  \uses{def:fiberExtVanishes}
  % NOTE (iter-029): the qcoh-transport discharging the hWqc hypothesis of
  % lem:ext1_X_pushforward_fibre_vanishes. Mathlib v4.30.0 has NO instance/lemma
  % giving IsQuasicoherent through Scheme.Modules.pullback / pushforward; the
  % general-S setting of ext1_X_pushforward_fibre_vanishes is not affine-local, so
  % the local tilde iso (thm:ega_fibre_tensor_residue) does not apply. Anchored as
  % the standard fact that for the (quasi-compact, quasi-separated) closed-fibre
  % inclusion j, the pushforward of a quasi-coherent sheaf is quasi-coherent. Small
  % sanctioned anchor, parallel to the other EGA fibre inputs. Standard EGA / [EGA
  % I, 9.2]; absent from Mathlib at this generality.
  \textit{Standard: for $j = f.\mathtt{fiber}\iota\,s$ the closed-fibre inclusion
  (quasi-compact, quasi-separated) and $F$ quasi-coherent, $j_*j^*F$ is
  quasi-coherent.}
  Let $f : X \to S$, $s \in S$, and $F$ a quasi-coherent $\Ox$-module. Writing
  $j = f.\mathtt{fiber}\iota\,s$, the $\Ox$-module $j_*j^*F$ is quasi-coherent.
  \begin{proof}
    The pullback $j^*F$ of a quasi-coherent sheaf is quasi-coherent, and the
    pushforward of a quasi-coherent sheaf along a quasi-compact, quasi-separated
    morphism is quasi-coherent. Standard; used here as an external input.
  \end{proof}
\end{theorem}

\begin{theorem}[Pullback preserves freeness (standard)]
  \label{thm:pullback_preserves_free}
  \lean{MR0555258CompactifyingPicard.External.pullback_isFreeMod}
  % NOTE (iter-017): residual of the 1.4 assembly. The pullback of a free
  % $\mathcal{O}_Y$-module along a scheme morphism $p : X \to Y$ is free because
  % $p^* \mathcal{O}_Y \cong \mathcal{O}_X$ (structure-sheaf pullback) and pullback,
  % being a left adjoint, preserves the coproduct defining $\free$. Mathlib v4.30.0
  % proves $p^*\mathcal{O}_Y \cong \mathcal{O}_X$ for Scheme.Modules.pullback only
  % under the auxiliary hypothesis that the site functor Opens.map p.base is FINAL
  % (SheafOfModules.pullbackObjFreeIso / pullbackObjUnitToUnit), which is absent for
  % a general scheme morphism (verified iter-016: failed to synthesize
  % (Opens.map p.base).Final). Anchored as the standard input the 1.4 proof uses
  % implicitly in its last step; flagged buildable from the structure-sheaf
  % pullback iso, to build/upstream later.
  % SOURCE: [Altman--Kleiman], §1, (1.4) proof p. 57 (read from references/MR0555258-compactifying-picard.pdf)
  % SOURCE QUOTE: "Since I is flat, the pullback of the sequence on X_0 is the
  % desired sequence on X."
  \textit{Source: as invoked in Altman--Kleiman, §1, (1.4) proof p. 57.}
  Let $p : X \to Y$ be a morphism of schemes and let $M$ be a free
  $\mathcal{O}_Y$-module (\cref{def:isFreeMod}). Then $p^* M$ is a free
  $\mathcal{O}_X$-module.
  \begin{proof}
    Immediate from $p^*\mathcal{O}_Y \cong \mathcal{O}_X$ and the fact that
    pullback preserves direct sums (it is a left adjoint). Used here as an
    external input at the sheaf level.
  \end{proof}
\end{theorem}

\begin{theorem}[Pullback preserves finite presentation (standard)]
  \label{thm:pullback_preserves_fp}
  \lean{MR0555258CompactifyingPicard.External.pullback_isLFP}
  % NOTE (iter-017): residual of the 1.4 assembly. Finite presentation of a sheaf
  % of modules is preserved by pullback (base change): a local presentation
  % $\mathcal{O}^{(\Lambda')} \to \mathcal{O}^{(\Lambda)} \to M \to 0$ pulls back to
  % a presentation of $p^*M$ since pullback is right exact and carries finite free
  % to finite free. Mathlib v4.30.0 has no instance/lemma supplying
  % SheafOfModules.IsFinitePresentation ((Scheme.Modules.pullback p).obj M) from
  % M.IsFinitePresentation (verified iter-016: infer_instance fails; the
  % local-presentation reconstruction through pullback is absent). Anchored as the
  % standard input the 1.4 proof uses implicitly; flagged buildable, to
  % build/upstream later.
  % SOURCE: [Altman--Kleiman], §1, (1.4) proof p. 57 (read from references/MR0555258-compactifying-picard.pdf)
  % SOURCE QUOTE: "Since I is flat, the pullback of the sequence on X_0 is the
  % desired sequence on X."
  \textit{Source: as invoked in Altman--Kleiman, §1, (1.4) proof p. 57.}
  Let $p : X \to Y$ be a morphism of schemes and let $M$ be a finitely presented
  $\mathcal{O}_Y$-module (\cref{def:isLFP}). Then $p^* M$ is a finitely presented
  $\mathcal{O}_X$-module.
  \begin{proof}
    A local finite free presentation of $M$ pulls back, via the right-exact
    functor $p^*$ which sends finite free to finite free, to a local finite free
    presentation of $p^* M$. Standard (finite presentation is stable under base
    change); used here as an external input at the sheaf level.
  \end{proof}
\end{theorem}

\begin{theorem}[Flat base change / limit isomorphism for $\Hom$, $q=0$ (EGA)]
  \label{thm:ega_basechange_q0}
  \lean{MR0555258CompactifyingPicard.External.basechange_q0}
  % NOTE: deferred — Lean decl not yet scaffolded (TODO block in Basic.lean); needs projective limits of schemes/modules + the comparison map, absent from Mathlib.
  % SOURCE: [Altman--Kleiman], §1, (1.5) proof p. 57 (read from references/MR0555258-compactifying-picard.pdf)
  % SOURCE QUOTE: "For q = 0, the assertion results from [EGA IV_3, 8.5.2 (i)]."
  \textit{Source: as invoked in Altman--Kleiman, §1, ref [EGA IV$_3$, 8.5.2(i)].}
  For a projective limit $S = \varprojlim S_\lambda$ of $S_0$-schemes with the
  flatness/finite-presentation hypotheses of \cref{lem:ext_limit}, the canonical
  map $\varprojlim \Hom_{X_\lambda}(I_\lambda, F_\lambda) \to \Hom_X(I,F)$ is an
  isomorphism (the case $q=0$ of the $\sExt$ limit isomorphism).
  \begin{proof}
    Established in [EGA IV$_3$, 8.5.2(i)]; used here as an external input.
  \end{proof}
\end{theorem}

\begin{theorem}\leanok
[Adjunction isomorphism for $\Hom$ (EGA)]
  \label{thm:ega_adjunction}
  \lean{MR0555258CompactifyingPicard.External.adjunction}
  \mathlibok
  % NOTE: REALIZED @iter-023 as a PROVED `noncomputable def` (NOT an axiom): the
  % degree-0 / hom-set form actually consumed by (1.3), `(I ⟶ j_* G) ≃ (j^* I ⟶ G)`
  % for a general `j : Y ⟶ X`, = `((Scheme.Modules.pullbackPushforwardAdjunction j).homEquiv I G).symm`.
  % Mathlib-backed (no new obligation), hence \mathlibok. The fibre-product SHEAF-Hom
  % form stated in the prose below (`Hom_X(I,(1×g)_* N) ≅ (1×g)_* Hom_{X_T}(I_T,N)`)
  % is a strictly stronger statement still absent from Mathlib; the Lean decl is its
  % degree-0 instance. PLANNER: reconcile the prose to the realized hom-set form.
  % SOURCE: [Altman--Kleiman], §1, (1.7) proof p. 58 (read from references/MR0555258-compactifying-picard.pdf)
  % SOURCE QUOTE: "For q = 0, the isomorphism (1.7.1) comes from the usual
  % adjunction isomorphism [EGA O_I, 4.4.3.1]."
  \textit{Source: as invoked in Altman--Kleiman, §1, ref [EGA O$_I$, 4.4.3.1].}
  For a base change $g : T \to S$ and a quasi-coherent $\Ot$-module $N$ there is a
  canonical adjunction isomorphism
  $\Hom_X(I, (1\times g)_* N) \cong (1\times g)_* \Hom_{X_T}(I_T, N)$.
  \begin{proof}
    Established in [EGA O$_I$, 4.4.3.1]; used here as an external input.
  \end{proof}
\end{theorem}

\begin{theorem}[Vanishing of a half-exact functor over a Noetherian local ring (OB)]
  \label{thm:ob_halfexact_free}
  \lean{MR0555258CompactifyingPicard.External.halfexact_free}
  % SOURCE: [Altman--Kleiman], §1, (1.3) proof p. 56 (read from references/MR0555258-compactifying-picard.pdf)
  % SOURCE QUOTE: "Since T is half-exact and since T(k(s)) is equal to zero,
  % T(M) is equal to zero for every finitely generated A-module M [OB, 2.1] or
  % [EGA III_2, 7.5.3])."
  \textit{Source: as invoked in Altman--Kleiman, §1, ref [OB, 2.1] / [EGA III$_2$, 7.5.3].}
  Let $A$ be a Noetherian local ring with residue field $k$, and let $T$ be a
  half-exact, additive functor from finitely generated $A$-modules to finitely
  generated $A$-modules. If $T(k) = 0$ then $T(M) = 0$ for every finitely
  generated $A$-module $M$; consequently the adjacent ($\Hom$-side) functor is
  exact, and the corresponding finitely presented $A$-module is free.

  \emph{(Formalization note.)} The Lean signature takes $T :
  \mathrm{Mod}_A^{(u)} \to \mathrm{Mod}_A^{(u+1)}$, i.e.\ the value category sits
  one universe up from the source. This is a faithful encoding, not a
  strengthening: the only instance of $T$ we feed (the functor of
  \cref{def:functorT}, $T(M) = \Ext^1_X(I, F\tensor_S\widetilde M)$) produces its
  values through the standard derived category of $X\text{-Mod}$, whose hom-groups
  live one universe above the objects, so $\Ext$ lands in $\mathrm{Ab}^{(u+1)}$.
  Mathematically $T$ is an endofunctor of finitely generated $A$-modules; the
  universe split records how the $\Ext$ groups are computed in Lean and the
  statement (OB 2.1 / EGA III$_2$ 7.5.3) is itself universe-agnostic.
  \begin{proof}
    Established in [OB, 2.1] / [EGA III$_2$, 7.5.3]; used here as an external input.
  \end{proof}
\end{theorem}

\begin{theorem}[Nakayama-type isomorphism for a limit-preserving functor (OB)]
  \label{thm:ob_nakayama_iso}
  \lean{MR0555258CompactifyingPicard.External.nakayama_iso}
  % NOTE: deferred — Lean decl not yet scaffolded (TODO block in Basic.lean); phrasable in pure module theory but needs the explicit base-change comparison natural transformation to state faithfully.
  % SOURCE: [Altman--Kleiman], §1, (1.9) proof p. 60 (read from references/MR0555258-compactifying-picard.pdf)
  % SOURCE QUOTE: "Therefore, by [OB, 4.1], the map, R(O_s) \otimes_{O_s} N ->
  % R(N), (1.9.1)"
  \textit{Source: as invoked in Altman--Kleiman, §1, ref [OB, 4.1].}
  Let $R$ be an additive functor from $\mathcal{O}_s$-modules to
  $\mathcal{O}_s$-modules that commutes with direct limits and sends finitely
  generated modules to finitely generated modules, where $\mathcal{O}_s$ is a
  local ring. If the natural map $R(\mathcal{O}_s) \tensor_{\mathcal{O}_s} k(s)
  \to R(k(s))$ is surjective, then the natural map
  $R(\mathcal{O}_s) \tensor_{\mathcal{O}_s} N \to R(N)$ is an isomorphism for
  every $\mathcal{O}_s$-module $N$.
  \begin{proof}
    Established in [OB, 4.1]; used here as an external input.
  \end{proof}
\end{theorem}

\begin{theorem}[Openness of the local Ext-vanishing locus (EGA)]
  \label{thm:ega_ext_vanishing_open}
  \lean{MR0555258CompactifyingPicard.External.ext_vanishing_open}
  % SOURCE: [Altman--Kleiman], §1, (1.10) proof p. 61, ref [EGA IV_3, 12.3.4] (read from references/MR0555258-compactifying-picard.pdf)
  % SOURCE QUOTE: "Then the proof of [EGA IV_3, 12.3.4] shows that U is open and
  % retrocompact, although this is not fully stated."
  \textit{Source: as invoked in Altman--Kleiman §1 (Thm 1.10), ref [EGA IV$_3$ 12.3.4].}
  Under the hypotheses of \cref{thm:ext_finite_flat} (so $f : X \to S$ finitely
  presented, $I, F$ locally finitely presented with $F$ $S$-flat, and the flatness
  hypotheses on $I, f$ in the relevant degrees), the locus
  \[
    \{\, s \in S : \sExt^q_{X(s)}\!\big(I(s), F(s)\big) = 0 \,\}
  \]
  is open and retrocompact in $S$.
  \begin{proof}
    Established in [EGA IV$_3$, 12.3.4]; used here as an external input.
  \end{proof}
\end{theorem}

\begin{theorem}[Coherence / local finite presentation of local Ext (EGA)]
  \label{thm:ega_ext_coherent_fp}
  \lean{MR0555258CompactifyingPicard.External.ext_coherent_fp}
  % SOURCE: [Altman--Kleiman], §1, (1.10) proof p. 61, ref [EGA III_2, 12.3.3] (read from references/MR0555258-compactifying-picard.pdf)
  % SOURCE QUOTE: "Finally, the assertion of local finite presentation is local on
  % S and is compatible with base-change. So we may assume S is affine and by
  % [EGA IV_3, Sects. 8, 11] Noetherian. Then the assertion holds by
  % [EGA III_2, 12.3.3]."
  \textit{Source: as invoked in Altman--Kleiman §1 (Thm 1.10), ref [EGA III$_2$ 12.3.3].}
  Over a Noetherian (affine) base $S$, with $f : X \to S$ finitely presented and
  proper and $I, F$ locally finitely presented, the relative local Ext module
  $\sExt^c_X(I, F)$ is coherent, i.e. locally finitely presented.
  \begin{proof}
    Established in [EGA III$_2$, 12.3.3]; used here as an external input.
  \end{proof}
\end{theorem}

%==============================================================================
% Project-local scaffolding (notation of (1.1)): the admissibility hypotheses,
% the base-change tensor F ⊗_S M, and the corepresented functor. These carry no
% mathematical content beyond the paper's (1.1) setup; they exist so the §1
% statements can be transcribed in the project's notation.
%==============================================================================

% The five blocks below name the standard sheaf-module notions Altman--Kleiman
% use in (1.1) (locally finitely presented, $S$-flat, free, invertible, and the
% tensor product of two $\Ox$-modules). They carry no content beyond their
% standard meaning; we record them with their intended Mathlib realization so the
% §1 statements have honest, constructible underpinnings (replacing the earlier
% opaque placeholders).

\begin{definition}

\leanok
\mathlibok
[Locally finitely presented $\Ox$-module]
  \label{def:isLFP}
  \lean{MR0555258CompactifyingPicard.IsLFP}
  An $\Ox$-module $M$ is \emph{locally finitely presented} when it admits, locally
  on $X$, a finite presentation $\mathcal{O}_X^{\,p} \to \mathcal{O}_X^{\,q} \to M
  \to 0$. Realized via Mathlib's \texttt{SheafOfModules.IsFinitePresentation} for
  the sheaf of modules underlying $M : X.\mathrm{Modules}$.
\end{definition}

\begin{definition}

\leanok
[$S$-flat $\Ox$-module]
  \label{def:isSFlat}
  \lean{MR0555258CompactifyingPicard.IsSFlat}
  For $f : X \to S$, an $\Ox$-module $M$ is \emph{flat over $S$} (along $f$) when,
  on affine opens $V = \Spec A \subseteq X$ over $U = \Spec R \subseteq S$, the
  $A$-module $M(V)$ is flat over $R$. Realized via stalkwise/affine-local flatness
  of the module over the base ring maps (cf.\ Mathlib's morphism-flatness
  \texttt{AlgebraicGeometry.Flat}, transported to module-level flatness over $f$).
\end{definition}

\begin{definition}

\leanok
[Tensor product of $\Ox$-modules $M \tensor N$]
  \label{def:tensorMod}
  \lean{MR0555258CompactifyingPicard.tensorMod}
  For $\Ox$-modules $M, N$ their tensor product $M \tensor_{\Ox} N$ is the
  sheafification of the presheaf tensor product. Realized via Mathlib's
  presheaf-level monoidal tensor \texttt{PresheafOfModules.Monoidal.tensorObj}
  followed by sheafification; note Mathlib has no \texttt{MonoidalCategory}
  instance on $X.\mathrm{Modules}$ directly, so the tensor must be assembled
  from the presheaf monoidal structure.
\end{definition}

%==============================================================================
% Symmetric monoidal scaffolding for the tensor of \cref{def:tensorMod}.
% PROJECT-BESPOKE (no external source): Mathlib supplies the monoidal structure
% on PRESHEAVES of modules (PresheafOfModules.monoidalCategory, with associators
% and unitors) but NOT on SheafOfModules over a sheaf of rings, and registers no
% braided/symmetric instance even at the presheaf level. These isomorphisms are
% built by transporting the pointwise module-level structure (TensorProduct.comm
% / .assoc / .lid) through the presheaf tensor and sheafification. They are the
% algebraic engine of the H-twist isomorphisms (1.1.2)/(1.1.3).
%==============================================================================

\begin{lemma}\leanok
[Braiding (symmetry) of the $\Ox$-tensor]
  \label{lem:tensorMod_braiding}
  \lean{MR0555258CompactifyingPicard.tensorBraiding}
  \uses{def:tensorMod, def:presheafTensorBraiding}
  For $\Ox$-modules $A, B$ there is a natural isomorphism
  $A \tensor_{\Ox} B \cong B \tensor_{\Ox} A$, the symmetry of the tensor.
  \begin{proof}
    The rings $\mathcal{O}_X(U)$ are commutative, so at the presheaf level the
    pointwise swap $m \tensor n \mapsto n \tensor m$ (\texttt{TensorProduct.comm})
    assembles into an isomorphism of presheaves of modules
    (\cref{def:presheafTensorBraiding}); applying sheafification (a functor) gives
    the asserted isomorphism of the sheafified tensors of \cref{def:tensorMod},
    natural in $A$ and $B$.
  \end{proof}
\end{lemma}

\begin{definition}

\leanok
[Presheaf-level braiding]
  \label{def:presheafTensorBraiding}
  \lean{MR0555258CompactifyingPicard.presheafTensorBraiding}
  For presheaves of modules $M, N$ over a presheaf of commutative rings, the
  presheaf-level swap $M \tensor N \cong N \tensor M$ of the monoidal tensor
  \texttt{PresheafOfModules.Monoidal.tensorObj}. Mathlib registers no
  braided/symmetric instance on \texttt{PresheafOfModules}, so it is built by hand:
  each section ring $R(U)$ is commutative, so the pointwise tensor carries the
  \texttt{ModuleCat} symmetric braiding $\beta$; these assemble (via
  \texttt{PresheafOfModules.isoMk}) into an isomorphism of presheaves of modules,
  naturality reducing on pure tensors to the defining swap.
\end{definition}

\begin{lemma}\leanok
[Associativity of the $\Ox$-tensor (conditional)]
  \label{lem:tensorMod_assoc}
  \lean{MR0555258CompactifyingPicard.tensorAssocOfLocallyBijective}
  \uses{def:tensorMod, lem:sheafifyTensorComparison, lem:sheafifyTensorComparisonLeft}
  For $\Ox$-modules $A, B, C$, \emph{given} that the four sheafification-unit
  whiskerings $\eta_{A\tensor B}\tensor\id_C$, $\id_A\tensor\eta_{B\tensor C}$
  (both handednesses) are locally bijective, there is a natural isomorphism
  $(A \tensor_{\Ox} B) \tensor_{\Ox} C \cong A \tensor_{\Ox} (B \tensor_{\Ox} C)$.
  (Project-bespoke; the local-bijectivity hypotheses isolate the genuine Mathlib
  gap — see \cref{lem:sheafifyTensorComparison}. The unconditional associator
  would follow from a $\texttt{MonoidalClosed (PresheafOfModules }R)$ instance.)
  \begin{proof}
    Mathlib's monoidal structure on presheaves of modules supplies the pointwise
    associator (\texttt{TensorProduct.assoc}). Because \cref{def:tensorMod}
    sheafifies the presheaf tensor, $(A \tensor B)\tensor C$ unfolds to the
    sheafification of a presheaf tensor one of whose factors,
    $(\widetilde{A\tensor B})$, is itself sheafified. The comparison
    \cref{lem:sheafifyTensorComparison} rewrites this back to the sheafification of
    the un-presheafified tensor; the presheaf associator then applies, and a second
    use of \cref{lem:sheafifyTensorComparison} on the other side reassembles the
    target. Concretely the chain is
    $\widetilde{(\widetilde{A\tensor B})\tensor C}
       \cong \widetilde{(A\tensor B)\tensor C}
       \cong \widetilde{A\tensor(B\tensor C)}
       \cong \widetilde{A\tensor(\widetilde{B\tensor C})}$.
  \end{proof}
\end{lemma}

%==============================================================================
% (1.1.2) naturality-layer scaffolding --- project-bespoke Lean-realization helpers.
% These carry NO Altman--Kleiman content; they package module sheafification and the
% presheaf tensor into the hand-built bifunctors and cheap-reduction bridge lemmas
% that the natural form of the comparison anchors and the naturality of tensorAssoc
% are stated and proved against (iters 030--032).
%==============================================================================

\begin{definition}

\leanok
[Coyoneda transported through a fully faithful endofunctor]
  \label{def:coyonedaCompFF}
  \lean{MR0555258CompactifyingPicard.coyonedaCompFF}
  For a fully faithful endofunctor $G$ of a category $\mathcal{C}$ and an object $I$,
  the $\mathrm{Type}$-valued natural isomorphism
  $G \circ \coyoneda(G\,I) \cong \coyoneda(I)$, i.e. $(G\,I \to G\,Z) \cong (I \to Z)$
  natural in $Z$. (Mathlib's \texttt{Functor.FullyFaithful.homNatIso'} lands in
  \texttt{ULift}-ed homs; this $\mathrm{Type}\,u$-valued form is what the $(1.1.2)$
  Yoneda step \cref{lem:H_tensor_invertible} needs.) Built from the full-faithfulness
  hom-equivalence, naturality by \texttt{aesop\_cat}.
\end{definition}

\begin{definition}

\leanok
[Module sheafification, packaged as a functor]
  \label{def:moduleSheafification}
  \lean{MR0555258CompactifyingPicard.moduleSheafification}
  The module-sheafification functor $\widetilde{(-)} : \PresheafOfModules \to
  \SheafOfModules$ over $\mathcal{O}_X$, repackaged in applied $\texttt{obj}/\texttt{map}$
  form (rather than the bare \texttt{PresheafOfModules.sheafification} functor), so the
  $\texttt{IsLocallyInjective}/\texttt{IsLocallySurjective}$ instance synthesis driving
  the comparison anchors resolves without leaving metavariables.
\end{definition}

\begin{definition}

\leanok
[Sheafify-then-forget]
  \label{def:sheafifyForget}
  \lean{MR0555258CompactifyingPicard.sheafifyForget}
  \uses{def:moduleSheafification}
  The endofunctor $P \mapsto (\widetilde{P})_{\mathrm{val}}$ on presheaves of
  $\mathcal{O}_X$-modules, sheafify then forget back to the underlying presheaf. Used to
  phrase the bifunctorial (natural) form of the comparison anchors.
\end{definition}

\begin{definition}

\leanok
[The bifunctor $(P,Q) \mapsto \widetilde{P \tensor Q}$]
  \label{def:tensorThenSheafify}
  \lean{MR0555258CompactifyingPicard.tensorThenSheafify}
  \uses{def:tensorMod, def:moduleSheafification}
  The bifunctor on presheaves of $\mathcal{O}_X$-modules sending $(P,Q)$ to
  $\widetilde{P \tensor Q}$ (presheaf monoidal tensor followed by sheafification) and a
  morphism $(f_1,f_2)$ to $\widetilde{f_1 \tensor f_2}$. The common right-hand side of
  both comparison anchors. Hand-built with an \emph{explicit} action on morphisms (rather
  than $\texttt{MonoidalCategory.tensor} \circ \texttt{moduleSheafification}$) so that
  $\texttt{.map}\,(a,b)$ reduces by a single $\texttt{rfl}$ to the hand-built
  $\texttt{tensorLeft}/\texttt{tensorRight}$ morphism maps; the monoidal-instance
  $\texttt{tensorHom}$ would blow up $\texttt{whnf}$. Objects are unchanged, so the
  anchors and $\texttt{tensorAssoc}$ (which use only $\texttt{.app}$) are unaffected.
\end{definition}

\begin{definition}

\leanok
[The bifunctor $(P,Q) \mapsto \widetilde{(\widetilde{P})_{\mathrm{val}} \tensor Q}$]
  \label{def:sheafifyFstTensorThenSheafify}
  \lean{MR0555258CompactifyingPicard.sheafifyFstTensorThenSheafify}
  \uses{def:tensorThenSheafify, def:moduleSheafification}
  Sheafify the first factor, then $\texttt{tensorThenSheafify}$: the left-hand side of
  $\texttt{External.sheafifyTensorComparison}$. Hand-built (explicit
  $\texttt{obj}/\texttt{map}$) for cheap reduction, as in \cref{def:tensorThenSheafify}.
\end{definition}

\begin{definition}

\leanok
[The bifunctor $(P,Q) \mapsto \widetilde{P \tensor (\widetilde{Q})_{\mathrm{val}}}$]
  \label{def:sheafifySndTensorThenSheafify}
  \lean{MR0555258CompactifyingPicard.sheafifySndTensorThenSheafify}
  \uses{def:tensorThenSheafify, def:moduleSheafification}
  The left-handed mirror of \cref{def:sheafifyFstTensorThenSheafify}: sheafify the second
  factor, then $\texttt{tensorThenSheafify}$. The left-hand side of
  $\texttt{External.sheafifyTensorComparisonLeft}$.
\end{definition}

\begin{lemma}

\leanok
[Naturality of the presheaf braiding in its first factor]
  \label{lem:presheafTensorBraiding_naturality}
  \lean{MR0555258CompactifyingPicard.presheafTensorBraiding_naturality}
  \uses{def:presheafTensorBraiding}
  For $f : M \to M'$ the square
  $(f \tensor \id_N) \circ \beta_{M',N} = \beta_{M,N} \circ (\id_N \tensor f)$ commutes,
  the pointwise \texttt{ModuleCat} braiding naturality
  (\texttt{BraidedCategory.braiding\_naturality}) assembled over the site. The
  first-variable-naturality input to \cref{def:tensorBraidingNatIso}.
\end{lemma}

\begin{definition}

\leanok
[The braiding as a natural isomorphism of endofunctors]
  \label{def:tensorBraidingNatIso}
  \lean{MR0555258CompactifyingPicard.tensorBraidingNatIso}
  \uses{lem:tensorMod_braiding, lem:presheafTensorBraiding_naturality, def:tensorRight}
  The $\mathcal{O}_X$-tensor braiding \cref{lem:tensorMod_braiding} packaged as a natural
  isomorphism of endofunctors $\texttt{tensorRight}\,B \cong \texttt{tensorLeft}\,B$ (i.e.
  $A \tensor B \cong B \tensor A$ natural in $A$); naturality in $A$ from
  \cref{lem:presheafTensorBraiding_naturality} and functoriality of sheafification.
\end{definition}

\begin{lemma}

\leanok
[Cheap unfold of $\texttt{tensorLeft}$ on morphisms]
  \label{lem:tensorLeft_map_eq}
  \lean{MR0555258CompactifyingPicard.tensorLeft_map_eq}
  \uses{def:tensorMod}
  $(\texttt{tensorLeft}\,F).\texttt{map}\,\varphi = \widetilde{\id_F \tensor \varphi}$, a
  single $\texttt{rfl}$. (Private bridge; lets the naturality of $\texttt{tensorAssoc}$ be
  rewritten against the anchor bifunctors without forming a concrete object equation.)
\end{lemma}

\begin{lemma}

\leanok
[Cheap unfold of $\texttt{tensorRight}$ on morphisms]
  \label{lem:tensorRight_map_eq}
  \lean{MR0555258CompactifyingPicard.tensorRight_map_eq}
  \uses{def:tensorRight}
  $(\texttt{tensorRight}\,L).\texttt{map}\,\varphi = \widetilde{\varphi \tensor \id_L}$, a
  single $\texttt{rfl}$. Private bridge, mirror of \cref{lem:tensorLeft_map_eq}.
\end{lemma}

\begin{lemma}

\leanok
[Cheap unfold of $\texttt{sheafifyFstTensorThenSheafify}$ on a generic morphism]
  \label{lem:sheafifyFst_map_eq}
  \lean{MR0555258CompactifyingPicard.sheafifyFst_map_eq}
  \uses{def:sheafifyFstTensorThenSheafify}
  For a \emph{generic} morphism $f$ (not a pair $(a,b)$),
  $(\texttt{sheafifyFstTensorThenSheafify}\,X).\texttt{map}\,f
  = \widetilde{(\widetilde{f_1})_{\mathrm{val}} \tensor f_2}$, a single $\texttt{rfl}$. The
  generic form keeps both sides syntactically $\widetilde{(-)}$, avoiding the $\texttt{isDefEq}$
  blowup the pair form $(a,b)$ triggers. Private bridge.
\end{lemma}

\begin{lemma}

\leanok
[Cheap unfold of $\texttt{sheafifySndTensorThenSheafify}$ on a generic morphism]
  \label{lem:sheafifySnd_map_eq}
  \lean{MR0555258CompactifyingPicard.sheafifySnd_map_eq}
  \uses{def:sheafifySndTensorThenSheafify}
  Mirror of \cref{lem:sheafifyFst_map_eq}:
  $(\texttt{sheafifySndTensorThenSheafify}\,X).\texttt{map}\,f
  = \widetilde{f_1 \tensor (\widetilde{f_2})_{\mathrm{val}}}$, a single $\texttt{rfl}$. Private bridge.
\end{lemma}

\begin{lemma}

\leanok
[Cheap unfold of $\texttt{tensorThenSheafify}$ on a generic morphism]
  \label{lem:tensorThenSheafify_map_eq}
  \lean{MR0555258CompactifyingPicard.tensorThenSheafify_map_eq}
  \uses{def:tensorThenSheafify}
  For a generic morphism $f$ of pairs,
  $(\texttt{tensorThenSheafify}\,X).\texttt{map}\,f = \widetilde{f_1 \tensor f_2}$, a single
  $\texttt{rfl}$. Private bridge (companion of \cref{lem:sheafifyFst_map_eq}), used to keep
  the associator-naturality transport on the abstract functor without forming a concrete
  object equation.
\end{lemma}

\begin{lemma}

\leanok
[Presheaf-level $\id \tensor \id = \id$]
  \label{lem:presheaf_id_tensorHom_id}
  \lean{MR0555258CompactifyingPicard.presheaf_id_tensorHom_id}
  \uses{def:tensorMod}
  At the presheaf-of-modules level, $\id_M \tensor \id_N = \id_{M \tensor N}$. Proved
  pointwise from Mathlib's \texttt{MonoidalCategory.id\_tensorHom\_id} (which here fires only
  as a term, not via \texttt{rw}/\texttt{simp}). Private helper feeding the sheafify-of-identity
  legs of \cref{lem:tensorAssocNatIso3} and \cref{lem:tensorAssocNatIso1}.
\end{lemma}

\begin{lemma}

\leanok
[Sheafification sends an identity to an identity (on underlying values)]
  \label{lem:sheafify_map_id_val}
  \lean{MR0555258CompactifyingPicard.sheafify_map_id_val}
  \uses{def:moduleSheafification}
  $(\widetilde{\id_P})_{\mathrm{val}} = \id_{(\widetilde{P})_{\mathrm{val}}}$, i.e.\ the
  module-sheafification functor carries the identity to the identity (read off on the
  underlying-presheaf value). Via \texttt{CategoryTheory.Functor.map\_id} and
  \texttt{SheafOfModules.id\_val}. Private helper for the sheafify-of-identity legs.
\end{lemma}

\begin{definition}

\leanok
[Index functors for middle-factor associator whiskering]
  \label{def:assocIdxL}
  \lean{MR0555258CompactifyingPicard.assocIdxL}
  \uses{def:tensorMod}
  The private index functor $Y \mapsto (F.\mathrm{val} \tensor Y.\mathrm{val},\ L.\mathrm{val})$
  on $\Ox$-modules, packaging the two product slots seen by one leg of the middle-factor
  associator naturality \cref{lem:tensorAssocNatIso2}. Together with \cref{def:assocIdxR} it
  lets each leg's naturality be a single whiskered comparison square.
\end{definition}

\begin{definition}

\leanok
[Index functor, right form, middle-factor]
  \label{def:assocIdxR}
  \lean{MR0555258CompactifyingPicard.assocIdxR}
  \uses{def:tensorMod}
  The private index functor $Y \mapsto (F.\mathrm{val},\ Y.\mathrm{val} \tensor L.\mathrm{val})$,
  the right-form companion of \cref{def:assocIdxL} for \cref{lem:tensorAssocNatIso2}.
\end{definition}

\begin{definition}

\leanok
[Index functors for last-factor associator whiskering]
  \label{def:assocIdxL3}
  \lean{MR0555258CompactifyingPicard.assocIdxL₃}
  \uses{def:tensorMod}
  The private index functor $Y \mapsto (F.\mathrm{val} \tensor L.\mathrm{val},\ Y.\mathrm{val})$
  for the last-factor associator naturality \cref{lem:tensorAssocNatIso3}.
\end{definition}

\begin{definition}

\leanok
[Index functor, right form, last-factor]
  \label{def:assocIdxR3}
  \lean{MR0555258CompactifyingPicard.assocIdxR₃}
  \uses{def:tensorMod}
  The private index functor $Y \mapsto (F.\mathrm{val},\ L.\mathrm{val} \tensor Y.\mathrm{val})$,
  the right-form companion of \cref{def:assocIdxL3} for \cref{lem:tensorAssocNatIso3}.
\end{definition}

\begin{lemma}
[Sheafification commutes with the presheaf tensor (unconditional) --- Mathlib gap anchor]
  \label{lem:sheafifyTensorComparison_uncond}
  \lean{MR0555258CompactifyingPicard.External.sheafifyTensorComparison}
  \uses{def:tensorMod, def:sheafifyTensorUnitWhiskerRight, def:sheafifyTensorUnitWhiskerLeft, def:sheafifyFstTensorThenSheafify, def:tensorThenSheafify}
  % NOTE: This is NOT from Altman--Kleiman; it is the standard general fact that the
  % module-sheafification functor ~(-) on presheaves of modules is monoidal, i.e.
  % the comparison maps ~((~P) (x) Q) -> ~(P (x) Q) and ~(P (x) (~Q)) -> ~(P (x) Q)
  % are isomorphisms. It is TRUE in full generality (sheafification is a monoidal
  % localization), but is ABSENT from Mathlib for PresheafOfModules: the naive
  % "tensoring preserves local isomorphisms" route fails on the injectivity half
  % (iter-008 finding, recorded at lem:sheafifyTensorComparison), and the
  % localization-monoidal route (CategoryTheory.Localization.Monoidal) needs
  % `W.IsMonoidal` + `L.IsLocalization W` for module sheafification, both absent
  % from Mathlib (verified iters 007--008). Anchored as a Mathlib gap, exactly as
  % `External.affine_fp_tilde` / `External.flat_tensor_exact` anchor true general
  % facts Mathlib lacks at the needed generality. It carries NO section-1 content.
  The module-sheafification functor $\widetilde{(-)}$ on presheaves of modules is
  monoidal: for presheaves of modules $P, Q$ the comparison maps assemble into
  natural isomorphisms
  \[
    \widetilde{\,(\widetilde{P}) \tensor Q\,} \;\cong\; \widetilde{P \tensor Q},
    \qquad
    \widetilde{\,P \tensor (\widetilde{Q})\,} \;\cong\; \widetilde{P \tensor Q}.
  \]
  This is the unconditional form of \cref{lem:sheafifyTensorComparison} /
  \cref{lem:sheafifyTensorComparisonLeft} (their local-bijectivity hypotheses, false
  in general, dropped). The Lean realization is a \emph{natural} isomorphism
  (an isomorphism in the functor category / \texttt{NatIso} between the two bifunctors),
  so that the associator, braiding and rearrangement built on top of it inherit
  naturality; the object-level comparison is recovered as its component
  ($\texttt{.app}$). This is the only change needed versus the object-level anchor of
  iter-031.
  \begin{proof}
    Standard: sheafification is a monoidal localization. Anchored as a Mathlib gap
    (absent from Mathlib for \texttt{PresheafOfModules}; see the note above). Used as
    an external input.
  \end{proof}
\end{lemma}

\begin{lemma}
[Sheafification--tensor comparison, left handedness (unconditional) --- Mathlib gap anchor]
  \label{lem:sheafifyTensorComparisonLeft_uncond}
  \lean{MR0555258CompactifyingPicard.External.sheafifyTensorComparisonLeft}
  \uses{def:tensorMod, lem:sheafifyTensorComparison_uncond, def:sheafifySndTensorThenSheafify}
  % NOTE: Left-handed companion of lem:sheafifyTensorComparison_uncond, anchored as the
  % same Mathlib gap (module sheafification is monoidal). Carries NO section-1 content.
  The left-handed unconditional comparison of \cref{lem:sheafifyTensorComparison_uncond}:
  for presheaves of modules $P, Q$ there is a natural isomorphism
  $\widetilde{\,P \tensor (\widetilde{Q})\,} \cong \widetilde{P \tensor Q}$.
  \begin{proof}
    Identical to \cref{lem:sheafifyTensorComparison_uncond} with the whiskering on the
    left factor. Anchored as a Mathlib gap; used as an external input.
  \end{proof}
\end{lemma}

\begin{lemma}

\leanok
[Associativity of the $\Ox$-tensor (unconditional)]
  \label{lem:tensorMod_assoc_uncond}
  \lean{MR0555258CompactifyingPicard.tensorAssoc}
  \uses{def:tensorMod, lem:sheafifyTensorComparison_uncond, lem:tensorMod_assoc}
  For $\Ox$-modules $A, B, C$ there is a natural isomorphism
  $(A \tensor_{\Ox} B) \tensor_{\Ox} C \cong A \tensor_{\Ox} (B \tensor_{\Ox} C)$.
  This is the unconditional associator: the conditional construction
  \cref{lem:tensorMod_assoc} run with the unconditional comparison
  \cref{lem:sheafifyTensorComparison_uncond} in place of the (false-hypothesis)
  conditional comparison.
  \begin{proof}
    \uses{lem:sheafifyTensorComparison_uncond, lem:tensorMod_assoc}
    Identical to the chain in \cref{lem:tensorMod_assoc}, with the two uses of the
    sheafification--tensor comparison supplied unconditionally by
    \cref{lem:sheafifyTensorComparison_uncond}; the presheaf associator
    (\texttt{TensorProduct.assoc}, from Mathlib's monoidal structure on presheaves of
    modules) supplies the middle step.
  \end{proof}
\end{lemma}

\begin{lemma}

\leanok
[Rearrangement of three $\Ox$-factors]
  \label{lem:tensorRearrange}
  \lean{MR0555258CompactifyingPicard.tensorRearrange}
  \uses{def:tensorMod, lem:tensorMod_assoc_uncond, lem:tensorMod_braiding}
  For $\Ox$-modules $F, Y, L$ there is an isomorphism
  \[
    (F \tensor_{\Ox} Y) \tensor_{\Ox} L \;\cong\; (F \tensor_{\Ox} L) \tensor_{\Ox} Y,
  \]
  the rearrangement swapping the two right factors. Instantiated at $Y = f^*M$ it is
  the rearrangement $r_M$ of the $(1.1.2)$ twist (\cref{lem:H_tensor_invertible}).
  \begin{proof}
    The chain $(F \tensor Y) \tensor L \cong F \tensor (Y \tensor L)
    \cong F \tensor (L \tensor Y) \cong (F \tensor L) \tensor Y$, using the
    unconditional associator \cref{lem:tensorMod_assoc_uncond} twice and the braiding
    \cref{lem:tensorMod_braiding} (on the inner factors $Y, L$) once.
  \end{proof}
\end{lemma}

\begin{definition}

\leanok
[Tensor-by-$L$ on the right]
  \label{def:tensorRight}
  \lean{MR0555258CompactifyingPicard.tensorRight}
  \uses{def:tensorMod}
  The ``tensor by $L$ on the right'' endofunctor on $\Ox$-modules, $N \mapsto N \tensor_{\Ox} L$,
  the right-handed companion of $\texttt{tensorLeft}$ (\cref{def:tensorMod}). On a morphism
  $\varphi$ it sheafifies the presheaf tensor $\varphi \tensor \id_L$; assembled by hand
  exactly as $\texttt{tensorLeft}$, since $\Ox$-modules carry no
  \texttt{MonoidalCategory} instance.
\end{definition}

\begin{lemma}

\leanok
[Naturality of $\texttt{tensorAssoc}$ in the middle factor]
  \label{lem:tensorAssocNatIso2}
  \lean{MR0555258CompactifyingPicard.tensorAssocNatIso₂}
  \uses{lem:tensorMod_assoc_uncond, lem:sheafifyTensorComparison_uncond, lem:sheafifyFst_map_eq, lem:tensorLeft_map_eq, lem:tensorRight_map_eq, lem:tensorThenSheafify_map_eq, def:assocIdxL, def:assocIdxR}
  For fixed $\Ox$-modules $F, L$ the associator \cref{lem:tensorMod_assoc_uncond} is
  natural in its middle factor $Y$: a natural isomorphism of endofunctors
  $\texttt{tensorLeft}\,F \circ \texttt{tensorRight}\,L \cong \dots$ at the associator step
  $Y \mapsto ((F \tensor Y) \tensor L \cong F \tensor (Y \tensor L))$.
  \begin{proof}
    The associator is built from the comparison anchor
    \cref{lem:sheafifyTensorComparison_uncond} (now a \emph{natural} iso of bifunctors) and
    the sheafified presheaf associator. Its naturality square in $Y$ is the
    $\texttt{NatTrans.naturality}$ of the comparison anchor, transported to the hand-built
    $\texttt{tensorLeft}/\texttt{tensorRight}$ forms. Crucially the transport is done by
    rewriting through the \emph{generic} bridge lemmas \cref{lem:sheafifyFst_map_eq},
    \cref{lem:tensorLeft_map_eq}, \cref{lem:tensorRight_map_eq} (which fire on the functor,
    keeping objects abstract), so no concrete object equation is ever formed and the
    $\texttt{isDefEq}$ blowup is avoided.
  \end{proof}
\end{lemma}

\begin{lemma}

\leanok
[Naturality of $\texttt{tensorAssoc}$ in the last factor]
  \label{lem:tensorAssocNatIso3}
  \lean{MR0555258CompactifyingPicard.tensorAssocNatIso₃}
  \uses{lem:tensorMod_assoc_uncond, lem:sheafifyTensorComparison_uncond, lem:sheafifySnd_map_eq, lem:tensorLeft_map_eq, lem:tensorRight_map_eq, lem:tensorThenSheafify_map_eq, lem:presheaf_id_tensorHom_id, lem:sheafify_map_id_val, def:assocIdxL3, def:assocIdxR3}
  Mirror of \cref{lem:tensorAssocNatIso2} in the last factor: the associator step
  $Y \mapsto (F \tensor (L \tensor Y) \cong (F \tensor L) \tensor Y)$ is natural in $Y$,
  proved through the generic bridge lemma \cref{lem:sheafifySnd_map_eq} the same way.
\end{lemma}

\begin{lemma}

\leanok
[Naturality of the rearrangement; the twist endofunctor isomorphism]
  \label{lem:tensorRearrange_natIso}
  \lean{MR0555258CompactifyingPicard.tensorRearrangeNatIso}
  \uses{lem:tensorRearrange, def:tensorRight, lem:tensorAssocNatIso2, lem:tensorAssocNatIso3, def:tensorBraidingNatIso}
  For fixed $\Ox$-modules $F, L$ the rearrangement \cref{lem:tensorRearrange} is natural
  in $Y$, packaging into a natural isomorphism of endofunctors on $\Ox$-modules
  \[
    \texttt{tensorLeft}\,F \ \circ\ \texttt{tensorRight}\,L
      \;\cong\; \texttt{tensorLeft}\,(F \tensor_{\Ox} L),
    \qquad Y \longmapsto \big((F \tensor Y) \tensor L \cong (F \tensor L) \tensor Y\big).
  \]
  \begin{proof}
    Compose the three natural steps: the middle-factor associator naturality
    \cref{lem:tensorAssocNatIso2}, the inner braiding naturality \cref{def:tensorBraidingNatIso}
    (whiskered into $\texttt{tensorLeft}\,F$, swapping $Y$ and $L$), and the last-factor
    associator naturality \cref{lem:tensorAssocNatIso3}. Concretely
    $\texttt{tensorRearrangeNatIso} := \texttt{tensorAssocNatIso2} \ggg
    \texttt{isoWhiskerRight}\,(\texttt{tensorBraidingNatIso}\,L)\,(\texttt{tensorLeft}\,F) \ggg
    \texttt{tensorAssocNatIso3}$.
  \end{proof}
\end{lemma}

\begin{definition}

\leanok
[Index functors for first-factor associator whiskering]
  \label{def:assocIdxL1}
  \lean{MR0555258CompactifyingPicard.assocIdxL₁}
  \uses{def:tensorMod}
  The private index functor $A \mapsto (A.\mathrm{val} \tensor L.\mathrm{val},\ L'.\mathrm{val})$
  for the first-factor associator naturality \cref{lem:tensorAssocNatIso1} (the first product
  slot now varies, the inner factors $L, L'$ are fixed). Built analogously to
  \cref{def:assocIdxL}.
\end{definition}

\begin{definition}

\leanok
[Index functor, right form, first-factor]
  \label{def:assocIdxR1}
  \lean{MR0555258CompactifyingPicard.assocIdxR₁}
  \uses{def:tensorMod}
  The private index functor $A \mapsto (A.\mathrm{val},\ L.\mathrm{val} \tensor L'.\mathrm{val})$,
  the right-form companion of \cref{def:assocIdxL1} for \cref{lem:tensorAssocNatIso1}.
\end{definition}

\begin{lemma}

\leanok
[Naturality of the left unitor in the tensored object]
  \label{lem:tensorLeftUnitor_natIso}
  \lean{MR0555258CompactifyingPicard.tensorLeftUnitorNatIso}
  \uses{def:moduleSheafification, lem:tensorMod_unitor}
  \textit{Project-local Lean realization (supporting \cref{lem:tensorRightUnitor_natIso}).}
  The left unitor $\mathcal{O}_X \tensor_{\Ox} A \cong A$ (\cref{lem:tensorMod_unitor}) is
  natural in $A$, packaging into a natural isomorphism of endofunctors on $\Ox$-modules
  \[
    \texttt{tensorLeft}\,(\mathcal{O}_X) \;\cong\; \mathbf{1}_{\Ox\text{-Mod}}.
  \]
  \begin{proof}
    The component at $A$ is the object-level left unitor. Naturality in $A$ combines the
    presheaf-level naturality of Mathlib's \texttt{MonoidalCategory.leftUnitor}
    (whose $\texttt{whiskerLeft}\,(\mathbf{1})\,\varphi$ is definitionally the
    $\texttt{tensorHom}$ used here) with the naturality of the sheafification counit
    (\cref{def:moduleSheafification}). Reassociation of the resulting sheafification square is
    discharged by a pure term-mode \texttt{Eq.trans} chain (the keyed \texttt{rw}/\texttt{simp}
    reassociation silently fails on these non-syntactic $\circ$ paths).
  \end{proof}
\end{lemma}

\begin{lemma}

\leanok
[Naturality of the right unitor in the tensored object]
  \label{lem:tensorRightUnitor_natIso}
  \lean{MR0555258CompactifyingPicard.tensorRightUnitorNatIso}
  \uses{def:tensorRight, lem:tensorLeftUnitor_natIso, def:tensorBraidingNatIso}
  The right unitor is natural in $A$, packaging into a natural isomorphism of endofunctors on
  $\Ox$-modules
  \[
    \mathbf{1}_{\Ox\text{-Mod}} \;\cong\; \texttt{tensorRight}\,(\mathcal{O}_X),
    \qquad A \longmapsto (A \cong A \tensor_{\Ox} \mathcal{O}_X).
  \]
  \begin{proof}
    One-liner: compose the natural braiding \cref{def:tensorBraidingNatIso} (at the unit) with
    the left-unitor natural isomorphism \cref{lem:tensorLeftUnitor_natIso} and take the inverse,
    $\texttt{tensorRightUnitorNatIso} := (\texttt{tensorBraidingNatIso}\,\mathcal{O}_X \ggg
    \texttt{tensorLeftUnitorNatIso})^{-1}$. Being a composite of natural isomorphisms it is
    natural in $A$.
  \end{proof}
\end{lemma}

\begin{lemma}

\leanok
[Naturality of $\texttt{tensorAssoc}$ in the first factor]
  \label{lem:tensorAssocNatIso1}
  \lean{MR0555258CompactifyingPicard.tensorAssocNatIso₁}
  \uses{lem:tensorMod_assoc_uncond, lem:sheafifyTensorComparison_uncond, lem:tensorRight_map_eq, lem:tensorThenSheafify_map_eq, lem:presheaf_id_tensorHom_id, lem:sheafify_map_id_val, def:assocIdxL1, def:assocIdxR1}
  For fixed $\Ox$-modules $L, L'$ the associator \cref{lem:tensorMod_assoc_uncond} is natural
  in its \emph{first} factor $A$: a natural isomorphism of endofunctors on $\Ox$-modules
  \[
    \texttt{tensorRight}\,(L \tensor_{\Ox} L') \;\cong\;
      \texttt{tensorRight}\,L \ \circ\ \texttt{tensorRight}\,L',
    \qquad A \longmapsto (A \tensor (L \tensor L') \cong (A \tensor L) \tensor L').
  \]
  \begin{proof}
    The first-factor mirror of \cref{lem:tensorAssocNatIso2}/\cref{lem:tensorAssocNatIso3}.
    Reuse the same three-\texttt{NatIso} decomposition, now with the varying factor in the
    first product slot, so the leg comparison morphism is $(\varphi.\mathrm{val},\
    \id_{L.\mathrm{val} \tensor L'.\mathrm{val}})$ and \emph{both} comparison legs sheafify
    the fixed inner factors $L, L'$, discharged by the sheafify-of-identity bridges
    \cref{lem:presheaf_id_tensorHom_id}, \cref{lem:sheafify_map_id_val} (the
    $\texttt{convert}\,\dots\,\texttt{using 4}$ trick used for \cref{lem:tensorAssocNatIso3}'s
    fixed leg). Index functors \cref{def:assocIdxL1}, \cref{def:assocIdxR1}.
  \end{proof}
\end{lemma}

\begin{lemma}

\leanok
[Whiskering an object isomorphism into $\texttt{tensorRight}$]
  \label{lem:tensorRightMapIso}
  \lean{MR0555258CompactifyingPicard.tensorRightMapIso}
  \uses{def:tensorRight, lem:tensorRight_map_eq, lem:tensorLeft_map_eq}
  \textit{Project-local Lean realization (supporting \cref{lem:tensorRight_invertible_equiv}).}
  An isomorphism $e : L_1 \cong L_2$ of $\Ox$-modules induces a natural isomorphism of
  endofunctors $\texttt{tensorRight}\,L_1 \cong \texttt{tensorRight}\,L_2$, with component at $A$
  the object isomorphism $A \tensor_{\Ox} L_1 \cong A \tensor_{\Ox} L_2$.
  \begin{proof}
    The component at $A$ is $(\texttt{tensorLeft}\,A).\texttt{mapIso}\,e$. Naturality in $A$ is
    the tensor interchange law: the presheaf-level square is closed pointwise by
    \texttt{MonoidalCategory.whisker\_exchange}, then sheafified through the generic bridge
    lemmas \cref{lem:tensorRight_map_eq}, \cref{lem:tensorLeft_map_eq}. (Note
    \texttt{tensorHom\_comp\_tensorHom} does \emph{not} fire on the hand-built tensor;
    \texttt{whisker\_exchange} is the working route.)
  \end{proof}
\end{lemma}

\begin{lemma}

\leanok
[An invertible sheaf gives a tensor equivalence]
  \label{lem:tensorRight_invertible_equiv}
  \lean{MR0555258CompactifyingPicard.tensorRightEquivOfInvertible}
  \uses{def:tensorRight, def:isInvertibleMod, lem:tensorMod_assoc_uncond, lem:tensorMod_unitor, lem:tensorMod_braiding, lem:tensorRightUnitor_natIso, lem:tensorAssocNatIso1, lem:tensorRightMapIso}
  If $L$ is invertible (\cref{def:isInvertibleMod}), with inverse $L'$ and
  $L \tensor_{\Ox} L' \cong \Ox$, then $\texttt{tensorRight}\,L$ is an equivalence of
  categories on $\Ox$-modules, with quasi-inverse $\texttt{tensorRight}\,L'$. In
  particular it is fully faithful, so it induces a bijection on Hom-sets
  $\Hom_X(A,B) \cong \Hom_X(A \tensor L, B \tensor L)$, natural in $A$ and $B$.
  \begin{proof}
    The unit/counit are the natural isomorphisms
    $A \cong (A \tensor L) \tensor L'$ and $(A \tensor L') \tensor L \cong A$ built from the
    \emph{first-factor} associator naturality \cref{lem:tensorAssocNatIso1}, the chosen iso
    $L \tensor L' \cong \Ox$ (and its braided companion $L' \tensor L \cong \Ox$), and the
    right-unitor naturality \cref{lem:tensorRightUnitor_natIso}; being composites of natural
    isomorphisms they are automatically natural in $A$. Concretely the unit is
    $\texttt{tensorRightUnitorNatIso} \ggg (-\tensor(L\tensor L')\text{ whiskered by }
    (L\tensor L'\cong\Ox)^{-1}) \ggg \texttt{tensorAssocNatIso1}$, and the counit is its
    mirror. Assemble with \texttt{CategoryTheory.Equivalence.mk} (no triangle/coherence
    obligation, verified iter-032). A fully faithful functor induces the stated Hom-bijection
    via \texttt{Equivalence.fullyFaithfulFunctor.homEquiv}.
  \end{proof}
\end{lemma}

\begin{lemma}

\leanok
[The $(1.1.2)$ twist isomorphism of $\homTensorFunctor$s]
  \label{lem:homTensorFunctor_twist_natIso}
  \lean{MR0555258CompactifyingPicard.homTensorFunctorTwistIso}
  \uses{def:homTensorFunctor, lem:tensorRearrange_natIso, lem:tensorRight_invertible_equiv, def:tensorBC}
  For $L$ invertible there is a natural isomorphism of functors $\Os\text{-Mod} \to \mathrm{Type}$
  \[
    \homTensorFunctor(f, I \tensor L, F \tensor L)
      \;\cong\; \homTensorFunctor(f, I, F).
  \]
  \begin{proof}
    Both functors factor through the pullback $f^*$, so it suffices to give a natural
    isomorphism of the $\Ox$-side functors $Y \mapsto \Hom_X(I \tensor L, (F \tensor L) \tensor Y)$
    and $Y \mapsto \Hom_X(I, F \tensor Y)$. Rearranging the target by
    \cref{lem:tensorRearrange_natIso} gives $(F \tensor L) \tensor Y \cong (F \tensor Y) \tensor L$,
    natural in $Y$; the tensor equivalence \cref{lem:tensorRight_invertible_equiv} then gives
    $\Hom_X(I \tensor L, (F \tensor Y) \tensor L) \cong \Hom_X(I, F \tensor Y)$, natural in
    $Y$. Whisker by $f^*$.
  \end{proof}
\end{lemma}

\begin{lemma}\leanok
[Sheafification commutes with the presheaf tensor (conditional, right)]
  \label{lem:sheafifyTensorComparison}
  \lean{MR0555258CompactifyingPicard.sheafifyTensorComparisonOfLocallyBijective}
  \uses{def:tensorMod, lem:sheafification_map_isIso, def:sheafifyTensorUnitWhiskerRight}
  For presheaves of modules $P, Q$, \emph{given} that the whiskered unit
  $\eta_P \tensor \id_Q$ is locally injective and locally surjective, the induced
  map exhibits an isomorphism
  $\widetilde{\,(\widetilde{P}) \tensor Q\,} \cong \widetilde{P \tensor Q}$.
  \begin{proof}
    \uses{lem:sheafification_map_isIso, def:sheafifyTensorUnitWhiskerRight}
    The comparison map is the sheafification of the whiskered unit
    $\eta_P \tensor \id_Q : P \tensor Q \to (\widetilde{P}) \tensor Q$
    (\cref{def:sheafifyTensorUnitWhiskerRight}). By the local-bijectivity
    hypotheses this whiskered map is a local isomorphism, so by
    \cref{lem:sheafification_map_isIso} its sheafification is an isomorphism; take
    its inverse.

    \textbf{Why the hypotheses cannot be discharged (iter-008 finding).} One would
    like ``tensoring by a fixed $Q$ sends local isomorphisms to local
    isomorphisms.'' This holds for local \emph{surjectivity} (right exactness of
    $-\tensor Q$ on sections), but is \emph{false} for local \emph{injectivity}:
    $-\tensor Q$ is not left exact (e.g.\ $\mathbb{Z}/2 \tensor (\cdot 2:
    \mathbb{Z}\to\mathbb{Z})$ is not injective). At the presheaf level the
    injectivity half genuinely fails; it is recovered only stalkwise or via a
    $\texttt{MonoidalClosed (PresheafOfModules }R)$ instance, neither of which is in
    Mathlib. Hence the hypotheses are left explicit. The unconditional engine
    \cref{lem:sheafification_map_isIso} \emph{is} proven outright.
  \end{proof}
\end{lemma}

\begin{lemma}\leanok
[Sheafification--tensor comparison, left handedness (conditional)]
  \label{lem:sheafifyTensorComparisonLeft}
  \lean{MR0555258CompactifyingPicard.sheafifyTensorComparisonLeftOfLocallyBijective}
  \uses{def:tensorMod, lem:sheafification_map_isIso, def:sheafifyTensorUnitWhiskerLeft}
  Symmetric to \cref{lem:sheafifyTensorComparison} on the left factor: given that
  $\id_P \tensor \eta_Q$ is locally bijective, the induced map exhibits
  $\widetilde{\,P \tensor (\widetilde{Q})\,} \cong \widetilde{P \tensor Q}$.
  \begin{proof}
    \uses{lem:sheafification_map_isIso, def:sheafifyTensorUnitWhiskerLeft}
    Identical to \cref{lem:sheafifyTensorComparison} with the whiskering on the left
    factor (\cref{def:sheafifyTensorUnitWhiskerLeft}).
  \end{proof}
\end{lemma}

\begin{lemma}\leanok
[Sheafification inverts locally bijective maps]
  \label{lem:sheafification_map_isIso}
  \lean{MR0555258CompactifyingPicard.sheafification_map_isIso_of_locallyBijective}
  For a locally bijective morphism of sheaves of rings and a presheaf-of-modules
  map $g$ that is locally injective and locally surjective, the sheafification of
  $g$ is an isomorphism. This is the unconditional ``sheafify inverts local isos''
  engine underlying both comparison lemmas.
  \begin{proof}
    Module sheafification $\widetilde{(-)}$ is a localization functor at the class
    $J.W$ of locally bijective maps (\texttt{ModuleCat.Sheaf.Localization}). A
    locally injective + locally surjective $g$ lies in $J.W$
    (\texttt{W\_of\_isLocallyBijective}, via \texttt{WEqualsLocallyBijective}), and a
    localization functor inverts its class (\texttt{Localization.inverts}); hence
    $\widetilde{g}$ is an isomorphism.
  \end{proof}
\end{lemma}

\begin{definition}\leanok
[Sheafification unit on a presheaf of modules]
  \label{def:sheafifyUnit}
  \lean{MR0555258CompactifyingPicard.sheafifyUnit}
  \uses{def:tensorMod}
  For a presheaf of modules $P$, the unit $\eta_P : P \to (\widetilde{P}).\mathrm{val}$
  of the module-sheafification adjunction (\texttt{PresheafOfModules.toSheafify}).
\end{definition}

\begin{definition}\leanok
[Whiskered sheafification unit, right]
  \label{def:sheafifyTensorUnitWhiskerRight}
  \lean{MR0555258CompactifyingPicard.sheafifyTensorUnitWhiskerRight}
  \uses{def:sheafifyUnit}
  The map $\eta_P \tensor \id_Q : P \tensor Q \to (\widetilde{P}) \tensor Q$
  obtained by whiskering the sheafification unit \cref{def:sheafifyUnit} on the
  right by $Q$ (\texttt{Monoidal.whiskerRight}).
\end{definition}

\begin{definition}\leanok
[Whiskered sheafification unit, left]
  \label{def:sheafifyTensorUnitWhiskerLeft}
  \lean{MR0555258CompactifyingPicard.sheafifyTensorUnitWhiskerLeft}
  \uses{def:sheafifyUnit}
  The map $\id_P \tensor \eta_Q : P \tensor Q \to P \tensor (\widetilde{Q})$
  obtained by whiskering \cref{def:sheafifyUnit} on the left by $P$.
\end{definition}

\begin{lemma}\leanok
[Unitor of the $\Ox$-tensor]
  \label{lem:tensorMod_unitor}
  \lean{MR0555258CompactifyingPicard.tensorLeftUnitor}
  \uses{def:tensorMod, def:presheafLeftUnitor, def:sheafifyValIso}
  For an $\Ox$-module $A$ there is a natural isomorphism
  $\mathcal{O}_X \tensor_{\Ox} A \cong A$ with the structure sheaf
  (\texttt{SheafOfModules.unit}) as tensor unit; the right unitor
  $A \tensor_{\Ox} \mathcal{O}_X \cong A$ follows by \cref{lem:tensorMod_braiding}.
  \begin{proof}
    Sheafify the presheaf-level left unitor (\cref{def:presheafLeftUnitor}, the
    pointwise $R(U) \tensor_{R(U)} M \cong M$, \texttt{TensorProduct.lid}), then
    identify the sheafification of the sheaf $A$ with $A$ itself
    (\cref{def:sheafifyValIso}).
  \end{proof}
\end{lemma}

\begin{definition}

\leanok
[Presheaf-level left unitor]
  \label{def:presheafLeftUnitor}
  \lean{MR0555258CompactifyingPicard.presheafLeftUnitor}
  The presheaf-level left unitor
  $\mathbf{1} \tensor M \cong M$ of \texttt{PresheafOfModules.Monoidal.tensorObj}.
  The monoidal unit of \texttt{PresheafOfModules} is definitionally
  \texttt{PresheafOfModules.unit}, so this is Mathlib's monoidal
  \texttt{leftUnitor}, restated so that its unit matches the \texttt{.val} of
  \texttt{SheafOfModules.unit} used by \cref{def:tensorMod}.
\end{definition}

\begin{definition}

\leanok
[Sheafification of a sheaf's underlying presheaf]
  \label{def:sheafifyValIso}
  \lean{MR0555258CompactifyingPicard.sheafifyValIso}
  For an $\Ox$-module $A$ (a sheaf of modules), the sheafification of its underlying
  presheaf is canonically isomorphic to $A$ itself, $\widetilde{A.\mathrm{val}}
  \cong A$. This is the counit of the reflective module-sheafification adjunction at
  the sheaf $A$, an isomorphism because the adjunction is reflective. A reusable
  building block for the unitor and associator.
\end{definition}

\begin{lemma}

\leanok
[Right unitor of the $\Ox$-tensor]
  \label{lem:tensorMod_rightUnitor}
  \lean{MR0555258CompactifyingPicard.tensorRightUnitor}
  \uses{lem:tensorMod_braiding, lem:tensorMod_unitor}
  For an $\Ox$-module $A$ there is a natural isomorphism
  $A \tensor_{\Ox} \mathcal{O}_X \cong A$.
  \begin{proof}
    Compose the braiding $A \tensor \mathcal{O}_X \cong \mathcal{O}_X \tensor A$
    (\cref{lem:tensorMod_braiding}) with the left unitor (\cref{lem:tensorMod_unitor}).
  \end{proof}
\end{lemma}

\begin{definition}\leanok
[The ``tensor by $F$'' endofunctor]
  \label{def:tensorLeft}
  \lean{MR0555258CompactifyingPicard.tensorLeft}
  \uses{def:tensorMod}
  For an $\Ox$-module $F$, the endofunctor $\mathrm{tensorLeft}\,F$ on
  $\Ox$-modules sends $N \mapsto F \tensor_{\Ox} N$ (\cref{def:tensorMod}) and a
  morphism $\varphi : N \to N'$ to the sheafification of the presheaf tensor
  $\mathrm{id}_F \tensor \varphi$. It is the functorial form of \cref{def:tensorMod}
  in its second argument; no \texttt{MonoidalCategory} instance on
  $X.\mathrm{Modules}$ is required.
\end{definition}

\begin{lemma}

\leanok
[Sheafification of presheaves of modules is additive]
  \label{lem:sheafification_additive}
  \lean{MR0555258CompactifyingPicard.presheafOfModules_sheafification_additive}
  Mathlib's sheafification functor on presheaves of modules
  (\texttt{PresheafOfModules.sheafification}) is additive: it preserves the
  addition of morphisms.
  \begin{proof}
    Sheafification is a left adjoint, hence preserves the (finite) biproducts that
    compute addition of morphisms in a preadditive category; a functor between
    preadditive categories preserving binary biproducts is additive. (Project-local
    Mathlib supplement: \texttt{PresheafOfModules.sheafification} carries
    \texttt{PreservesZeroMorphisms} but not yet \texttt{Additive}.)
  \end{proof}
\end{lemma}

\begin{lemma}

\leanok
[$\mathrm{tensorLeft}\,F$ is additive]
  \label{lem:tensorLeft_additive}
  \lean{MR0555258CompactifyingPicard.tensorLeft_additive}
  \uses{def:tensorLeft, lem:sheafification_additive}
  The endofunctor $\mathrm{tensorLeft}\,F$ (\cref{def:tensorLeft}) is additive.
  \begin{proof}
    On the presheaf level $\mathrm{id}_F \tensor (-)$ is additive because
    $\mathrm{ModuleCat}$ is monoidal-preadditive
    ($\mathrm{tensorHom}\,(\mathrm{id}_F)\,(\varphi + \psi) =
    \mathrm{tensorHom}\,(\mathrm{id}_F)\,\varphi +
    \mathrm{tensorHom}\,(\mathrm{id}_F)\,\psi$ pointwise), and sheafification is
    additive (\cref{lem:sheafification_additive}); the composite is additive.
  \end{proof}
\end{lemma}

\begin{definition}\leanok
[Inhabitedness of $\Ox$-modules]
  \label{def:nonempty_modules}
  \lean{MR0555258CompactifyingPicard.instNonemptyModules}
  The category of $\Ox$-modules is inhabited, the zero module being an object.
  This is a technical witness (edge-free by design) needed to state the
  placeholder declarations whose result type is an $\Ox$-module.
\end{definition}

\begin{definition}

\leanok
[Free $\Ox$-module]
  \label{def:isFreeMod}
  \lean{MR0555258CompactifyingPicard.IsFreeMod}
  An $\Ox$-module $J$ is \emph{free} when $J \cong \mathcal{O}_X^{(\Lambda)}$ for
  some index set $\Lambda$. Realized via Mathlib's \texttt{SheafOfModules.free}
  (existence of an iso to a coproduct of copies of the unit).
\end{definition}

\begin{definition}

\leanok
[Invertible $\Ox$-module]
  \label{def:isInvertibleMod}
  \uses{def:tensorMod}
  \lean{MR0555258CompactifyingPicard.IsInvertibleMod}
  An $\Ox$-module $L$ is \emph{invertible} when it is locally free of rank $1$,
  equivalently when there is $L^{-1}$ with $L \tensor_{\Ox} L^{-1} \cong
  \mathcal{O}_X$. Realized as an invertible object for the tensor of
  \cref{def:tensorMod}.
\end{definition}

% --- Project-local (Mathlib-gap) predicates of strands (a)/(b). These are the
% "opaque placeholder" Props of the §1 statements; the bottom-up P1b-iii(a) layer
% realizes the tractable batch as genuine definitions. The fibrewise condition
% FiberExtVanishes is realized via the scheme fibre X(s) (Mathlib's
% Scheme.Hom.fiber) and the absolute Ext group; the relative-Ext sheaf relExt and
% the LFP+flat-on-an-open predicate IsLFPAndFlatOn stay opaque pending the
% relative-Ext / proper-pushforward (R^q f_*) infrastructure absent from Mathlib. ---

\begin{definition}\leanok
[Retrocompact open]
  \label{def:isRetrocompact}
  \lean{MR0555258CompactifyingPicard.IsRetrocompact}
  An open $U \subseteq S$ is \emph{retrocompact} when its inclusion is
  quasi-compact (equivalently, $U \cap V$ is quasi-compact for every
  quasi-compact open $V$). Realized via Mathlib's quasi-compactness of the open
  immersion $U \hookrightarrow S$.
\end{definition}

\begin{definition}\leanok
[Locally free of finite rank on an open]
  \label{def:isLocallyFreeOfFiniteRankOn}
  \lean{MR0555258CompactifyingPicard.IsLocallyFreeOfFiniteRankOn}
  \uses{def:isFreeMod}
  An $\Os$-module $M$ is \emph{locally free of finite rank on} an open
  $U \subseteq S$ when the restriction $M|_U$ is, locally on $U$, isomorphic to a
  finite free $\mathcal{O}_U$-module. Realized via Mathlib's
  \texttt{SheafOfModules.free} on the restriction together with a finiteness
  condition on the index type.
\end{definition}

\begin{definition}[Absolute Ext group $\Ext^q_Y(M,N)$]\leanok
  \label{def:extGroup}
  \lean{MR0555258CompactifyingPicard.extGroup}
  The (absolute) Ext group $\Ext^q_Y(M,N)$ of two $\mathcal{O}_Y$-modules, an
  abelian group, used to phrase acyclicity in \cref{lem:acyclic_surjection}.
  Realized via Mathlib's \texttt{CategoryTheory.Abelian.Ext} in the abelian
  category $Y.\mathrm{Modules}$ (the existing \texttt{Scheme.Modules.instAbelian}
  instance), endowing it with the standard derived category
  (\texttt{HasDerivedCategory.standard}) and the resulting
  \texttt{hasExt\_of\_hasDerivedCategory} structure, packaged as an object of
  \texttt{Ab}. The Ext universe is the object universe of $Y.\mathrm{Modules}$,
  so the group lands in $\mathtt{Ab}.\{u{+}1\}$; the signature is taken there
  (every use site phrases only $\mathtt{IsZero}(\cdot)$, which is
  universe-agnostic, so no dependent constrains the universe). This avoids the
  \texttt{IsGrothendieckAbelian (SheafOfModules R)} instance entirely (its
  \texttt{AB5OfSize}/\texttt{HasSeparator} fields are absent from Mathlib).
\end{definition}

\begin{definition}[Sheaf cohomology $H^q(U,F)$]\leanok
  \label{def:sheafCohomology}
  \lean{MR0555258CompactifyingPicard.sheafCohomology}
  \uses{def:extGroup}
  The sheaf cohomology $H^q(U,F)$ of an $\mathcal{O}_U$-module $F$, an abelian
  group, used only to state affine acyclicity (\cref{thm:ega_ext_affine_acyclic}).
  Realized as $\Ext^q_U(\mathcal{O}_U, F)$ (the derived functor of global sections
  expressed through \cref{def:extGroup}), with $\mathcal{O}_U$ the unit
  $\mathcal{O}_U$-module (\texttt{SheafOfModules.unit}). Lands in
  $\mathtt{Ab}.\{u{+}1\}$, inheriting the universe of \cref{def:extGroup}.
\end{definition}

\begin{definition}[Internal hom sheaf $\mathcal{H}om_{\Ox}(I,-)$]
  \label{def:internalHom}
  \lean{MR0555258CompactifyingPicard.internalHom}
  \textit{Project-local (Mathlib gap): the sheaf-Hom functor, foundational
  infrastructure for the relative local Ext sheaf of Altman--Kleiman, §1.}
  For a fixed $\Ox$-module $I$, the internal hom functor
  $\mathcal{H}om_{\Ox}(I,-) : \mathbf{Mod}_{\Ox} \to \mathbf{Mod}_{\Ox}$ sends an
  $\Ox$-module $M$ to the sheaf-Hom $\mathcal{H}om_{\Ox}(I,M)$, the
  $\Ox$-module whose sections over an open $U$ are the $\Ox|_U$-linear maps
  $I|_U \to M|_U$. It is additive and left exact (a right adjoint of $-\tensor I$
  at the presheaf level; the section presheaf $U \mapsto (I|_U \to M|_U)$ is already
  a sheaf when $M$ is, so no sheafification is needed — in contrast to $\tensor$).
  This functor is absent from Mathlib v4.30.0 (no \texttt{internalHom}/\texttt{sheafHom}
  for \texttt{SheafOfModules}); it is built project-side from the
  \texttt{PresheafOfModules} restriction functors and the module structure on
  morphism sets, and is the first foundational brick for \cref{def:relExt}.
\end{definition}

\begin{definition}[Relative local Ext sheaf $\sExt^q_{X/S}(I,F)$]
  \label{def:relExt}
  \lean{MR0555258CompactifyingPicard.relExt}
  \uses{def:internalHom}
  \textit{Project-local realization of the notation $\sExt^q_{X/S}(I,F)$ of Altman--Kleiman, §1.}
  The relative local Ext sheaf $\sExt^q_{X/S}(I,F)$, an $\Os$-module, defined as the
  $q$-th right derived functor of the composite
  $F \mapsto f_*\,\mathcal{H}om_{\Ox}(I,F)$ — the internal hom
  \cref{def:internalHom} followed by the pushforward
  $f_* : \mathbf{Mod}_{\Ox} \to \mathbf{Mod}_{\Os}$ (Mathlib's
  \texttt{SheafOfModules.pushforward}). \emph{Currently opaque}: a faithful
  realization requires \cref{def:internalHom} (under construction) and the right
  derived functor $R^q$ of the composite, the latter reusing the derived-category
  setup already used by \cref{def:extGroup}
  (\texttt{HasDerivedCategory.standard} on \texttt{X.Modules}). The declaration
  carries the correct type so the §1 statements that mention it can be transcribed,
  and is realized once \cref{def:internalHom} and the derived-functor brick land.
\end{definition}

\begin{definition}[Fibrewise Ext-vanishing $\Ext^q_{X(s)}(I(s),F(s)) = 0$]\leanok
  \label{def:fiberExtVanishes}
  \lean{MR0555258CompactifyingPicard.FiberExtVanishes}
  \uses{def:extGroup}
  \textit{Project-local realization of the fibrewise condition $\Ext^q_{X(s)}(I(s),F(s))=0$ used throughout Altman--Kleiman, §1 (e.g. (1.3), (1.9), (1.10); the (1.3) instance is quoted verbatim in \cref{thm:H_locallyFree}).}
  The fibrewise vanishing condition $\Ext^q_{X(s)}(I(s),F(s)) = 0$ at a point
  $s \in S$, where $X(s)$ is the scheme-theoretic fibre of $f$ over $s$ and
  $I(s), F(s)$ are the restrictions of $I, F$ to $X(s)$. Realized faithfully: take
  the fibre $X(s) = $ \texttt{Scheme.Hom.fiber}\,$f\,s$ with its inclusion
  $j = $ \texttt{fiber\(\iota\)}\,$f\,s : X(s) \to X$ (both from Mathlib's
  \texttt{AlgebraicGeometry.Fiber}); set $I(s) = j^* I$, $F(s) = j^* F$ via
  \texttt{Scheme.Modules.pullback}; the condition is
  $\mathtt{IsZero}\big(\Ext^q_{X(s)}(I(s), F(s))\big)$ through the absolute Ext
  group of \cref{def:extGroup} on the fibre scheme. This replaces the former opaque
  placeholder by a genuine definition built from existing Mathlib + project pieces.
\end{definition}

\begin{definition}[Locally finitely presented and flat on an open]
  \label{def:isLFPAndFlatOn}
  \lean{MR0555258CompactifyingPicard.IsLFPAndFlatOn}
  \uses{def:relExt, def:isLFP, def:isSFlat}
  \textit{Project-local realization of the conclusion of Altman--Kleiman, §1, (1.10)(i).}
  The predicate that the restriction of the relative Ext sheaf
  $\sExt^c_{X/S}(I,F)$ (\cref{def:relExt}) to an open $V \subseteq S$ is locally
  finitely presented and $S$-flat. \emph{Currently opaque}: it references
  \cref{def:relExt}, whose realization is gated on the absent $R^q f_*$
  infrastructure; the declaration carries the correct type so (1.10)(i) can be
  transcribed, and is realized together with \cref{def:relExt}.
\end{definition}

\begin{definition}\leanok
[Half-exact additive functor]
  \label{def:isHalfExact}
  \lean{MR0555258CompactifyingPicard.IsHalfExact}
  An additive endofunctor $T$ of finitely generated $A$-modules is
  \emph{half-exact} when, for every short exact sequence
  $0 \to A' \to B \to C \to 0$, the middle sequence $T(A') \to T(B) \to T(C)$ is
  exact at $T(B)$. Realized directly as exactness (at the middle) of the image
  short complex under $T$ for every short exact \texttt{ShortComplex}; no packaged
  Mathlib predicate exists, so the definition is given explicitly.
  In Lean the functor is typed $T : \mathrm{Mod}_A^{(u)} \to \mathrm{Mod}_A^{(u+1)}$
  (the value category one universe up), matching the $\Ext$-valued instance of
  \cref{def:functorT}; see the formalization note of \cref{thm:ob_halfexact_free}.
  The image short complex $T(A')\to T(B)\to T(C)$ then lives in
  $\mathrm{Mod}_A^{(u+1)}$, where exactness is taken.
\end{definition}

\begin{definition}\leanok
[Admissibility hypotheses of (1.1)]
  \label{def:admissible}
  \lean{MR0555258CompactifyingPicard.Admissible}
  \uses{def:isLFP, def:isSFlat}
  \textit{Source: Altman--Kleiman, §1, (1.1), p. 54 (standing hypotheses; verbatim quoted in \cref{def:H}).}
  The data $(f, I, F)$ are \emph{admissible} when $f : X \to S$ is proper and
  locally of finite presentation, $I$ and $F$ are locally finitely presented
  $\Ox$-modules, and $F$ is flat over $S$. These are the standing hypotheses of
  (1.1) under which $\HIF(I,F)$ and $h(I,F)$ are defined.
\end{definition}

\begin{definition}

\leanok
[Base-change tensor functor $M \mapsto F \tensor_S M$]
  \label{def:tensorBaseChangeFunctor}
  \lean{MR0555258CompactifyingPicard.tensorBaseChangeFunctor}
  \uses{def:tensorMod, def:tensorLeft}
  \textit{Source: Altman--Kleiman, §1, (1.1), p. 54 (the functor $M \mapsto F \tensor_S M$).}
  For $f : X \to S$ and an $\Ox$-module $F$, the \emph{base-change tensor functor}
  is the covariant functor $\Os\text{-mod} \to \Ox\text{-mod}$ sending an
  $\Os$-module $M$ to $F \tensor_{\Ox} f^* M$, the tensor over $\Ox$ of $F$ with
  the pullback $f^* M$. Realized as the composite of Mathlib's pullback functor
  $f^* = $ \texttt{Scheme.Modules.pullback}\,$f : \Os\text{-mod} \to \Ox\text{-mod}$
  with the ``tensor by $F$'' endofunctor on $\Ox$-modules built from the presheaf
  monoidal tensor (\cref{def:tensorMod}, in its functorial form
  \texttt{Monoidal.tensorHom}) followed by sheafification. (No
  \texttt{MonoidalCategory} instance on $X.\mathrm{Modules}$ is required.)
\end{definition}

\begin{lemma}\leanok
[$\mathrm{tensorLeft}\,F$ preserves zero morphisms]
  \label{lem:tensorLeft_preservesZeroMorphisms}
  \lean{MR0555258CompactifyingPicard.tensorLeft_preservesZeroMorphisms}
  \uses{def:tensorLeft}
  The endofunctor $\mathrm{tensorLeft}\,F$ (\cref{def:tensorLeft}) sends the zero
  morphism to the zero morphism. (Technical instance, needed to form the image of
  a short complex under the base-change tensor functor.)
  \begin{proof}
    The presheaf tensor $\mathrm{id}_F \tensor 0$ is zero because $\mathrm{ModuleCat}$
    is monoidal-preadditive ($\mathrm{tensorHom}\,(\mathrm{id}_F)\,0 = 0$ pointwise),
    and sheafification, being additive, preserves the zero morphism.
  \end{proof}
\end{lemma}

\begin{lemma}\leanok
[The base-change tensor functor preserves zero morphisms]
  \label{lem:tensorBaseChangeFunctor_preservesZeroMorphisms}
  \lean{MR0555258CompactifyingPicard.tensorBaseChangeFunctor_preservesZeroMorphisms}
  \uses{def:tensorBaseChangeFunctor, lem:tensorLeft_preservesZeroMorphisms}
  The base-change tensor functor $F \tensor_S (-)$
  (\cref{def:tensorBaseChangeFunctor}) sends the zero morphism to the zero
  morphism. (Technical instance, needed to state the image
  $\mathtt{sc.map}\,(F \tensor_S -)$ of a short complex.)
  \begin{proof}
    It is the composite of \texttt{Scheme.Modules.pullback}\,$f$ (which preserves
    zero morphisms) with $\mathrm{tensorLeft}\,F$
    (\cref{lem:tensorLeft_preservesZeroMorphisms}); a composite of zero-preserving
    functors preserves zero morphisms.
  \end{proof}
\end{lemma}

\begin{definition}\leanok
[Base-change tensor $F \tensor_S M$]
  \label{def:tensorBC}
  \lean{MR0555258CompactifyingPicard.tensorBC}
  \uses{def:tensorBaseChangeFunctor}
  \textit{Source: Altman--Kleiman, §1, (1.1), p. 54 (the notation $F \tensor_S M$).}
  For $f : X \to S$, an $\Ox$-module $F$ and an $\Os$-module $M$, write
  $F \tensor_S M$ for the $\Ox$-module obtained by pulling $M$ back along $f$ and
  tensoring with $F$ over $\Ox$. This is the value at $M$ of the base-change
  tensor functor $M \mapsto F \tensor_S M$ (\cref{def:tensorBaseChangeFunctor})
  from $\Os$-modules to $\Ox$-modules.
\end{definition}

\begin{definition}\leanok
[The corepresented functor $M \mapsto \Hom_X(I, F \tensor_S M)$]
  \label{def:homTensorFunctor}
  \lean{MR0555258CompactifyingPicard.homTensorFunctor}
  \uses{def:tensorBC}
  \textit{Source: Altman--Kleiman, §1, (1.1), p. 54 (the functor corepresented by $H(I,F)$).}
  The covariant functor on quasi-coherent $\Os$-modules
  \[
    M \;\longmapsto\; \Hom_X(I, F \tensor_S M),
  \]
  valued in types. By (1.1) it is corepresentable, and $\HIF(I,F)$ is by
  definition a corepresenting object.
\end{definition}

%==============================================================================
% Strand (a): the representing module H(I,F) and its local-freeness criterion.
%==============================================================================

\begin{definition}\leanok
[The $\Os$-module $H(I,F)$]
  \label{def:H}
  \lean{MR0555258CompactifyingPicard.H}
  \uses{thm:ega_H_existence, def:homTensorFunctor, def:admissible}
  % SOURCE: [Altman--Kleiman], §1, (1.1), p. 54 (read from references/MR0555258-compactifying-picard.pdf)
  % SOURCE QUOTE: "(1.1) (The O_S-Module H(I,F)). Let f : X -> S be a finitely
  % presented, proper morphism of schemes, and let I and F be two locally
  % finitely presented O_X-Modules, with F flat over S. Then there exist a
  % locally finitely presented O_S-Module H(I,F) and an element h(I,F) of
  % Hom_X(I, F (x)_S H(I,F)) which represent the (covariant) functor,
  % M |-> Hom_X(I, F (x)_S M), defined on the category of quasi-coherent
  % O_S-Modules M, and the formation of the pair commutes with base change;"
  \textit{Source: Altman--Kleiman, §1, (1.1), p. 54.}
  Let $f : X \to S$ be a finitely presented, proper morphism of schemes, and let
  $I$ and $F$ be locally finitely presented $\Ox$-modules with $F$ flat over $S$.
  Then $\HIF(I,F)$ denotes the locally finitely presented $\Os$-module that
  represents the covariant functor
  \[
    M \;\longmapsto\; \Hom_X(I, F \tensor_S M)
  \]
  on the category of quasi-coherent $\Os$-modules, whose existence is supplied by
  \cref{thm:ega_H_existence} (over a Noetherian base) together with descent.
\end{definition}

\begin{definition}\leanok
[The universal element $h(I,F)$]
  \label{def:h}
  \lean{MR0555258CompactifyingPicard.h}
  \uses{def:H}
  \textit{Source: Altman--Kleiman, §1, (1.1), p. 54.}
  The universal element $h(I,F) \in \Hom_X\!\big(I, F \tensor_S \HIF(I,F)\big)$ is
  the element corresponding, under the representing isomorphism of \cref{def:H}
  for $M = \HIF(I,F)$, to the identity of $\HIF(I,F)$.
\end{definition}

\begin{theorem}\leanok
[Representability and base-change compatibility of $(H,h)$]
  \label{thm:H_represents}
  \lean{MR0555258CompactifyingPicard.H_represents}
  \uses{def:H, def:h}
  % SOURCE: [Altman--Kleiman], §1, (1.1.1), p. 54 (read from references/MR0555258-compactifying-picard.pdf)
  % SOURCE QUOTE: "in other words, the Yoneda map defined by h(I,F),
  % y : Hom_T(H(I,F)_T, M) -> Hom_{X_T}(I_T, F (x)_S M), (1.1.1)
  % is an isomorphism for every S-scheme T and every quasi-coherent
  % O_T-Module M."
  \textit{Source: Altman--Kleiman, §1, (1.1.1), p. 54.}
  For every $S$-scheme $T$ and every quasi-coherent $\Ot$-module $M$, the Yoneda
  map defined by $h(I,F)$,
  \[
    y : \Hom_T\!\big(\HIF(I,F)_T, M\big) \;\longrightarrow\;
        \Hom_{X_T}\!\big(I_T, F \tensor_S M\big),
        \tag{1.1.1}
  \]
  is an isomorphism. Equivalently, the formation of the pair $(\HIF(I,F), h(I,F))$
  commutes with base change: for $g : T \to S$ one has
  $\HIF(I,F)_T \cong \HIF(I_T, F_T)$ compatibly with the universal elements.
  \begin{proof}\leanok
    \uses{thm:ega_H_existence, thm:ega_descent_noetherian}
    Representability is local on the base $S$, so we may assume $S$ affine; by
    \cref{thm:ega_descent_noetherian} the data $X, I, F$ descend to a finite-type
    $\mathbf{Z}$-scheme $S_0$. Over the Noetherian base $S_0$ the pair
    $(\HIF(I_0, F_0), h(I_0,F_0))$ exists and its formation commutes with
    arbitrary base change (\cref{thm:ega_H_existence}). Base-changing back to $S$
    and using the universal property of $h$ gives the Yoneda isomorphism
    \textnormal{(1.1.1)} for all $T$ and $M$; its naturality in $T$ is exactly the
    asserted base-change compatibility.
  \end{proof}
\end{theorem}

\begin{lemma}\leanok
[Invariance of $H$ under twist by an invertible sheaf]
  \label{lem:H_tensor_invertible}
  \lean{MR0555258CompactifyingPicard.H_tensor_invertible}
  \uses{def:H, thm:H_represents, def:isInvertibleMod, def:tensorMod, def:homTensorFunctor, lem:tensorRearrange, lem:tensorRearrange_natIso, lem:tensorRight_invertible_equiv, lem:homTensorFunctor_twist_natIso, def:coyonedaCompFF}
  % SOURCE: [Altman--Kleiman], §1, (1.1.2), p. 54--55 (read from references/MR0555258-compactifying-picard.pdf)
  % SOURCE QUOTE: "For any invertible O_X-Module L, there is a canonical
  % isomorphism, H(I (x) L, F (x) L) = H(I,F), (1.1.2) because tensoring by L
  % gives a map, Hom_X(I, F (x)_S M) -> Hom_X(I (x) L, (F (x) L) (x)_S M), with
  % an inverse given by tensoring by L^{-1}."
  \textit{Source: Altman--Kleiman, §1, (1.1.2), p. 54--55.}
  For any invertible $\Ox$-module $L$ there is a canonical isomorphism
  \[
    \HIF(I \tensor L,\; F \tensor L) \;=\; \HIF(I,F).
    \tag{1.1.2}
  \]
  \begin{proof}
    \uses{lem:tensorMod_assoc_uncond, lem:tensorMod_braiding, lem:tensorMod_unitor, thm:H_represents}
    The proof is by \emph{anchor-and-assemble} (the same pattern used for (1.2),
    \cref{lem:acyclic_surjection}, and (1.4), \cref{lem:free_resolution_step}): the
    only genuine input absent from Mathlib is the \emph{symmetric-monoidal structure}
    on $\Ox$-modules, which is a true, general fact (sheafification of presheaves of
    modules is monoidal) anchored as \cref{lem:tensorMod_assoc_uncond}; everything
    else is built.

    \emph{Rearrangement.} For an $\Os$-module $M$ write $F \tensor_S M = F \tensor f^*M$
    (\cref{def:tensorBC}). The (unconditional) associator \cref{lem:tensorMod_assoc_uncond}
    together with the braiding \cref{lem:tensorMod_braiding} gives a natural isomorphism
    \[
      r_M : (F \tensor_S M) \tensor L \;=\; (F \tensor f^*M) \tensor L
            \;\cong\; (F \tensor L) \tensor f^*M
            \;=\; (F \tensor L) \tensor_S M,
    \]
    natural in $M$ (it is the $\Ox$-module rearrangement of the three factors
    $F, f^*M, L$).

    \emph{Twist on Hom-sets.} Since $L$ is invertible (\cref{def:isInvertibleMod}),
    tensoring a homomorphism by $L$ is a bijection on Hom-sets, with inverse tensoring
    by $L^{-1}$. Composing with $r_M$ gives, naturally in $M$, a bijection
    \[
      \Phi_M : \Hom_X(I, F \tensor_S M)
        \;\xrightarrow{\;-\,\tensor\,\id_L\;}\;
        \Hom_X(I \tensor L, (F \tensor_S M) \tensor L)
        \;\xrightarrow{\; r_M \circ -\;}\;
        \Hom_X(I \tensor L, (F \tensor L) \tensor_S M),
    \]
    i.e. a natural isomorphism of functors
    $\homTensorFunctor(f, I\tensor L, F\tensor L) \cong \homTensorFunctor(f, I, F)$
    (\cref{def:homTensorFunctor}).

    \emph{Yoneda.} By representability \cref{thm:H_represents} both functors are
    corepresented, by $\HIF(I\tensor L, F\tensor L)$ and $\HIF(I,F)$ respectively;
    since $\mathrm{coyoneda}$ is fully faithful it reflects the isomorphism $\Phi$,
    yielding $\HIF(I\tensor L, F\tensor L) \cong \HIF(I,F)$, which is
    \textnormal{(1.1.2)}.

    \emph{Realization (the naturality layer).} The two displays above are object-level;
    making $r_M$ and $\Phi_M$ \emph{natural in $M$} is the actual Lean content, since
    $\Ox$-modules carry no \texttt{MonoidalCategory} instance. It is supplied by:
    \cref{lem:tensorRearrange_natIso} (the rearrangement $r$ as a natural isomorphism of
    endofunctors, which rests on the sheafification--tensor comparison
    \cref{lem:sheafifyTensorComparison_uncond} being a \emph{natural} isomorphism), and
    \cref{lem:tensorRight_invertible_equiv} (the twist $-\tensor L$ as a categorical
    equivalence from invertibility, giving the natural Hom-bijection). These assemble
    into the natural isomorphism \cref{lem:homTensorFunctor_twist_natIso}, to which the
    Yoneda step is applied.
  \end{proof}
\end{lemma}

\begin{lemma}
[Pullback commutes with tensor ($f^*$ is monoidal) --- Mathlib gap anchor]
  \label{lem:pullbackTensorComparison}
  \lean{MR0555258CompactifyingPicard.External.pullbackTensorComparison}
  \uses{def:tensorMod}
  % NOTE: This is NOT from Altman--Kleiman; it is the standard general fact that the
  % module-pullback functor f^* = Scheme.Modules.pullback f along a morphism of schemes is
  % (strong) monoidal: f^*(A (x)_S B) ~= (f^* A) (x)_X (f^* B), natural in A and B. It is TRUE
  % in full generality -- f^* is the composite forget . presheaf-pullback . sheafification;
  % presheaf base change R (x)_{phi^* S} phi^*(-) is monoidal pointwise, and sheafification is
  % monoidal (the lem:sheafifyTensorComparison_uncond anchor) -- but is ABSENT from Mathlib:
  % neither Scheme.Modules.pullback nor PresheafOfModules.pullback carries any
  % MonoidalFunctor/monoidal structure (loogle-confirmed iter-035). Anchored as a Mathlib gap,
  % exactly parallel to External.sheafifyTensorComparison. It carries NO section-1 content.
  \textit{General fact ($f^*$ is monoidal); anchored as a Mathlib gap. Not from Altman--Kleiman.}
  For a morphism $f : X \to S$ of schemes and $\Os$-modules $A, B$ the module pullback
  $f^* = \texttt{Scheme.Modules.pullback}\,f$ commutes with the tensor product: there is an
  isomorphism of $\Ox$-modules
  \[
    (f^* A) \tensor_{\Ox} (f^* B) \;\cong\; f^*(A \tensor_S B),
  \]
  natural in $A$ and $B$. The Lean realization (re-typed iter-036, following the
  iter-031$\to$032 pattern of \cref{lem:sheafifyTensorComparison_uncond}) is a
  \emph{single-variable} \texttt{NatIso}: for a fixed first slot $A$ it is the isomorphism of
  functors $\Smod \to \Xmod$
  \[
    f^*(-) \,\circ\, \mathtt{tensorLeft}\,(f^* A) \;\cong\;
      \mathtt{tensorLeft}\,A \,\circ\, f^*(-),
  \]
  natural in the second slot $B$ (the only naturality the $(1.1.3)$ comparison ever uses, since
  there the first slot is the fixed $\HIF(I,F)$). The object form
  $(f^* A) \tensor_{\Ox} (f^* B) \cong f^*(A \tensor_S B)$ is recovered as the component at $B$.
  \begin{proof}
    Standard: pullback of modules is strong monoidal (presheaf base change is monoidal
    pointwise; sheafification is monoidal). Anchored as a Mathlib gap (absent from Mathlib;
    see the note above). Used as an external input.
  \end{proof}
\end{lemma}

\begin{definition}

\leanok
[The canonical comparison map of $(1.1.3)$]
  \label{def:htensorComparison}
  \lean{MR0555258CompactifyingPicard.htensorComparison}
  \uses{def:H, def:h, thm:H_represents, def:homTensorFunctor, def:tensorBC, def:tensorRight, lem:tensorMod_assoc_uncond, lem:pullbackTensorComparison}
  \textit{Project-local construction (the canonical map underlying \cref{lem:H_tensor}).}
  The canonical morphism of $\Os$-modules
  \[
    \theta \;:\; \HIF(I \tensor_S N,\, F) \;\longrightarrow\; \HIF(I,F) \tensor_S N .
  \]
  Since $\HIF(I \tensor_S N, F)$ corepresents the functor
  $M \mapsto \Hom_X(I \tensor_S N,\ F \tensor_S M)$ (\cref{thm:H_represents}), giving a map
  \emph{out} of it to $Y := \HIF(I,F) \tensor_S N$ is the same as giving an element of
  $\Hom_X(I \tensor_S N,\ F \tensor_S Y)$ (this is the constructible direction; the final
  isomorphism is its inverse). Take the element
  \[
    I \tensor_S N
      \xrightarrow{\ h(I,F)\,\tensor\,\id_N\ } (F \tensor_S \HIF(I,F)) \tensor_S N
      \xrightarrow{\ \cong\ } F \tensor_S (\HIF(I,F) \tensor_S N) = F \tensor_S Y ,
  \]
  where $h(I,F)$ is the universal element (\cref{def:h}); the first map is
  $\texttt{tensorRight}\,(f^* N)$ applied to $h(I,F)$ (so its source is
  $I \tensor_{\Ox} f^* N = \tensorBC(f, I, N)$), and the isomorphism is the associator
  \cref{lem:tensorMod_assoc_uncond} followed by the $f^*$-monoidality comparison
  \cref{lem:pullbackTensorComparison} applied inside $F \tensor_{\Ox} (-)$ (identifying
  $f^*\HIF(I,F) \tensor_{\Ox} f^* N \cong f^*(\HIF(I,F) \tensor_S N)$, whence the target is
  $\tensorBC(f, F, Y) = F \tensor_S Y$). Then $\theta$ is the image of this element under the
  inverse of the corepresenting isomorphism \cref{thm:H_represents} of $\HIF(I \tensor_S N, F)$,
  evaluated at $Y$.
\end{definition}

\begin{definition}

\leanok
[The first-variable action of $H(-,F)$ on a morphism]
  \label{def:HMapFst}
  \lean{MR0555258CompactifyingPicard.HMapFst}
  \uses{def:H, def:homTensorFunctor, def:tensorBaseChangeFunctor, thm:H_represents}
  \textit{Project-local construction (the per-morphism substrate for \cref{lem:H_tensor}).}
  Fix $F$. For a morphism $\varphi : J \to J'$ of $\Ox$-modules and admissibility
  witnesses $\mathrm{Admissible}(f,J,F)$, $\mathrm{Admissible}(f,J',F)$, the morphism
  \[
    \mathtt{HMapFst}\,f\,F\,\varphi \;:\; \HIF(J,F) \longrightarrow \HIF(J',F)
  \]
  is defined as follows. Precomposition with $\varphi$ gives a natural transformation
  $\mathtt{homTensorFunctor}\,f\,J'\,F \to \mathtt{homTensorFunctor}\,f\,J\,F$ of the
  corepresented functors $M \mapsto \Hom_X(J^{(\prime)}, F \tensor_S M)$
  (\cref{def:homTensorFunctor}). Conjugating by the two corepresentation isomorphisms
  $\HIF(J^{(\prime)},F) \cong \mathtt{homTensorFunctor}\,f\,J^{(\prime)}\,F$
  (\cref{thm:H_represents}) and reflecting through the full-faithful
  $\mathtt{coyoneda}$ (the corepresenting object is unique up to unique isomorphism)
  produces the stated morphism; the direction is covariant in $J$ (precomposition is
  contravariant, and $\mathtt{coyoneda}$/$(-)^{\mathrm{op}}$ flips it again). Unlike
  \cref{def:HFunctorFst} this carries no totality hypothesis: only the two admissibility
  witnesses at $J$ and $J'$ are needed, so it is the genuinely reusable substrate the
  $(1.1.3)$ naturality (\cref{lem:htensorComparison_naturality}) is phrased against.
\end{definition}

\begin{lemma}[$H(-,F)$ preserves identities and composites in its first variable]
  \label{lem:HMapFst_functorial}
  \lean{MR0555258CompactifyingPicard.HMapFst_id, MR0555258CompactifyingPicard.HMapFst_comp}
  \uses{def:HMapFst}
  \textit{Project-local (functoriality laws for \cref{def:HMapFst};
  Lean: \texttt{HMapFst\_id} and \texttt{HMapFst\_comp}).}
  The action $\mathtt{HMapFst}$ of \cref{def:HMapFst} satisfies
  $\mathtt{HMapFst}\,f\,F\,(\id_J) = \id_{\HIF(J,F)}$ and
  $\mathtt{HMapFst}\,f\,F\,(\varphi \circ \psi)
    = \mathtt{HMapFst}\,f\,F\,\varphi \circ \mathtt{HMapFst}\,f\,F\,\psi$.
  \begin{proof}
    Both follow from the functoriality of precomposition, $\mathtt{coyoneda}$ and
    $(-)^{\mathrm{op}}$, together with the cancellation of the corepresentation
    isomorphisms: for the identity the middle precomposition transformation is the
    identity and the two isomorphisms cancel; for a composite the corepresentation
    isomorphisms at the middle object cancel and the remaining factors compose by
    $(\varphi \circ \psi)^{\mathrm{op}} = \psi^{\mathrm{op}} \circ \varphi^{\mathrm{op}}$.
  \end{proof}
\end{lemma}

\begin{definition}

\leanok
[$H(-,F)$ as a covariant functor in its first variable]
  \label{def:HFunctorFst}
  \lean{MR0555258CompactifyingPicard.HFunctorFst}
  \uses{def:H, def:homTensorFunctor, thm:H_represents, def:HMapFst, lem:HMapFst_functorial}
  \textit{Project-local construction (functorial substrate for \cref{lem:H_tensor}).}
  Fix $F$ (and the admissibility data). For a morphism $\varphi : J \to J'$ of
  $\Ox$-modules, precomposition gives a natural transformation
  $\Hom_X(J', F \tensor_S -) \to \Hom_X(J, F \tensor_S -)$ of the corepresented
  functors $\mathtt{homTensorFunctor}\,f\,J'\,F \to \mathtt{homTensorFunctor}\,f\,J\,F$
  (\cref{def:homTensorFunctor}); since $\HIF(J,F)$ corepresents
  $\mathtt{homTensorFunctor}\,f\,J\,F$ (\cref{thm:H_represents}) and the corepresenting
  object is unique up to a unique isomorphism, this transformation is induced by a
  unique morphism $\HIF(J,F) \to \HIF(J',F)$. These assemble into a covariant functor
  \[
    \mathtt{HFunctorFst}\,f\,F \;:\; \Xmod \longrightarrow \Smod, \qquad
    J \longmapsto \HIF(J,F),
  \]
  whose value on objects agrees with $\HIF(-,F)$. (This packages the
  first-variable functoriality of $H$ that \cref{lem:H_tensor} states informally as
  ``$H(I,F)$ is functorial in $I$''.)
\end{definition}

\begin{definition}

\leanok
[The domain functor $N \mapsto \HIF(I \tensor_S N, F)$]
  \label{def:htensorDomainFunctor}
  \lean{MR0555258CompactifyingPicard.htensorDomainFunctor}
  \uses{def:HFunctorFst, def:tensorBaseChangeFunctor, def:tensorBC}
  \textit{Project-local construction (functorial substrate for \cref{lem:H_tensor}).}
  Fix $I, F$ on $X$ and the proper map $f : X \to S$. The assignment
  $N \mapsto \HIF(I \tensor_S N, F) = \HIF(\tensorBC(f, I, N), F)$ is the composite
  functor
  \[
    \mathtt{htensorDomainFunctor}\,f\,I\,F
      \;=\; (\mathtt{tensorBaseChangeFunctor}\,f\,I) \;\circ\; (\mathtt{HFunctorFst}\,f\,F)
      \;:\; \Smod \longrightarrow \Smod,
  \]
  where $\mathtt{tensorBaseChangeFunctor}\,f\,I = f^*(-) \tensor_{\Ox}$-twist sends
  $N \mapsto I \tensor_{\Ox} f^* N = \tensorBC(f, I, N)$ (\cref{def:tensorBaseChangeFunctor},
  \cref{def:tensorBC}) and $\mathtt{HFunctorFst}\,f\,F$ is \cref{def:HFunctorFst}.
\end{definition}

\begin{lemma}

\leanok
[Naturality in $N$ of the comparison map $\theta$ (per-morphism)]
  \label{lem:htensorComparison_naturality}
  \lean{MR0555258CompactifyingPicard.htensorComparison_naturality}
  \uses{def:htensorComparison, def:HMapFst, def:tensorBaseChangeFunctor, def:tensorLeft, def:h, lem:tensorMod_assoc_uncond, lem:pullbackTensorComparison, lem:compose_three_squares}
  \textit{Project-local construction (the per-morphism naturality the $(1.1.3)$
  presentation argument needs).}
  The pointwise comparison map $\theta_N$ of \cref{def:htensorComparison} is natural in
  $N$, stated \emph{per-morphism} (not as a transformation of total endofunctors, which
  is impossible: the domain functor $N \mapsto \HIF(I \tensor_S N, F)$ of
  \cref{def:htensorDomainFunctor} requires admissibility for every object, an
  unsatisfiable hypothesis since it forces every $\Ox$-module to be locally finitely
  presented). For a morphism $\psi : N \to N'$ of $\Os$-modules, with admissibility
  witnesses for $I, F$ and for $I \tensor_S N$, $I \tensor_S N'$, the square
  \[
    \begin{array}{ccc}
      \HIF(I \tensor_S N, F) & \xrightarrow{\ \theta_N\ } & \HIF(I,F) \tensor_S N \\[2pt]
      {\scriptstyle \mathtt{HMapFst}\,((\mathtt{tensorBaseChangeFunctor}\,f\,I)(\psi))}
        \big\downarrow & & \big\downarrow {\scriptstyle (\mathtt{tensorLeft}\,\HIF(I,F))(\psi)} \\[2pt]
      \HIF(I \tensor_S N', F) & \xrightarrow{\ \theta_{N'}\ } & \HIF(I,F) \tensor_S N'
    \end{array}
  \]
  commutes; here the left vertical is the first-variable action \cref{def:HMapFst} of
  $H(-,F)$ on $(\mathtt{tensorBaseChangeFunctor}\,f\,I)(\psi) : I \tensor_S N \to
  I \tensor_S N'$, and the right vertical is $\HIF(I,F) \tensor_S \psi$.
  \begin{proof}
    \uses{thm:H_represents}
    Both verticals composed with the relevant $\theta$ are morphisms \emph{out of} the
    corepresenting object $\HIF(I \tensor_S N, F)$, so by \cref{thm:H_represents} the two
    composites are equal iff the corresponding universal elements in
    $\Hom_X(I \tensor_S N,\ F \tensor_S (\HIF(I,F) \tensor_S N'))$ agree. Applying the
    corepresentation isomorphism reduces the square to an identity of universal elements,
    which is then assembled from: the naturality of the universal element $h(I,F)$
    (\cref{def:h}); the naturality of the associator
    (\cref{lem:tensorMod_assoc_uncond}); and the second-slot naturality of the
    $f^*$-monoidality comparison (\cref{lem:pullbackTensorComparison}, now a
    single-variable \texttt{NatIso} in $N$). The first-variable action $\mathtt{HMapFst}$
    is expanded through its defining $\mathtt{coyoneda}$-preimage to expose the same
    precomposition transformation. (Per the prover hand-off the square does not close by
    $\mathtt{simp}$/$\mathtt{aesop}$; it is the element-level corepresentation reflection
    just described.)
  \end{proof}
\end{lemma}

\begin{lemma}

\leanok
[Compose three commuting squares]
  \label{lem:compose_three_squares}
  \lean{MR0555258CompactifyingPicard.compose_three_squares}
  \textit{Project-local plumbing (used in \cref{lem:htensorComparison_naturality}).}
  In any category, given four composable morphisms $A, B, C, D$ and a factorisation of the
  intermediate composites into three commuting squares
  $A \circ B = L_1 \circ M_1$, $M_1 \circ C = L_2 \circ M_2$, $M_2 \circ D = L_3 \circ Q$,
  one has $A \circ B \circ C \circ D = (L_1 \circ L_2 \circ L_3) \circ Q$.
  \begin{proof}
    Immediate by substituting the three square identities in turn and reassociating; no
    hypotheses beyond the category axioms.
  \end{proof}
\end{lemma}

\begin{theorem}[Base-change unit $f^*\Os \cong \Ox$ (standard)]
  \label{thm:ega_pullbackUnit}
  \lean{MR0555258CompactifyingPicard.External.pullbackUnit}
  \textit{Source: standard (pullback of the structure sheaf); used here as an external input.}
  For a morphism of schemes $f : X \to S$, the pullback of the structure sheaf $\Os$
  (the tensor unit of $\Os$-modules) is canonically isomorphic to the structure sheaf $\Ox$
  (the tensor unit of $\Ox$-modules):
  \[
    f^{*}\Os \;\cong\; \Ox .
  \]
  \begin{proof}
    Standard: $f^{*}$ is a (strong) monoidal functor, so it carries the tensor unit to the
    tensor unit. Mathlib registers no monoidal structure on
    $\mathtt{Scheme.Modules.pullback}$, so the unit isomorphism is recorded as an external
    input (parallel to \cref{lem:pullbackTensorComparison}, the $f^*$-monoidality comparison).
  \end{proof}
\end{theorem}

\begin{theorem}[Eilenberg--Watts comparison for $\theta$ (external input)]
  \label{thm:ew_htensorComparison}
  \lean{MR0555258CompactifyingPicard.External.eilenbergWatts_htensorComparison}
  \uses{def:htensorComparison, lem:htensorComparison_naturality}
  % SOURCE: [Altman--Kleiman], §1, (1.1.3), p. 55 (read from references/MR0555258-compactifying-picard.pdf)
  % SOURCE QUOTE: "Moreover it is clearly right exact in each variable. In particular, the
  % functor   N |--> H(I (x)_S N, F)   is covariant and right exact. So we have a canonical
  % isomorphism,   H(I,F) (x) N = H(I (x)_S N, F).   (1.1.3)"
  \textit{Source: Altman--Kleiman, §1, (1.1.3), p. 55 (the ``So we have'' step), which invokes
  the Eilenberg--Watts theorem (Watts, \emph{Intrinsic characterizations of some additive
  functors}, Proc. AMS 11 (1960)).}
  Let $f : X \to S$ with $I, F$ admissible. The two functors $N \mapsto \HIF(I \tensor_S N, F)$
  and $N \mapsto \HIF(I,F) \tensor_S N$ of a quasi-coherent $\Os$-module $N$ are covariant,
  additive, right exact and sum-preserving, and $\theta$ (\cref{def:htensorComparison}) is the
  canonical comparison transformation between them (natural by
  \cref{lem:htensorComparison_naturality}). If $\theta_{\Os}$ is an isomorphism then $\theta_N$
  is an isomorphism for every quasi-coherent $N$.
  \begin{proof}
    This is the Eilenberg--Watts propagation principle: a natural transformation between two
    right-exact, direct-sum-preserving functors on quasi-coherent $\Os$-modules that is an
    isomorphism on the unit $\Os$ is an isomorphism on every quasi-coherent object. Indeed,
    locally choose a free presentation $\Os^{(J)} \to \Os^{(K)} \to N \to 0$; $\theta$ is an
    isomorphism on each free $\Os^{(\bullet)}$ (a direct sum of $\theta_{\Os}$, using
    sum-preservation), hence on the cokernel $N$ by right exactness and the five lemma. Mathlib
    has neither the Eilenberg--Watts theorem nor the free-presentation machinery for
    $\Os$-modules, and the source functor is not total (admissibility constrains the objects), so
    the principle is recorded as an external input \emph{conditional on $\theta_{\Os}$ invertible};
    the project supplies that hypothesis directly from the unitors and
    \cref{thm:ega_pullbackUnit}. This is strictly weaker than \cref{lem:H_tensor} (which is
    unconditional) and is the general theorem Altman--Kleiman cite, not a restatement of it.
  \end{proof}
\end{theorem}

\begin{theorem}[$f^*$ right-unitality coherence (standard)]
  \label{thm:ega_pullback_rightUnitality}
  \lean{MR0555258CompactifyingPicard.External.pullback_rightUnitality}
  \uses{lem:pullbackTensorComparison, thm:ega_pullbackUnit}
  \textit{Source: standard monoidal-functor coherence; used here as an external input.}
  For a morphism of schemes $f : X \to S$ and an $\Os$-module $A$, the right-unitality square of
  the (strong) monoidal functor $f^*$ commutes. Writing $\mu_{A,\Os} : f^*A \tensor f^*\Os \to
  f^*(A \tensor_S \Os)$ for the $f^*$-laxity comparison (\cref{lem:pullbackTensorComparison}),
  $\varepsilon : f^*\Os \cong \Ox$ for the unit comparison (\cref{thm:ega_pullbackUnit}), and
  $\rho$ for the right unitor, one has (in diagrammatic, left-to-right order)
  \[
    \mu_{A,\Os} \,;\, f^*(\rho_A) \;=\; (\id_{f^*A} \tensor \varepsilon) \,;\, \rho_{f^*A}.
  \]
  \begin{proof}
    Standard: this is one of the two unit coherence axioms of a (strong/lax) monoidal functor.
    Mathlib registers no monoidal structure on $\mathtt{Scheme.Modules.pullback}$, so — exactly as
    for the laxity $\mu$ (\cref{lem:pullbackTensorComparison}) and the unit $\varepsilon$
    (\cref{thm:ega_pullbackUnit}) — the coherence relating these two otherwise-independent
    comparison data is recorded as an external input. It mentions neither $\HIF$ nor admissibility,
    so it is strictly weaker than \cref{lem:H_tensor} and in-family with the existing
    $f^*$-monoidality anchors.
  \end{proof}
\end{theorem}

\begin{lemma}
[Right-unitor triangle of the sheafified tensor (internal monoidal coherence) --- Mathlib gap anchor]
  \label{lem:internal_tensorMod_rightUnitality}
  \lean{MR0555258CompactifyingPicard.External.tensorMod_rightUnitality}
  \uses{def:tensorMod, lem:tensorMod_assoc_uncond, lem:tensorMod_unitor, lem:sheafifyTensorComparison_uncond}
  % NOTE: This is NOT from Altman--Kleiman. It is the right-unitor coherence (the
  % "triangle" identity alpha ; (id (x) rho) = rho) of the hand-built sheafified-tensor
  % monoidal structure on X.Modules. It is TRUE in any monoidal category, but is ABSENT
  % from Mathlib here because no `MonoidalCategory X.Modules` instance exists: the whole
  % layer (associator tensorAssoc, unitor tensorRightUnitor, ...) is built ON TOP of the
  % anchored comparison `External.sheafifyTensorComparison{,Left}` (lem:sheafifyTensorComparison_uncond)
  % as bare NatIsos WITHOUT monoidal-functor coherence data. Unfolding the triangle leaves
  % the LHS carrying the two OPAQUE comparison anchors (inside tensorAssoc) against a
  % comparison-free RHS; the comparisons cannot cancel without a coherence axiom about
  % them (verified empirically iter-039: coherence/monoidal_coherence/aesop all stall).
  % It is in-family with lem:sheafifyTensorComparison_uncond (the same "module
  % sheafification is monoidal" package) and STRICTLY WEAKER than lem:H_tensor (it
  % mentions neither H nor admissibility nor Eilenberg--Watts). Anchoring the single
  % triangle H_tensor actually needs is leaner and more faithful than building the full
  % MonoidalCategory instance, which would require anchoring sheafification AS A monoidal
  % functor (all its coherences), bottoming out at the same comparison mismatch.
  For $\Ox$-modules $M, N$, the right-unitor triangle of the sheafified tensor on
  $X.\mathrm{Modules}$ commutes: writing $\alpha$ for the associator
  (\cref{lem:tensorMod_assoc_uncond}) and $\rho$ for the right unitor
  (\cref{lem:tensorMod_unitor}), in diagrammatic order
  \[
    \alpha_{M,N,\Ox} \,;\, (M \tensor \rho_N) \;=\; \rho_{M \tensor N}
    \quad:\quad (M \tensor N) \tensor \Ox \longrightarrow M \tensor N.
  \]
  \begin{proof}
    Standard monoidal coherence (a consequence of the triangle/pentagon axioms), holding
    in any monoidal category. Mathlib registers no \texttt{MonoidalCategory} instance on
    \texttt{X.Modules} (the layer is built on the anchored comparison
    \cref{lem:sheafifyTensorComparison_uncond} as bare natural isomorphisms), so the
    coherence relating the otherwise-opaque comparison data is recorded as an external
    input, exactly as for \cref{lem:sheafifyTensorComparison_uncond} itself. Used as an
    external input.
  \end{proof}
\end{lemma}

\begin{definition}\leanok
[Base-change unit collapse $I \tensor_S \Os \cong I$]
  \label{def:tensorUnitBC}
  \lean{MR0555258CompactifyingPicard.tensorUnitBC}
  \uses{def:tensorBC, thm:ega_pullbackUnit, lem:tensorMod_unitor}
  \textit{Project-local construction (substrate for \cref{lem:htensorComparison_unit_isIso}).}
  For $f : X \to S$ and an $\Ox$-module $I$, the base-change tensor at the unit collapses to $I$:
  \[
    I \tensor_S \Os \;=\; I \tensor f^*\Os \;\cong\; I \tensor \Ox \;\cong\; I,
  \]
  the first isomorphism whiskering the unit comparison $f^*\Os \cong \Ox$ (\cref{thm:ega_pullbackUnit})
  into the left factor, the second the right unitor (\cref{lem:tensorMod_unitor}).
\end{definition}

\begin{lemma}\leanok
[Admissibility transports to $I \tensor_S \Os$]
  \label{lem:admissible_tensorUnit}
  \lean{MR0555258CompactifyingPicard.admissible_tensorUnit}
  \uses{def:admissible, def:tensorUnitBC}
  \textit{Project-local construction (substrate for \cref{lem:htensorComparison_unit_isIso}).}
  If $(f, I, F)$ is admissible (\cref{def:admissible}) then so is $(f, I \tensor_S \Os, F)$.
  \begin{proof}
    Every field is inherited from the admissibility of $I$ except the local finite presentation of
    the first argument, which transports from that of $I$ along the isomorphism
    $I \tensor_S \Os \cong I$ (\cref{def:tensorUnitBC}), finite presentation being closed under
    isomorphism.
  \end{proof}
\end{lemma}

\begin{lemma}\leanok
[$\theta$ is an isomorphism at the unit]
  \label{lem:htensorComparison_unit_isIso}
  \lean{MR0555258CompactifyingPicard.htensorComparison_unit_isIso}
  \uses{def:htensorComparison, thm:H_represents, def:tensorUnitBC, lem:admissible_tensorUnit, lem:tensorMod_unitor, lem:pullbackTensorComparison, thm:ega_pullbackUnit, thm:ega_pullback_rightUnitality, lem:internal_tensorMod_rightUnitality}
  \textit{Project-local construction (the $\theta_{\Os}$ unit-agreement discharging \cref{thm:ew_htensorComparison} in \cref{lem:H_tensor}).}
  For $f : X \to S$ with $I, F$ admissible, the comparison map $\theta$ of \cref{def:htensorComparison}
  is an isomorphism at the unit $N = \Os$: the map
  $\theta_{\Os} : \HIF(I \tensor_S \Os, F) \to \HIF(I,F) \tensor_S \Os$ is an isomorphism.
  \begin{proof}
    By corepresentability (\cref{thm:H_represents}) it suffices to show $\theta_{\Os}$ is carried by
    the coyoneda embedding to an isomorphism; since coyoneda reflects isomorphisms, this reduces to
    identifying $\mathrm{coyoneda}(\theta_{\Os})$ with the explicit comparison iso $\Psi$ built from
    the right unitor $\rho_{\HIF(I,F)}$ (\cref{lem:tensorMod_unitor}), the genuine corepresentation of
    $\mathrm{homTensorFunctor}(I,F)$, and the unit collapse $u : I \tensor_S \Os \cong I$
    (\cref{def:tensorUnitBC}). Evaluating both sides at the identity (the Yoneda lemma), the claim
    reduces to the single coherence identity asserting that the defining universal element $e_{\Os}$ of
    $\theta_{\Os}$ — namely $(h(I,F) \tensor \id) \,;\, \mathrm{assoc} \,;\, (F \tensor \mu_{H,\Os})$ —
    equals $h(I,F)$ transported along the unit collapse $u$ and the right unitor. This is precisely
    the right-unitality coherence of the monoidal functor $f^*$ relating the laxity $\mu$
    (\cref{lem:pullbackTensorComparison}) and the unit $\varepsilon$ (\cref{thm:ega_pullbackUnit}): it
    follows from \cref{thm:ega_pullback_rightUnitality} by the naturality of the right-unit collapse
    at $h(I,F)$ (a whisker-exchange square together with right-unitor naturality, mirroring the
    construction of \cref{lem:htensorComparison_naturality}), the associator naturality in the third
    slot at $\varepsilon$, and the internal right-unitor triangle of the sheafified tensor
    (\cref{lem:internal_tensorMod_rightUnitality}, which turns the $f^*$-level coherence into the
    associator form). The admissibility of $I \tensor_S \Os$ used throughout is
    \cref{lem:admissible_tensorUnit}.
  \end{proof}
\end{lemma}

\begin{lemma}\leanok
[Right exactness of $H$ in the second-variable tensor]
  \label{lem:H_tensor}
  \lean{MR0555258CompactifyingPicard.H_tensor}
  \uses{def:H, thm:H_represents, def:tensorMod, def:tensorBC, def:htensorComparison, lem:htensorComparison_naturality, lem:pullbackTensorComparison, lem:tensorMod_unitor, thm:ega_pullbackUnit, thm:ew_htensorComparison, lem:admissible_tensorUnit, lem:htensorComparison_unit_isIso}
  % SOURCE: [Altman--Kleiman], §1, (1.1.3), p. 55 (read from references/MR0555258-compactifying-picard.pdf)
  % SOURCE QUOTE: "The O_S-Module H(I,F) is obviously functorial in I and F; it is
  % covariant in I and contravariant in F. Moreover it is clearly right exact in each
  % variable. In particular, the functor   N |--> H(I (x)_S N, F)   is covariant and
  % right exact. So we have a canonical isomorphism,   H(I,F) (x) N = H(I (x)_S N, F).   (1.1.3)"
  \textit{Source: Altman--Kleiman, §1, (1.1.3), p. 55.}
  The $\Os$-module $\HIF(I,F)$ is functorial in $I$ and $F$ (covariant in $I$,
  contravariant in $F$) and right exact in each variable. In particular, for a
  quasi-coherent $\Os$-module $N$ there is a canonical isomorphism
  \[
    \HIF(I,F) \tensor N \;=\; \HIF(I \tensor_S N,\; F).
    \tag{1.1.3}
  \]
  \begin{proof}
    The claimed isomorphism is the inverse of the canonical comparison map
    $\theta : \HIF(I \tensor_S N, F) \to \HIF(I,F) \tensor_S N$ of
    \cref{def:htensorComparison}, which is natural in $N$. It remains to show $\theta$ is an
    isomorphism.

    \emph{Both sides are right-exact, sum-preserving functors of $N$.} The functor
    $N \mapsto \HIF(I,F) \tensor_S N$ is right exact and commutes with arbitrary direct sums,
    being a tensor product. The functor $N \mapsto \HIF(I \tensor_S N, F)$ is the composite of
    $N \mapsto I \tensor_S N$ (right exact, sum-preserving) with $J \mapsto \HIF(J, F)$; the
    latter is right exact because by \cref{thm:H_represents} it corepresents
    $M \mapsto \Hom_X(J, F \tensor_S M)$, a functor that is left exact and contravariant in $J$,
    so by Yoneda $\HIF(-, F)$ turns right-exact sequences into right-exact ones and preserves
    direct sums. Hence both functors of $N$ are covariant, additive, right exact and
    sum-preserving.

    \emph{Agreement on the unit.} For $N = \Os$ one has $I \tensor_S \Os \cong I$ and
    $\HIF(I,F) \tensor_S \Os \cong \HIF(I,F)$ (the unitors, \cref{lem:tensorMod_unitor}, together
    with the base-change unit $f^*\Os \cong \Ox$, \cref{thm:ega_pullbackUnit}), and under these
    identifications the defining element $h(I,F) \tensor \id$ of \cref{def:htensorComparison}
    reduces to $h(I,F)$ itself; thus $\theta_{\Os}$ corresponds under \cref{thm:H_represents} to the
    universal element $h(I,F)$, i.e. is the canonical isomorphism
    $\HIF(I \tensor_S \Os, F) \cong \HIF(I,F) \cong \HIF(I,F) \tensor_S \Os$. This is the only
    isomorphism the assembly proves directly; it is recorded as $\theta_{\Os}$ being invertible.

    \emph{$\theta$ is an isomorphism.} A natural transformation between two right-exact,
    direct-sum-preserving functors that is an isomorphism on $\Os$ is an isomorphism on every
    quasi-coherent $N$: locally on $S$ (affine) choose a presentation
    $\Os^{(J)} \to \Os^{(K)} \to N \to 0$ by direct sums of $\Os$; $\theta$ is an isomorphism on
    each $\Os^{(\bullet)}$ (a direct sum of the isomorphism on $\Os$), hence on the cokernel $N$
    by right exactness. This is the Eilenberg--Watts comparison of two right-exact, sum-preserving
    functors agreeing on the unit, supplied as the external input \cref{thm:ew_htensorComparison}:
    the per-morphism naturality \cref{lem:htensorComparison_naturality} witnesses that $\theta$ is
    the comparison transformation, and the unit-agreement $\theta_{\Os}$ above discharges its
    hypothesis, so $\theta_N$ is an isomorphism for every quasi-coherent $N$. Therefore $\theta$ is
    an isomorphism, giving \textnormal{(1.1.3)}.
  \end{proof}
\end{lemma}

\begin{definition}\leanok
[Restriction ring-sheaf datum]
  \label{def:restrictionRingSheafHom}
  \lean{AlgebraicGeometry.Scheme.Modules.restrictionRingSheafHom}
  \textit{Project-local Lean realization (supporting \cref{def:extensionByZero}).}
  For an open immersion $j : U \hookrightarrow X$ of schemes,
  $\mathtt{restrictionRingSheafHom}\,j$ is the morphism of sheaves of rings
  $\Ox|U \to (j_*)\,\mathcal{O}_U$ that Mathlib keeps \emph{inline} inside the
  definition of $\mathtt{restrictFunctor}\,j$, exposed here as a named datum
  $\varphi$ so that the identification $j^{*} = \mathtt{pushforward}\,\varphi$
  (\cref{lem:restrictFunctor_eq_pushforward}) can be stated and the
  extension-by-zero adjunction typed. Built from the inverse structure-sheaf
  isomorphisms $\mathtt{appIso}$ of $j$ on opens, whiskered through
  $\mathrm{forget}_2\,\mathsf{CommRing}\,\mathsf{Ring}$. Project-local because
  Mathlib does not name it.
\end{definition}

\begin{lemma}\leanok
[Restriction is pushforward of the ring-sheaf datum]
  \label{lem:restrictFunctor_eq_pushforward}
  \lean{AlgebraicGeometry.Scheme.Modules.restrictFunctor_eq_pushforward}
  \uses{def:restrictionRingSheafHom}
  \textit{Project-local Lean realization (supporting \cref{lem:extensionByZeroAdjunction}).}
  For an open immersion $j : U \hookrightarrow X$, the restriction functor
  $j^{*} = \mathtt{restrictFunctor}\,j$ on sheaves of modules is
  \emph{definitionally equal} to the sheaf-of-modules pushforward
  $\mathtt{pushforward}\,(\mathtt{restrictionRingSheafHom}\,j)$ along the
  continuous open-immersion site map. The proof is $\mathtt{rfl}$. This is the
  defeq that lets $j_!$ be defined as $\mathtt{pullback}$ of the same datum and
  the adjunction $j_! \dashv j^{*}$ be read off
  $\mathtt{pullbackPushforwardAdjunction}$.
\end{lemma}

\begin{definition}\leanok
[Extension by zero $j_!$]
  \label{def:extensionByZero}
  \lean{AlgebraicGeometry.Scheme.Modules.extensionByZero}
  For an open immersion $j : U \hookrightarrow X$ of schemes, the
  \emph{extension by zero} functor
  $j_! : \ModuleCat(U) \to \ModuleCat(X)$ is the left adjoint of the restriction
  functor $j^{*} = \mathtt{restrictFunctor}\,j$. Concretely, for an
  $\mathcal{O}_U$-module $M$, $j_! M$ is the sheafification of the presheaf of
  modules on $X$ given by
  \[
    V \longmapsto
      \begin{cases} M(V) & V \subseteq U,\\[2pt] 0 & V \not\subseteq U,\end{cases}
  \]
  whose stalks are $(j_! M)_x = M_x$ for $x \in U$ and $(j_! M)_x = 0$ for
  $x \notin U$. This is the Lean realization of the operation $(j_U)_!$ used in
  \cref{lem:acyclic_surjection}. It is \emph{absent} from Mathlib v4.30.0, which
  provides $j^{*} = \mathtt{restrictFunctor}\,j$ and its \emph{right} adjoint
  $j_{*} = \mathtt{pushforward}\,j$ (with $\mathtt{restrictAdjunction}$), but no
  left adjoint to $j^{*}$.
\end{definition}

\begin{lemma}\leanok
[Adjunction $j_! \dashv j^{*}$]
  \label{lem:extensionByZeroAdjunction}
  \lean{AlgebraicGeometry.Scheme.Modules.extensionByZeroAdjunction}
  \uses{def:extensionByZero}
  For an open immersion $j : U \hookrightarrow X$ there is an adjunction
  $j_! \dashv j^{*}$, i.e. a natural bijection
  $\Hom_X(j_! M, F) \cong \Hom_U(M, j^{*} F)$ in $M \in \ModuleCat(U)$ and
  $F \in \ModuleCat(X)$.
  \begin{proof}
    \uses{def:extensionByZero, def:restrictionRingSheafHom, lem:restrictFunctor_eq_pushforward}
    The adjunction is read off Mathlib's site-pushforward machinery through a
    definitional identification, avoiding any hand-rolled Kan extension or
    sheafification. Mathlib's restriction functor $j^{*} =
    \mathtt{restrictFunctor}\,j$ is, \emph{definitionally}, the sheaf-of-modules
    pushforward $\mathtt{pushforward}\,\varphi$ along the canonical ring-sheaf
    morphism $\varphi = \mathtt{restrictionRingSheafHom}\,j$ that $j$ induces (this
    is \cref{lem:restrictFunctor_eq_pushforward}, which holds by $\mathtt{rfl}$).
    Because the base sites are the \emph{small} categories $\mathrm{Opens}(U)$,
    $\mathrm{Opens}(X)$, Mathlib's $\mathtt{PullbackContinuous}$ already equips
    $\mathtt{pushforward}\,\varphi$ with a left adjoint, namely
    $\mathtt{pullback}\,\varphi$, together with the adjunction
    $\mathtt{pullbackPushforwardAdjunction}\,\varphi :
    \mathtt{pullback}\,\varphi \dashv \mathtt{pushforward}\,\varphi$. Setting
    $j_! := \mathtt{pullback}\,\varphi$ (this is \cref{def:extensionByZero}), the
    adjunction $j_! \dashv j^{*}$ is exactly
    $\mathtt{pullbackPushforwardAdjunction}\,\varphi$ re-read through the defeq
    $j^{*} = \mathtt{pushforward}\,\varphi$. In particular $j_!$, as a left
    adjoint, preserves all colimits — used in \cref{lem:acyclic_surjection} to
    commute $j_!$ past direct sums.
  \end{proof}
\end{lemma}

\begin{lemma}

\leanok
[Restriction $j^{*}$ is exact]
  \label{lem:restrictFunctor_exact}
  \lean{AlgebraicGeometry.Scheme.Modules.restrictFunctor_preservesLimits}
  \uses{lem:extensionByZeroAdjunction}
  \textit{Project-local Lean realization (the $j^{*}$ half of $j_!/j^{*}$ exactness).}
  For an open immersion $j : U \hookrightarrow X$ the restriction functor
  $j^{*} = \mathtt{restrictFunctor}\,j$ preserves all limits.
  \begin{proof}
    \uses{lem:extensionByZeroAdjunction}
    $j^{*}$ is the \emph{right} adjoint of the adjunction $j_! \dashv j^{*}$
    (\cref{lem:extensionByZeroAdjunction}), and right adjoints preserve limits.
    (Since $j^{*}$ is moreover a left adjoint, $\mathtt{restrictAdjunction} :
    j^{*} \dashv j_{*}$, it preserves colimits too, hence is exact.)
  \end{proof}
\end{lemma}

\begin{lemma}

\leanok
[$j_!$ exactness reduces to mono-preservation]
  \label{lem:extensionByZero_exact_of_mono}
  \lean{AlgebraicGeometry.Scheme.Modules.extensionByZero_preservesFiniteLimits_of_preservesMono}
  \uses{def:extensionByZero, lem:extensionByZeroAdjunction, lem:restrictFunctor_exact}
  \textit{Project-local Lean realization (reduction step for \cref{lem:extensionByZero_exact}).}
  For an open immersion $j$, \emph{if} $j_!$ preserves monomorphisms then $j_!$
  preserves finite limits.
  \begin{proof}
    \uses{lem:extensionByZeroAdjunction, lem:restrictFunctor_exact}
    $j_!$ is a left adjoint (\cref{lem:extensionByZeroAdjunction}), hence
    right exact (preserves finite colimits); it is additive ($\mathtt{left\_adjoint\_additive}$),
    as is $j^{*}$ (additive because it preserves finite products, by
    \cref{lem:restrictFunctor_exact}). For an additive right-exact functor of
    abelian categories, preservation of finite limits is equivalent to
    preservation of monomorphisms
    ($\mathtt{preservesFiniteLimits\_iff\_forall\_exact\_map\_and\_mono}$).
  \end{proof}
\end{lemma}

\begin{lemma}

\leanok
[$j_!$ preserves monomorphisms (presheaf-pullback reduction)]
  \label{lem:extensionByZero_preservesMono}
  \lean{AlgebraicGeometry.Scheme.Modules.extensionByZero_preservesMonomorphisms_of_presheafPullback}
  \uses{def:extensionByZero, lem:restrictFunctor_exact}
  \textit{Project-local Lean realization (conditional reduction of $j_!$
  mono-preservation; Mathlib-gap infra).}
  For an open immersion $j : U \hookrightarrow X$, \emph{if} the presheaf-level
  pullback $\mathtt{PresheafOfModules.pullback}\,\varphi$ along
  $\varphi = \mathtt{restrictionRingSheafHom}\,j$ preserves monomorphisms, then so
  does $j_!$.
  \begin{proof}
    Factor $j_! = \mathtt{pullback}\,\varphi$ through
    $\mathtt{SheafOfModules.pullbackIso}$ as
    $\mathtt{forget} \circ \mathtt{PresheafOfModules.pullback}\,\varphi \circ
    \mathtt{sheafification}$. The forgetful functor is a right adjoint and
    sheafification is left exact, so both preserve monomorphisms unconditionally;
    hence $j_!$ preserves monomorphisms as soon as the middle presheaf-pullback
    factor does.
  \end{proof}
\end{lemma}

\begin{lemma}

\leanok
[Extension by zero is exact (presheaf-pullback reduction)]
  \label{lem:extensionByZero_exact}
  \lean{AlgebraicGeometry.Scheme.Modules.extensionByZero_preservesFiniteLimits_of_presheafPullback}
  \uses{def:extensionByZero, lem:extensionByZero_exact_of_mono, lem:extensionByZero_preservesMono, lem:restrictFunctor_exact}
  \textit{Project-local Lean realization (conditional exactness of $j_!$;
  Mathlib-gap infra).}
  For an open immersion $j : U \hookrightarrow X$, \emph{if} the presheaf-level
  pullback along $\varphi = \mathtt{restrictionRingSheafHom}\,j$ preserves
  monomorphisms, then $j_!$ preserves finite limits; together with the colimit
  preservation it has as a left adjoint (\cref{lem:extensionByZeroAdjunction}),
  $j_!$ is then \emph{exact}. The restriction $j^{*} = \mathtt{restrictFunctor}\,j$
  is unconditionally exact (\cref{lem:restrictFunctor_exact}).
  \begin{proof}
    \uses{lem:extensionByZero_exact_of_mono, lem:extensionByZero_preservesMono}
    By \cref{lem:extensionByZero_exact_of_mono} it suffices that $j_!$ preserves
    monomorphisms, which is \cref{lem:extensionByZero_preservesMono} under the
    same presheaf-pullback hypothesis.
  \end{proof}
\end{lemma}

\begin{theorem}[Existence of an extension-by-zero surjection (standard)]
  \label{thm:exists_extensionByZero_surjection}
  \lean{MR0555258CompactifyingPicard.External.exists_extensionByZero_surjection}
  \uses{def:extensionByZero, def:extensionByZeroCoprod}
  % SOURCE: [Altman--Kleiman], §1, (1.2) proof, p. 55 (read from references/MR0555258-compactifying-picard.pdf)
  % SOURCE QUOTE: "For any element f in any stalk of I, there is an affine
  % neighborhood U of the stalk and an element g in Gamma(U, I) whose image in the
  % stalk is equal to f. So there are a family of affine open sets U and a
  % surjection J = ||_U J_U -> I, where J_U denotes the extension by zero
  % (j_U)_!(O_X | U), where j_U denotes the inclusion of U in X."
  \textit{Source: as invoked in Altman--Kleiman, §1, in the proof of (1.2).}
  Let $X$ be a scheme and $I$ an arbitrary $\Ox$-module. Then there exist a family
  $(U_i)_{i \in \iota}$ of affine open subsets of $X$ (with inclusions
  $j_{U_i} : U_i \hookrightarrow X$) and an epimorphism
  \[
    \pi : \;\coprod_{i \in \iota} (j_{U_i})_!\bigl(\mathcal{O}_{U_i}\bigr)
      \;\longrightarrow\; I .
  \]
  \begin{proof}
    For every point and every germ of $I$ there is an affine open $U$ over which
    the germ is the image of a section $s \in \Gamma(U, I)$; via the adjunction
    $j_! \dashv j^{*}$ (\cref{lem:extensionByZeroAdjunction}) the section $s$
    corresponds to a map $(j_U)_!(\mathcal{O}_U) \to I$. Indexing over all such
    pairs $(U, s)$ and taking the coproduct of these maps yields a morphism
    $\pi$ whose image contains every germ of $I$, hence an epimorphism. This is
    the standard ``sections over affine opens generate'' fact, used by
    Altman--Kleiman without further comment; it is anchored here as a Mathlib-gap
    input (the stalkwise surjectivity is not available for
    $\mathtt{SheafOfModules}$ in Mathlib v4.30.0), to build/upstream later.
  \end{proof}
\end{theorem}

\begin{theorem}[Acyclicity of the pullback of an extension-by-zero sum (standard)]
  \label{thm:extensionByZero_coprod_acyclic}
  \lean{MR0555258CompactifyingPicard.External.extensionByZero_coprod_acyclic}
  \uses{def:extensionByZero, def:extensionByZeroCoprod, def:extGroup, thm:ega_ext_affine_acyclic}
  % SOURCE: [Altman--Kleiman], §1, (1.2) proof, p. 55 (read from references/MR0555258-compactifying-picard.pdf)
  % SOURCE QUOTE: "Then g*J is equal to ||_U J_{g^{-1}U} because pullback commutes
  % with direct sum and with extension by zero. Hence we have Hom_Y(g*J, F) = prod
  % Hom_Y(J_{g^{-1}U}, F) = prod Hom_{g^{-1}U}(O_{g^{-1}U}, F | g^{-1}(U)) = prod
  % Gamma(g^{-1}U, F). Therefore we have Ext^q_Y(g*J, F) = prod H^q(g^{-1}U, F |
  % g^{-1}U). Since g^{-1}U is affine and F is quasi-coherent, the right-hand side
  % is equal to zero for q > 0."
  \textit{Source: as invoked in Altman--Kleiman, §1, in the proof of (1.2).}
  Let $(U_i)_{i \in \iota}$ be a family of affine open subsets of a scheme $X$,
  let $g : Y \to X$ be an affine morphism, let $F$ be a quasi-coherent
  $\Oy$-module, and let $q > 0$. Writing
  $J = \coprod_i (j_{U_i})_!(\mathcal{O}_{U_i})$, the pullback $g^{*}J$ is acyclic
  for $\Hom_Y(-,F)$:
  \[
    \Ext^q_Y\bigl(g^{*}J,\; F\bigr) \;=\; 0 .
  \]
  \begin{proof}
    Pullback commutes with coproducts and with extension by zero, so
    $g^{*}J \cong \coprod_i (j_{g^{-1}U_i})_!(\mathcal{O}_{g^{-1}U_i})$ with each
    $g^{-1}U_i$ affine. Since $\Ext^q$ turns the coproduct into a product and the
    adjunction $j_! \dashv j^{*}$ passes to derived categories,
    $\Ext^q_Y(g^{*}J, F) = \prod_i \Ext^q_{g^{-1}U_i}(\mathcal{O}_{g^{-1}U_i},
    F|g^{-1}U_i) = \prod_i H^q(g^{-1}U_i, F|g^{-1}U_i)$, which vanishes for $q>0$
    because each $g^{-1}U_i$ is affine and $F$ quasi-coherent
    (\cref{thm:ega_ext_affine_acyclic}). This bundles the standard base-change,
    derived-adjunction and additivity facts Altman--Kleiman invoke without proof;
    it is anchored as a Mathlib-gap input. The underlying $j_!$-exactness route
    (\cref{lem:extensionByZero_exact}) reduces to a pointwise Lan formula for
    $\mathtt{PresheafOfModules.pullback}$ absent from Mathlib v4.30.0, and the
    derived adjunction needs derivability data ($\mathtt{Adjunction.derived}$)
    likewise absent for these sheafification localizations — both to
    build/upstream later.
  \end{proof}
\end{theorem}

\begin{definition}\leanok
[The extension-by-zero coproduct $J = \coprod_U J_U$]
  \label{def:extensionByZeroCoprod}
  \lean{MR0555258CompactifyingPicard.extensionByZeroCoprod}
  \uses{def:extensionByZero}
  For a scheme $X$, an index type $\iota$ and a family $(U_i)_{i\in\iota}$ of open
  subsets, the $\Ox$-module
  $\coprod_{i} (j_{U_i})_!\bigl(\mathcal{O}_{U_i}\bigr)$, where
  $j_{U_i}$ is the inclusion of $U_i$ and $(j_{U_i})_!$ the extension by zero
  (\cref{def:extensionByZero}). This is the Lean realization of the coproduct
  $J = \coprod_U J_U$ that appears in the surjection of \cref{lem:acyclic_surjection};
  it is shared (as a private abbreviation) by \cref{thm:exists_extensionByZero_surjection}
  and \cref{thm:extensionByZero_coprod_acyclic} so that the witness and the acyclicity
  conclusion match definitionally.
\end{definition}

\begin{lemma}\leanok
[Existence of an acyclic surjection]
  \label{lem:acyclic_surjection}
  \lean{MR0555258CompactifyingPicard.exists_acyclic_surjection}
  \uses{thm:ega_ext_affine_acyclic, def:extensionByZero, thm:exists_extensionByZero_surjection, thm:extensionByZero_coprod_acyclic}
  % SOURCE: [Altman--Kleiman], §1, (1.2), p. 55 (read from references/MR0555258-compactifying-picard.pdf)
  % SOURCE QUOTE: "(1.2) LEMMA. Let X be a scheme and let I be an arbitrary
  % O_X-Module. Then there exists a surjection J -> I in which J is an
  % O_X-Module such that for each affine morphism g : Y -> X and each
  % quasi-coherent O_Y-Module F, the pullback g*J is acyclic for the functor
  % Hom_Y(-, F)."
  \textit{Source: Altman--Kleiman, §1, (1.2), p. 55.}
  Let $X$ be a scheme and $I$ an arbitrary $\Ox$-module. Then there is a
  surjection $J \to I$ with $J$ an $\Ox$-module such that for every affine
  morphism $g : Y \to X$ and every quasi-coherent $\Oy$-module $F$, the pullback
  $g^* J$ is acyclic for $\Hom_Y(-,F)$; that is, $\Ext^q_Y(g^* J, F) = 0$ for
  $q > 0$.
  % SOURCE QUOTE PROOF: "Proof. For any element f in any stalk of I, there is an
  % affine neighborhood U of the stalk and an element g in Gamma(U, I) whose
  % image in the stalk is equal to f. So there are a family of affine open sets
  % U and a surjection J = ||_U J_U -> I, where J_U denotes the extension by
  % zero (j_U)_!(O_X | U), where j_U denotes the inclusion of U in X. Then g*J
  % is equal to ||_U J_{g^{-1}U} because pullback commutes with direct sum and
  % with extension by zero. Hence we have Hom_Y(g*J, F) = prod Hom_Y(J_{g^{-1}U},
  % F) = prod Hom_{g^{-1}U}(O_{g^{-1}U}, F | g^{-1}(U)) = prod Gamma(g^{-1}U, F).
  % Therefore we have Ext^q_Y(g*J, F) = prod H^q(g^{-1}U, F | g^{-1}U). Since
  % g^{-1}U is affine and F is quasi-coherent, the right-hand side is equal to
  % zero for q > 0."
  \begin{proof}
    \uses{thm:ega_ext_affine_acyclic, def:extensionByZero, thm:exists_extensionByZero_surjection, thm:extensionByZero_coprod_acyclic}
    For each section of $I$ choose an affine open $U \subseteq X$ over which it
    extends; the inclusion $j_U : U \hookrightarrow X$ gives the extension by zero
    $J_U = (j_U)_!(\Ox|U)$. Taking $J = \coprod_U J_U$ over this family of affine
    opens yields a surjection $J \to I$. For an affine morphism $g : Y \to X$,
    pullback commutes with direct sums and with extension by zero, so
    $g^* J = \coprod_U J_{g^{-1}U}$ with $g^{-1}U$ affine. Hence
    \[
      \Hom_Y(g^* J, F) = \prod_U \Hom_{g^{-1}U}(\mathcal{O}_{g^{-1}U},
        F|g^{-1}U) = \prod_U \Gamma(g^{-1}U, F),
    \]
    and therefore $\Ext^q_Y(g^* J, F) = \prod_U H^q(g^{-1}U, F|g^{-1}U)$. Since
    each $g^{-1}U$ is affine and $F$ is quasi-coherent, this vanishes for $q>0$ by
    \cref{thm:ega_ext_affine_acyclic}.
  \end{proof}
\end{lemma}

\begin{lemma}\leanok
[Half-exactness of $N \mapsto \Ext^q_X(I, F\tensor_S N)$]
  \label{lem:ext_tensor_halfexact}
  \lean{MR0555258CompactifyingPicard.ext_tensor_halfexact}
  \uses{def:extGroup, def:tensorBaseChangeFunctor, def:isSFlat, thm:flat_tensor_exact, lem:tensorBaseChangeFunctor_preservesZeroMorphisms}
  % NOTE: Lean decl realized @iter-023 in CONDITIONAL form — the S-flatness premise
  % `IsSFlat f F` is NOT consumed directly; instead the SES-preservation it implies is
  % taken as an explicit hypothesis `hSE : (sc.map (tensorBaseChangeFunctor f F)).ShortExact`,
  % and the conclusion is the exactness-at-the-middle (membership) form matching Mathlib's
  % `CategoryTheory.Abelian.Ext.covariant_sequence_exact₂`. To make it UNCONDITIONAL the
  % planner must supply the deferred anchor `External.flat_tensor_exact` (S-flat ⟹ hSE),
  % a from-scratch stalkwise-flatness gap parallel to `tilde_exact`. \leanok reflects the
  % conditional decl compiling sorry-free, not the full S-flat statement in the prose.
  % SOURCE: [Altman--Kleiman], §1, (1.3) proof, p. 56 (read from references/MR0555258-compactifying-picard.pdf)
  % SOURCE QUOTE: "Since T is half-exact and since T(k(s)) is equal to zero, T(M)
  % is equal to zero for every finitely generated A-module M [OB, 2.1] or
  % [EGA III_2, 7.5.3])."
  \textit{Project-local homological core of \cref{thm:H_locallyFree}: the
  half-exactness of the functor $T$ that Altman--Kleiman invoke in the (1.3)
  proof.}
  Let $f : X \to S$, let $I, F$ be $\Ox$-modules with $F$ $S$-flat
  (\cref{def:isSFlat}), and let $q \in \mathbb{N}$. Then the functor
  \[
    N \longmapsto \Ext^q_X\bigl(I,\; F \tensor_S N\bigr)
  \]
  from $\Os$-modules to abelian groups is half-exact: for every short exact
  sequence $0 \to N' \to N \to N'' \to 0$ of $\Os$-modules the induced sequence
  \[
    \Ext^q_X(I, F\tensor_S N') \to \Ext^q_X(I, F\tensor_S N) \to
      \Ext^q_X(I, F\tensor_S N'')
  \]
  is exact (at the middle term).
  \begin{proof}
    \uses{def:extGroup, def:tensorBaseChangeFunctor, def:isSFlat}
    Because $F$ is $S$-flat, the base-change tensor $F \tensor_S (-) =
    (\mathtt{tensorBaseChangeFunctor}\,f\,F)$ carries the short exact sequence
    $0 \to N' \to N \to N'' \to 0$ to a short exact sequence
    $0 \to F\tensor_S N' \to F\tensor_S N \to F\tensor_S N'' \to 0$ of
    $\Ox$-modules (flatness is checked stalkwise, where it is the exactness of
    $M \mapsto F_x \tensor_{\mathcal{O}_{S,f(x)}} M$). Applying the covariant long
    exact sequence of $\Ext^{\bullet}_X(I, -)$ to this short exact sequence
    (Mathlib's \texttt{CategoryTheory.Abelian.Ext.covariant\_sequence\_exact}
    family) yields exactness at the middle term, which is precisely the asserted
    half-exactness.
  \end{proof}
\end{lemma}

\begin{lemma}

\leanok
[Unconditional half-exactness for $F$ $S$-flat]
  \label{lem:ext_tensor_halfexact_of_flat}
  \lean{MR0555258CompactifyingPicard.ext_tensor_halfexact_of_flat}
  \uses{lem:ext_tensor_halfexact, thm:flat_tensor_exact, def:isSFlat, def:tensorBaseChangeFunctor}
  % SOURCE: [Altman--Kleiman], §1, (1.3) proof, p. 56 (read from references/MR0555258-compactifying-picard.pdf)
  % SOURCE QUOTE: "Since T is half-exact and since T(k(s)) is equal to zero, T(M)
  % is equal to zero for every finitely generated A-module M [OB, 2.1] or
  % [EGA III_2, 7.5.3])."
  \textit{Project-local corollary discharging the hypothesis $hSE$ of
  \cref{lem:ext_tensor_halfexact}.}
  Let $f : X \to S$, let $I, F$ be $\Ox$-modules with $F$ $S$-flat
  (\cref{def:isSFlat}), and let $q \in \mathbb{N}$. Then $N \mapsto
  \Ext^q_X(I, F\tensor_S N)$ is half-exact \emph{unconditionally}: for every short
  exact sequence of $\Os$-modules the induced $\Ext$-sequence is exact at the middle
  term.
  \begin{proof}
    \uses{lem:ext_tensor_halfexact, thm:flat_tensor_exact}
    Feed the flatness anchor \cref{thm:flat_tensor_exact} (which, from $F$ $S$-flat,
    produces the short exactness of $\mathtt{sc.map\,(tensorBaseChangeFunctor}\,f\,F)$)
    as the hypothesis $hSE$ of \cref{lem:ext_tensor_halfexact}.
  \end{proof}
\end{lemma}

\begin{lemma}

\leanok
[Module-valued lift of an additive functor out of $\mathrm{Mod}_A$]
  \label{lem:additive_functor_module_lift}
  \lean{MR0555258CompactifyingPicard.liftAdditiveToModuleCat}
  \uses{lem:end_ring_hom, lem:module_of_lift}
  \textit{Project-bespoke category-theory supplement (no external source): the
  plumbing that turns the $\Ext$-valued functor of (1.3) into an
  $A$-module-valued one, so it can be fed to \cref{thm:ob_halfexact_free}.}
  Let $A$ be a commutative ring and let $G : \mathrm{Mod}_A \to \mathrm{Ab}$ be an
  \emph{additive} functor. Then $G$ lifts to a functor $\widehat G :
  \mathrm{Mod}_A \to \mathrm{Mod}_A$ such that the underlying abelian group of
  $\widehat G(M)$ is $G(M)$, and the underlying group homomorphism of $\widehat
  G(\varphi)$ is $G(\varphi)$.
  \begin{proof}
    For each $A$-module $M$ the scalar map $a \mapsto a\cdot\mathrm{id}_M$ is a ring
    homomorphism $A \to \mathrm{End}_{\mathrm{Mod}_A}(M)$ (commutativity of $A$
    makes every $a\cdot\mathrm{id}_M$ an $A$-linear endomorphism, and the
    $\mathrm{Mod}_A$-category is $A$-linear). Because $G$ is additive it induces a
    ring homomorphism $\mathrm{End}(M) \to \mathrm{End}(G(M))$ (it preserves
    addition by additivity and composition/identity by functoriality). Composing,
    $A \to \mathrm{End}(G(M))$ is a ring homomorphism from a commutative ring,
    equipping $G(M)$ with an $A$-module structure. For a morphism $\varphi : M \to
    N$ the naturality square for $a\cdot\mathrm{id}$ shows $G(\varphi)$ is
    $A$-linear, so $M \mapsto G(M)$, $\varphi \mapsto G(\varphi)$ is a functor
    $\mathrm{Mod}_A \to \mathrm{Mod}_A$ lifting $G$.
  \end{proof}
\end{lemma}

\begin{definition}

\leanok
[The functor $T(M) = \Ext^1_X(I, F\tensor_S\widetilde M)$]
  \label{def:functorT}
  \lean{MR0555258CompactifyingPicard.functorT}
  \uses{lem:additive_functor_module_lift, lem:covariant_ext_functor, lem:tensor_base_change_additive, def:tensorBaseChangeFunctor, def:tildeMod, def:extGroup}
  % SOURCE: [Altman--Kleiman], §1, (1.3) proof, p. 55--56 (read from references/MR0555258-compactifying-picard.pdf)
  % SOURCE QUOTE: "Consider the functor, T(M) = Ext^1_X(I, F (x)_S M~). from the
  % category of finitely generated A-modules M to itself."
  \textit{Source: Altman--Kleiman, §1, (1.3) proof.}
  For $A$ Noetherian local, $f : X \to \Spec A$, and $\Ox$-modules $I$, $F$, the
  functor $T : \mathrm{Mod}_A \to \mathrm{Mod}_A$ is the
  \cref{lem:additive_functor_module_lift}-lift of the additive composite
  \[
    \mathrm{Mod}_A \xrightarrow{\ \widetilde{(-)}\ } (\Spec A)\text{-Mod}
      \xrightarrow{\ F\tensor_S(-)\ } X\text{-Mod}
      \xrightarrow{\ \Ext^1_X(I,-)\ } \mathrm{Ab},
  \]
  i.e.\ $T(M) = \Ext^1_X\bigl(I, F\tensor_S\widetilde M\bigr)$ with the $A$-module
  structure read off from functoriality in $M$.
\end{definition}

\begin{lemma}

\leanok
[$T$ is additive]
  \label{lem:functorT_additive}
  \lean{MR0555258CompactifyingPicard.functorT_additive}
  \uses{def:functorT, lem:additive_functor_module_lift, lem:lift_additive, lem:tensor_base_change_additive, lem:covariant_ext_functor_additive, lem:tensorLeft_additive}
  \textit{Project-bespoke: additivity of the functor $T$ of \cref{def:functorT}.}
  The functor $T$ of \cref{def:functorT} is additive.
  \begin{proof}
    \uses{def:functorT}
    $T$ is the module-valued lift (\cref{lem:additive_functor_module_lift}) of a
    composite of additive functors: $\widetilde{(-)}$ is additive (Mathlib's
    $\mathtt{tilde}$ is a left adjoint), $F\tensor_S(-) =
    \mathtt{tensorBaseChangeFunctor}\,f\,F$ is additive (\cref{lem:tensorLeft_additive}
    through the base-change pullback), and the covariant $\Ext^1_X(I,-)$ is additive.
    The lift forgets to the additive composite, and the forgetful functor
    $\mathrm{Mod}_A \to \mathrm{Ab}$ is faithfully additive, so $T$ is additive.
  \end{proof}
\end{lemma}

\begin{lemma}\leanok
[Ring homomorphism $A \to \mathrm{End}(G(M))$ from an additive functor]
  \label{lem:end_ring_hom}
  \lean{MR0555258CompactifyingPicard.endRingHom}
  \textit{Project-bespoke (multiplicative core of \cref{lem:additive_functor_module_lift}).}
  For an additive functor $G : \mathrm{Mod}_A \to \mathrm{Ab}$ and an object $M$,
  the assignment $a \mapsto G(a\cdot\mathrm{id}_M)$ is a ring homomorphism
  $A \to \mathrm{End}(G(M))$.
  \begin{proof}
    Additivity and the unit come from $G(\mathrm{id}) = \mathrm{id}$ and
    $G$ preserving addition; multiplicativity uses $G$ preserving composition
    together with the commutativity of $A$, via
    $(ab)\cdot\mathrm{id}_M = (b\cdot\mathrm{id}_M)\circ(a\cdot\mathrm{id}_M)$.
  \end{proof}
\end{lemma}

\begin{lemma}\leanok
[$A$-module structure induced by the End-action]
  \label{lem:module_of_lift}
  \lean{MR0555258CompactifyingPicard.moduleOfLift}
  \uses{lem:end_ring_hom}
  \textit{Project-bespoke (object half of \cref{lem:additive_functor_module_lift}).}
  Composing the ring homomorphism \cref{lem:end_ring_hom} with the canonical
  $\mathrm{End}(G(M))$-action equips $G(M)$ with an $A$-module structure.
  \begin{proof}
    Restriction of scalars along the ring homomorphism of \cref{lem:end_ring_hom}.
  \end{proof}
\end{lemma}

\begin{lemma}\leanok
[The lift is additive]
  \label{lem:lift_additive}
  \lean{MR0555258CompactifyingPicard.liftAdditiveToModuleCat_additive}
  \uses{lem:additive_functor_module_lift}
  \textit{Project-bespoke.}
  The lift $\widehat G$ of \cref{lem:additive_functor_module_lift} is additive: on
  morphisms it acts as $G$, which is additive.
  \begin{proof}
    Its action on morphisms is the underlying group homomorphism of $G(\varphi)$,
    additive in $\varphi$ because $G$ is additive.
  \end{proof}
\end{lemma}

\begin{lemma}\leanok
[Covariant $\Ext$ as a functor]
  \label{lem:covariant_ext_functor}
  \lean{MR0555258CompactifyingPicard.covariantExtFunctor}
  \uses{def:extGroup}
  \textit{Project-bespoke: the functorial form of \cref{def:extGroup} in its second
  variable, absent from Mathlib v4.30.0.}
  For $I$ an $\Ox$-module on $Y$ and $q \in \mathbb{N}$, the assignment
  $N \mapsto \Ext^q_Y(I, N)$ is a functor $Y\text{-Mod} \to \mathrm{Ab}$; on a
  morphism $f : N \to N'$ it is $x \mapsto x \circ \mathrm{mk}_0(f)$. (As with
  \cref{def:extGroup} the groups live one universe up.)
  \begin{proof}
    Functoriality is identity ($x\circ\mathrm{mk}_0(\mathrm{id}) = x$) and
    composition ($\mathrm{mk}_0(f\circ g) = \mathrm{mk}_0(f)\circ\mathrm{mk}_0(g)$
    with associativity of $\Ext$-composition).
  \end{proof}
\end{lemma}

\begin{lemma}\leanok
[Covariant $\Ext$ functor is additive]
  \label{lem:covariant_ext_functor_additive}
  \lean{MR0555258CompactifyingPicard.covariantExtFunctor_additive}
  \uses{lem:covariant_ext_functor}
  \textit{Project-bespoke.}
  The functor \cref{lem:covariant_ext_functor} is additive.
  \begin{proof}
    $\mathrm{mk}_0(f+g) = \mathrm{mk}_0(f) + \mathrm{mk}_0(g)$ and
    $\Ext$-composition is additive in the right argument.
  \end{proof}
\end{lemma}

\begin{lemma}\leanok
[$F\tensor_S(-)$ is additive]
  \label{lem:tensor_base_change_additive}
  \lean{MR0555258CompactifyingPicard.tensorBaseChangeFunctor_additive}
  \uses{def:tensorBaseChangeFunctor, lem:tensorLeft_additive}
  \textit{Project-bespoke.}
  The base-change tensor functor $\mathtt{tensorBaseChangeFunctor}\,f\,F$ is
  additive.
  \begin{proof}
    It is the composite of the additive pullback $f^*$ (a left adjoint) and the
    additive $\mathtt{tensorLeft}\,F$ (\cref{lem:tensorLeft_additive}).
  \end{proof}
\end{lemma}

\begin{lemma}

\leanok
[$T$ is half-exact]
  \label{lem:functorT_halfexact}
  \lean{MR0555258CompactifyingPicard.functorT_halfExact}
  \uses{def:functorT, def:isHalfExact, lem:ext_tensor_halfexact_of_flat, lem:tilde_exact, lem:functorT_additive}
  % SOURCE: [Altman--Kleiman], §1, (1.3) proof p. 56 (read from references/MR0555258-compactifying-picard.pdf)
  % SOURCE QUOTE: "Consider the functor, T(M) = Ext^1_X(I, F (x)_S M~). from the
  % category of finitely generated A-modules M to itself. ... T is half-exact"
  \textit{Source: Altman--Kleiman, §1, (1.3) proof (``$T$ is half-exact'').}
  For $F$ $S$-flat, the functor $T = $ \cref{def:functorT} is half-exact: it sends
  every short exact sequence of finitely generated $A$-modules to a complex exact
  at its middle term.
  \begin{proof}
    \uses{lem:ext_tensor_halfexact_of_flat, lem:tilde_exact}
    Let $0 \to M' \to M \to M'' \to 0$ be a short exact sequence of $A$-modules.
    Applying $\widetilde{(-)}$ keeps it exact (\cref{lem:tilde_exact}), and $F$
    being $S$-flat keeps $F\tensor_S(-)$ exact, so we obtain a short exact sequence
    of $\Ox$-modules. The long exact sequence of $\Ext^\bullet_X(I,-)$ is then
    exact at $\Ext^1_X(I, F\tensor_S\widetilde M)$ — precisely the unconditional
    half-exactness of \cref{lem:ext_tensor_halfexact_of_flat} (degree $q=1$). Read
    through the $A$-module structure of $T$ (\cref{def:functorT}), this is exactness
    of $T(M') \to T(M) \to T(M'')$ at $T(M)$.
  \end{proof}
\end{lemma}

\begin{lemma}\leanok
[Homological core: $T(k(s)) = 0 \Rightarrow T(M) = 0$ for all $M$]
  \label{lem:functorT_eq_zero_of_residue}
  \lean{MR0555258CompactifyingPicard.functorT_eq_zero_of_residue}
  \uses{def:functorT, lem:functorT_halfexact, lem:functorT_additive, thm:ob_halfexact_free}
  % SOURCE: [Altman--Kleiman], §1, (1.3) proof p. 56 (read from references/MR0555258-compactifying-picard.pdf)
  % SOURCE QUOTE: "Since T is half-exact and since T(k(s)) is equal to zero, T(M)
  % is equal to zero for every finitely generated A-module M [OB, 2.1] or
  % [EGA III_2, 7.5.3])."
  \textit{Source: Altman--Kleiman, §1, (1.3) proof.}
  Let $A$ be a Noetherian local ring, $f : X \to \Spec A$, and $I, F$
  $\Ox$-modules with $F$ $S$-flat. If $T(k) = 0$ for the residue field $k$ of $A$
  (\cref{def:functorT}, $T = $ \texttt{functorT}), then $T(M) = 0$ for every
  finitely generated $A$-module $M$.
  \begin{proof}
    \uses{lem:functorT_halfexact, thm:ob_halfexact_free}
    The functor $T$ is half-exact (\cref{lem:functorT_halfexact}) and additive
    (\cref{lem:functorT_additive}); feeding this together with the residue-field
    vanishing $T(k) = 0$ to the [OB, 2.1] vanishing principle
    (\cref{thm:ob_halfexact_free}) gives $T(M) = 0$ for every finitely generated
    $A$-module $M$. This packages the homological heart of the (1.3) argument; the
    hypothesis $T(k) = 0$ is discharged separately by
    \cref{lem:ext1_vanishing_of_acyclic_chase} (brick (iii)).
  \end{proof}
\end{lemma}

\begin{lemma}

\leanok
[Diagram-chase vanishing of $\Ext^1$ via an acyclic surjection]
  \label{lem:ext1_vanishing_of_acyclic_chase}
  \lean{MR0555258CompactifyingPicard.ext1_vanishing_of_acyclic_chase}
  \uses{def:extGroup, lem:acyclic_surjection, thm:ega_adjunction}
  % SOURCE: [Altman--Kleiman], §1, (1.3) proof p. 56 (read from references/MR0555258-compactifying-picard.pdf)
  % SOURCE QUOTE: "Consider an exact sequence, 0 -> K -> J -> I -> 0, in which J is
  % as specified in (1.2). Since I is S-flat, the sequence remains exact when
  % restricted to X(s). So it yields a commutative diagram with exact rows, [...]
  % where j is the inclusion map of the closed fiber X(s) into X. The two vertical
  % maps are the adjunction isomorphisms. Now, j*I is equal to I(s), and
  % Ext^1_{X(s)}(I(s), F(s)) is equal to zero. Hence Ext^1_X(I, j_*F(s)) is equal
  % to zero. However, the latter Ext is just T(k(s))."
  \textit{Project-bespoke homological core formalizing the ``$T(k(s)) = 0$'' step
  of Altman--Kleiman, §1, (1.3) proof: the abstract diagram chase that makes
  $\Ext^1_X(I, Y) = 0$ fall out of an acyclic surjection and the surjectivity of
  the $\Hom$-level transition map.}
  Let $\mathcal{C}$ be an abelian category with enough Ext-theory
  (\texttt{HasExt}), let $0 \to K \to J \to I \to 0$ be a short exact sequence in
  $\mathcal{C}$, and let $Y$ be an object of $\mathcal{C}$. Suppose
  \begin{enumerate}
    \item $\Ext^1_{\mathcal{C}}(J, Y) = 0$ (the middle object $J$ is acyclic for
      $\Ext^1(-, Y)$), and
    \item the transition map $\Hom_{\mathcal{C}}(J, Y) \to \Hom_{\mathcal{C}}(K, Y)$
      (precomposition with $K \to J$, i.e. $\Ext^0$ in the first variable) is
      surjective.
  \end{enumerate}
  Then $\Ext^1_{\mathcal{C}}(I, Y) = 0$.
  \begin{proof}
    This is the abstract content of the diagram chase in the (1.3) proof, read off
    the contravariant long exact $\Ext$-sequence of the short exact sequence
    $0 \to K \to J \to I \to 0$ (Mathlib's
    \texttt{CategoryTheory.Abelian.Ext.contravariant\_sequence\_exact} family).
    Take any class $x \in \Ext^1(I, Y)$. Its image in $\Ext^1(J, Y)$ under the
    transition map (precomposition with $K \to J$) is $0$ by hypothesis (1), so by
    exactness of the sequence
    $\Ext^0(K, Y) \xrightarrow{\delta} \Ext^1(I, Y) \to \Ext^1(J, Y)$
    (\texttt{contravariant\_sequence\_exact₃}) there is
    $w \in \Ext^0(K, Y) = \Hom(K, Y)$ with $\delta(w) = x$, where $\delta$ is
    composition with the extension class of $0 \to K \to J \to I \to 0$. By
    hypothesis (2) write $w$ as the image of some $v \in \Hom(J, Y)$ under the
    transition map. But the composite ``extension class after the transition map''
    is zero (consecutive maps in the long exact sequence compose to zero — the
    $\texttt{.zero}$ of the relevant short complex), so
    $x = \delta(w) = \delta(\text{transition}(v)) = 0$. Hence every class in
    $\Ext^1(I, Y)$ vanishes, i.e. $\Ext^1(I, Y) = 0$.

    In the (1.3) application $\mathcal{C} = \Ox\text{-Mod}$,
    $Y = j_*F(s) = F\tensor_S\widetilde{k(s)}$, the surjection $J \to I$ is the
    acyclic one of \cref{lem:acyclic_surjection} (so (1) holds because $J$ is
    acyclic and $j_*F(s)$ is quasi-coherent), and (2) is obtained by transporting
    the surjectivity of $\Hom_{X(s)}(j^*J, F(s)) \to \Hom_{X(s)}(j^*K, F(s))$ — a
    consequence of $\Ext^1_{X(s)}(I(s), F(s)) = 0$ via the same contravariant
    sequence on the fibre $X(s)$ — across the degree-$0$ adjunction isomorphism
    $\Hom_X(-, j_*F(s)) \cong \Hom_{X(s)}(j^*-, F(s))$
    (\cref{thm:ega_adjunction}). The conclusion $\Ext^1_X(I, j_*F(s)) = 0$ is
    exactly $T(k(s)) = 0$ (\cref{def:functorT}).
  \end{proof}
\end{lemma}

% ---------------------------------------------------------------------------
% Coverage-debt helper blocks for the iter-027 chase decls (Lean↔blueprint 1:1).
% Project-bespoke (no external source): packaging lemmas around extGroup IsZero.
% ---------------------------------------------------------------------------

\begin{lemma}

\leanok
[$\mathtt{extGroup}$ vanishing as a membership condition]
  \label{lem:isZero_extGroup_iff}
  \lean{MR0555258CompactifyingPicard.isZero_extGroup_iff}
  \uses{def:extGroup}
  \textit{Project-bespoke bridge between the categorical $\mathtt{IsZero}$ predicate
  on $\mathtt{extGroup}\,q\,M\,N$ and the element-level statement.}
  For $q \in \mathbb{N}$ and objects $M, N$, the abelian group
  $\mathtt{extGroup}\,q\,M\,N$ (\cref{def:extGroup}) is a zero object if and only if
  every element of it equals $0$.
  \begin{proof}
    A (small) abelian group is a zero object iff it is a subsingleton; a
    subsingleton with a chosen zero is exactly the condition that every element
    equals $0$. (\texttt{AddCommGrpCat.isZero\_iff\_subsingleton} +
    \texttt{subsingleton\_iff\_forall\_eq}.)
  \end{proof}
\end{lemma}

\begin{lemma}

\leanok
[Scheme-level diagram-chase vanishing of $\Ext^1$]
  \label{lem:ext1_isZero_of_acyclic_chase}
  \lean{MR0555258CompactifyingPicard.ext1_isZero_of_acyclic_chase}
  \uses{lem:ext1_vanishing_of_acyclic_chase, lem:isZero_extGroup_iff, def:extGroup}
  \textit{Project-bespoke $\mathtt{extGroup}$-packaging of
  \cref{lem:ext1_vanishing_of_acyclic_chase} over a scheme's module category
  $Y\text{-Mod}$, the form the (1.3) assembly consumes.}
  Let $0 \to K \to J \to I \to 0$ be a short exact sequence of $\Oy$-modules and
  $W$ an $\Oy$-module. If $\mathtt{extGroup}\,1\,J\,W$ is a zero object and the
  transition map $\Hom_Y(J, W) \to \Hom_Y(K, W)$ is surjective, then
  $\mathtt{extGroup}\,1\,I\,W$ is a zero object.
  \begin{proof}
    Rewrite both the hypothesis and the goal with \cref{lem:isZero_extGroup_iff}
    into the membership form and apply \cref{lem:ext1_vanishing_of_acyclic_chase}.
    The $\mathtt{HasDerivedCategory.standard}$ / $\mathtt{HasExt}$ instances baked
    into $\mathtt{extGroup}$ make its carrier definitionally the $\mathtt{Abelian.Ext}$
    group on which the generic chase is stated.
  \end{proof}
\end{lemma}

\begin{lemma}

\leanok
[First-argument iso-invariance of $\mathtt{extGroup}$ vanishing]
  \label{lem:isZero_extGroup_of_iso_left}
  \lean{MR0555258CompactifyingPicard.isZero_extGroup_of_iso_left}
  \uses{lem:isZero_extGroup_iff, def:extGroup}
  \textit{Project-bespoke transport of $\mathtt{extGroup}$ vanishing along an
  isomorphism in the first argument.}
  If $\mathtt{extGroup}\,q\,M\,N$ is a zero object and $e : M \cong M'$ is an
  isomorphism, then $\mathtt{extGroup}\,q\,M'\,N$ is a zero object.
  \begin{proof}
    Transport an element $y \in \Ext^q(M', N)$ by
    $y = (\mathtt{mk_0}\,e^{-1}) \circ (\mathtt{mk_0}\,e) \circ y$
    (\texttt{comp\_assoc\_of\_second\_deg\_zero}, \texttt{mk\_0\_comp\_mk\_0},
    $e.\mathtt{inv\_hom\_id}$, \texttt{mk\_0\_id\_comp}); the inner factor
    $(\mathtt{mk_0}\,e) \circ y \in \Ext^q(M, N)$ vanishes by hypothesis, so
    $y = 0$ by \texttt{Ext.comp\_zero}.
  \end{proof}
\end{lemma}

\begin{lemma}\leanok
[Second-argument iso-invariance of $\mathtt{extGroup}$ vanishing]
  \label{lem:isZero_extGroup_of_iso_right}
  \lean{MR0555258CompactifyingPicard.isZero_extGroup_of_iso_right}
  \uses{lem:isZero_extGroup_iff, def:extGroup}
  \textit{Project-bespoke transport of $\mathtt{extGroup}$ vanishing along an
  isomorphism in the second argument (the covariant mirror of
  \cref{lem:isZero_extGroup_of_iso_left}).}
  If $\mathtt{extGroup}\,q\,M\,N$ is a zero object and $e : N \cong N'$ is an
  isomorphism, then $\mathtt{extGroup}\,q\,M\,N'$ is a zero object.
  \begin{proof}
    Transport an element $y \in \Ext^q(M, N')$ by
    $y = y \circ (\mathtt{mk_0}\,e^{-1}) \circ (\mathtt{mk_0}\,e)$
    (\texttt{comp\_assoc\_of\_third\_deg\_zero}, \texttt{mk\_0\_comp\_mk\_0},
    $e.\mathtt{inv\_hom\_id}$, \texttt{comp\_mk\_0\_id}); the factor
    $y \circ (\mathtt{mk_0}\,e^{-1}) \in \Ext^q(M, N)$ vanishes by hypothesis, so
    $y = 0$ by \texttt{Ext.zero\_comp}.
  \end{proof}
\end{lemma}

\begin{lemma}

\leanok
[$\Ext^1$ vanishing for the acyclic middle term]
  \label{lem:acyclic_ext1_vanishes}
  \lean{MR0555258CompactifyingPicard.acyclic_ext1_vanishes}
  \uses{lem:acyclic_surjection, lem:isZero_extGroup_of_iso_left, def:extGroup}
  \textit{Project-bespoke discharge of the acyclicity hypothesis ($\mathtt{hJ}$) of
  \cref{lem:ext1_isZero_of_acyclic_chase}: the middle term $J$ produced by
  \cref{lem:acyclic_surjection} is $\Ext^1$-acyclic.}
  For the acyclic object $J$ of \cref{lem:acyclic_surjection} and any
  quasi-coherent $\Ox$-module $W$, $\mathtt{extGroup}\,1\,J\,W$ is a zero object.
  \begin{proof}
    Specialize the acyclicity output of \cref{lem:acyclic_surjection} to the affine
    morphism $g = \mathrm{id}_X$; strip the resulting $\mathrm{id}_X^*$ pullback with
    $\mathtt{Scheme.Modules.pullbackId}$ and the iso-invariance
    \cref{lem:isZero_extGroup_of_iso_left}.
  \end{proof}
\end{lemma}

% ---------------------------------------------------------------------------
% iter-028 targets: the remaining buildable bricks discharging hsurj and
% assembling Ext^1_X(I, j_*F(s)) = 0 (all axiom-clean, no new axiom).
% ---------------------------------------------------------------------------

\begin{lemma}

\leanok
[Acyclic short exact sequence $0 \to K \to J \to I \to 0$]
  \label{lem:acyclic_ses}
  \lean{MR0555258CompactifyingPicard.acyclic_ses}
  \uses{lem:acyclic_surjection}
  % SOURCE: [Altman--Kleiman], §1, (1.3) proof p. 56 (read from references/MR0555258-compactifying-picard.pdf)
  % SOURCE QUOTE: "Consider an exact sequence, 0 -> K -> J -> I -> 0, in which J is
  % as specified in (1.2)."
  \textit{Source: Altman--Kleiman, §1, (1.3) proof p. 56 (``Consider an exact
  sequence $0 \to K \to J \to I \to 0$ in which $J$ is as specified in (1.2)'').}
  For an $\Ox$-module $I$, the acyclic surjection $\pi : J \to I$ of
  \cref{lem:acyclic_surjection} extends to a short exact sequence
  $0 \to K \to J \xrightarrow{\pi} I \to 0$ of $\Ox$-modules, where $K = \ker\pi$.
  \begin{proof}
    $\Ox\text{-Mod}$ is abelian; take $K = \ker\pi$ with inclusion the kernel
    morphism. The short complex $(\ker\pi \to J \xrightarrow{\pi} I)$ is short exact
    because $\pi$ is epi (\cref{lem:acyclic_surjection}), the kernel inclusion is
    mono, and the complex is exact at $J$ by the kernel universal property
    (\texttt{ShortComplex.ShortExact.mk} + \texttt{exact\_kernel}).
  \end{proof}
\end{lemma}

\begin{lemma}

\leanok
[Degree-$0$ $\Ext$ vs.\ $\Hom$ precomposition surjectivity]
  \label{lem:ext0_precomp_surjective_iff}
  \lean{MR0555258CompactifyingPicard.ext0_precomp_surjective_iff}
  \uses{def:extGroup}
  \textit{Project-bespoke bridge between the degree-$0$ $\Ext$ form (consumed by
  \cref{lem:ext1_isZero_of_acyclic_chase} via the contravariant LES) and the plain
  $\Hom$-set form (produced by the adjunction ladder).}
  For $\varphi : K \to J$ in $\Oy$-modules and a target $N$, precomposition with
  $\mathtt{mk_0}\,\varphi$ on $\Ext^0(J,N) \to \Ext^0(K,N)$ is surjective if and
  only if precomposition with $\varphi$ on $\Hom_Y(J,N) \to \Hom_Y(K,N)$ is
  surjective.
  \begin{proof}
    The bijection $\mathtt{mk_0} : (J \to N) \simeq \Ext^0(J,N)$
    (\texttt{Ext.mk\_0\_bijective}) conjugates the two precomposition maps:
    $\mathtt{mk_0}\,\varphi$ composed with $\mathtt{mk_0}\,v$ equals
    $\mathtt{mk_0}\,(\varphi \circ v)$ (\texttt{Ext.mk\_0\_comp\_mk\_0}), and
    $\mathtt{homEquiv_0}$ is its inverse. Hence the two surjectivity statements are
    equivalent.
  \end{proof}
\end{lemma}

\begin{lemma}

\leanok
[Fibre-side $\Hom$ surjectivity from $\Ext^1_{X(s)}(I(s),F(s)) = 0$]
  \label{lem:fibre_hom_surjective}
  \lean{MR0555258CompactifyingPicard.fibre_hom_surjective}
  \uses{lem:acyclic_ses, def:fiberExtVanishes, lem:ext0_precomp_surjective_iff, lem:isZero_extGroup_iff, def:extGroup}
  % SOURCE: [Altman--Kleiman], §1, (1.3) proof p. 56 (read from references/MR0555258-compactifying-picard.pdf)
  % SOURCE QUOTE: "Since I is S-flat, the sequence remains exact when restricted to
  % X(s). So it yields a commutative diagram with exact rows, [...] Now, j*I is equal
  % to I(s), and Ext^1_{X(s)}(I(s), F(s)) is equal to zero."
  \textit{Source: Altman--Kleiman, §1, (1.3) proof p. 56 (the fibre row of the
  comparison diagram).}
  Let $j = f.\mathtt{fiber}\iota\,s : X(s) \hookrightarrow X$ be the closed-fibre
  inclusion and assume $I$ is $S$-flat (\cref{def:isSFlat}). Restricting the short
  exact sequence $0 \to K \to J \to I \to 0$ (\cref{lem:acyclic_ses}) to $X(s)$ stays
  short exact (Altman--Kleiman: ``since $I$ is $S$-flat, the sequence remains exact
  when restricted to $X(s)$'' — the existing anchor \cref{thm:flat_pullback_exact}
  applied along $j$, with the $S$-flatness of the pulled-back right-hand term $j^*I$
  supplied by $I$ $S$-flat). With $\Ext^1_{X(s)}(j^*I, F(s)) = 0$
  (\cref{def:fiberExtVanishes} at $q = 1$), the contravariant $\Ext$ long exact
  sequence gives surjectivity of the transition map
  \[
    \Hom_{X(s)}(j^*J, F(s)) \longrightarrow \Hom_{X(s)}(j^*K, F(s)).
  \]
  \begin{proof}
    \uses{thm:flat_pullback_exact, def:fiberExtVanishes, def:isSFlat}
    The pulled-back sequence $0 \to j^*K \to j^*J \to j^*I \to 0$ is short exact by
    \cref{thm:flat_pullback_exact} (its $\mathtt{X_3} = j^*I$ is $S$-flat because $I$
    is). Take $w \in \Hom_{X(s)}(j^*K, F(s)) = \Ext^0$. Its
    image under the connecting map $\delta$ lies in $\Ext^1_{X(s)}(j^*I, F(s)) = 0$,
    hence is $0$; by exactness at the $\Ext^0(j^*K)$ spot
    (\texttt{Ext.contravariant\_sequence\_exact₁}, $n_0 = 0$, $n_1 = 1$) $w$ is the
    image of some element of $\Hom_{X(s)}(j^*J, F(s))$ under precomposition with
    $j^*(K \to J)$. Hence the transition map is surjective.
  \end{proof}
\end{lemma}

\begin{lemma}

\leanok
[Adjunction-ladder transport of $\Hom$ surjectivity]
  \label{lem:adjunction_ladder_surjective}
  \lean{MR0555258CompactifyingPicard.adjunction_ladder_surjective}
  \uses{thm:ega_adjunction, lem:fibre_hom_surjective}
  % SOURCE: [Altman--Kleiman], §1, (1.3) proof p. 56 (read from references/MR0555258-compactifying-picard.pdf)
  % SOURCE QUOTE: "it yields a commutative diagram with exact rows, [...] where j is
  % the inclusion map of the closed fiber X(s) into X. The two vertical maps are the
  % adjunction isomorphisms."
  \textit{Source: Altman--Kleiman, §1, (1.3) proof p. 56 (the two vertical
  adjunction isomorphisms of the comparison diagram).}
  Let $j : Y \to X$ be a morphism of schemes, $G$ a $\Oy$-module, and
  $\varphi : K \to J$ a morphism of $\Ox$-modules. If precomposition with $j^*\varphi$
  is surjective $\Hom_Y(j^*J, G) \to \Hom_Y(j^*K, G)$, then precomposition with
  $\varphi$ is surjective $\Hom_X(J, j_*G) \to \Hom_X(K, j_*G)$.
  \begin{proof}
    \uses{thm:ega_adjunction}
    The degree-$0$ adjunction $\Hom_X(-, j_*G) \cong \Hom_Y(j^*-, G)$
    (\cref{thm:ega_adjunction}, the hom-set equivalence of
    $\mathtt{pullbackPushforwardAdjunction}\,j$) is natural in the first argument
    (\texttt{Adjunction.homEquiv\_naturality\_left}), so the square comparing
    precomposition by $\varphi$ (top) and by $j^*\varphi$ (bottom) commutes.
    Conjugating the surjection of the bottom row by the two vertical bijections
    yields surjectivity of the top row.
  \end{proof}
\end{lemma}

\begin{lemma}

\leanok
[$\Ext^1_X(I, j_*F(s)) = 0$ (the $T(k(s)) = 0$ assembly, up to the fibre identification)]
  \label{lem:ext1_X_pushforward_fibre_vanishes}
  \lean{MR0555258CompactifyingPicard.ext1_X_pushforward_fibre_vanishes}
  \uses{lem:ext1_isZero_of_acyclic_chase, lem:acyclic_ses, lem:acyclic_ext1_vanishes, lem:adjunction_ladder_surjective, lem:fibre_hom_surjective, thm:ega_flat_fibre_restriction, def:fiberExtVanishes, def:isSFlat}
  % SOURCE: [Altman--Kleiman], §1, (1.3) proof p. 56 (read from references/MR0555258-compactifying-picard.pdf)
  % SOURCE QUOTE: "Hence Ext^1_X(I, j_*F(s)) is equal to zero. However, the latter
  % Ext is just T(k(s))."
  \textit{Source: Altman--Kleiman, §1, (1.3) proof p. 56 (``Hence
  $\Ext^1_X(I, j_*F(s)) = 0$'').}
  Let $f : X \to S$, $I, F$ be $\Ox$-modules with $I$ $S$-flat
  ($\mathtt{hI : IsSFlat}\,f\,I$, \cref{def:isSFlat}), $s \in S$ with
  $\Ext^1_{X(s)}(I(s), F(s)) = 0$ (\cref{def:fiberExtVanishes}). Then
  $\mathtt{extGroup}\,1\,I\,(j_*j^*F) = 0$, where $j = f.\mathtt{fiber}\iota\,s$.
  \begin{proof}
    \uses{thm:ega_flat_fibre_restriction}
    Apply \cref{lem:ext1_isZero_of_acyclic_chase} to the acyclic short exact sequence
    of \cref{lem:acyclic_ses} with $W = j_*j^*F$. The restriction of that sequence to
    the closed fibre stays short exact because $I$ is $S$-flat
    (\cref{thm:ega_flat_fibre_restriction}, the correct flatness mechanism — flatness
    of the base term $I$, not of the fibre map $j$). Its hypothesis $\mathtt{hJ}$
    ($\Ext^1_X(J, W) = 0$) is \cref{lem:acyclic_ext1_vanishes} ($W$ quasi-coherent as
    a pushforward of a pullback); its hypothesis $\mathtt{hsurj}$ is
    \cref{lem:adjunction_ladder_surjective} applied to $\varphi = (K \to J)$ and
    $G = j^*F = F(s)$, whose fibre-side surjectivity is
    \cref{lem:fibre_hom_surjective}. The conclusion is $\Ext^1_X(I, j_*F(s)) = 0$,
    which is $T(k(s)) = 0$ once the fibre identification $j_*j^*F \cong F\tensor_S
    \widetilde{k(s)}$ (\cref{thm:ega_fibre_tensor_residue}) is supplied.
  \end{proof}
\end{lemma}

\begin{lemma}\leanok
[$\mathtt{IsZero}$ reflection through $\mathtt{liftAdditiveToModuleCat}$]
  \label{lem:liftAdditiveToModuleCat_isZero_obj_iff}
  \lean{MR0555258CompactifyingPicard.liftAdditiveToModuleCat_isZero_obj_iff}
  \uses{def:functorT}
  \textit{Project-bespoke: the $A$-module lift of an additive functor reflects
  $\mathtt{IsZero}$ on objects.}
  For an additive functor $G$ into $\mathrm{Ab}$ and an object $M$,
  $(\mathtt{liftAdditiveToModuleCat}\,G).\mathtt{obj}\,M$ is a zero object of
  $\mathrm{Mod}_A$ if and only if $G.\mathtt{obj}\,M$ is a zero object of
  $\mathrm{Ab}$.
  \begin{proof}
    By construction $(\mathtt{liftAdditiveToModuleCat}\,G).\mathtt{obj}\,M =
    \mathtt{ModuleCat.of}\,A\,(G.\mathtt{obj}\,M)$ as the underlying additive
    group, so via \texttt{ModuleCat.isZero\_of\_iff\_subsingleton} and
    \texttt{AddCommGrpCat.isZero\_iff\_subsingleton} both sides reduce to
    $\mathtt{Subsingleton}\,(G.\mathtt{obj}\,M)$.
  \end{proof}
\end{lemma}

\begin{lemma}\leanok
[$\mathtt{IsZero}$ reflection through the $A$-module lift of $T$]
  \label{lem:functorT_obj_isZero_iff}
  \lean{MR0555258CompactifyingPicard.functorT_obj_isZero_iff}
  \uses{def:functorT, def:extGroup, def:tensorBaseChangeFunctor, lem:liftAdditiveToModuleCat_isZero_obj_iff}
  \textit{Project-bespoke bridge: $(\texttt{functorT}\,A\,f\,I\,F).\mathtt{obj}\,M$ is a zero
  object of $\mathrm{Mod}_A$ if and only if the underlying $\mathtt{extGroup}$ is a
  zero object of $\mathrm{Ab}$.}
  For $A$ a commutative ring, $f : X \to \Spec A$, $\Ox$-modules $I, F$, and an
  $A$-module $M$, $(\texttt{functorT}\,A\,f\,I\,F).\mathtt{obj}\,M$ is a zero object iff
  $\mathtt{extGroup}\,1\,I\,\big((\mathtt{tensorBaseChangeFunctor}\,f\,F).\mathtt{obj}\,(\widetilde M)\big)$
  is a zero object.
  \begin{proof}
    $\texttt{functorT}$ is the End-action lift (\cref{def:functorT}) of the composite
    $\widetilde{(-)} \circ (\mathtt{tensorBaseChangeFunctor}\,f\,F) \circ
    (\texttt{covariantExtFunctor}\,I\,1)$; the lift only re-structures the underlying
    abelian group with an $A$-action, so $\mathtt{IsZero}$ of the lifted
    $\mathrm{Mod}_A$ object is equivalent to $\mathtt{IsZero}$ (subsingleton) of the
    underlying $\mathtt{extGroup}$, which is exactly the stated $\Ext^1$ group.
  \end{proof}
\end{lemma}

\begin{lemma}\leanok
[$T(k(s)) = 0$: the residue value of the functor $T$]
  \label{lem:functorT_residue_eq_zero}
  \lean{MR0555258CompactifyingPicard.functorT_residue_eq_zero}
  \uses{def:functorT, lem:ext1_X_pushforward_fibre_vanishes, thm:ega_fibre_tensor_residue, thm:ega_fibre_pushforward_qcoh, lem:functorT_obj_isZero_iff, lem:isZero_extGroup_iff, lem:isZero_extGroup_of_iso_right, def:isSFlat}
  % SOURCE: [Altman--Kleiman], §1, (1.3) proof p. 56 (read from references/MR0555258-compactifying-picard.pdf)
  % SOURCE QUOTE: "Now, j*I is equal to I(s), and Ext^1_{X(s)}(I(s), F(s)) is equal
  % to zero. Hence Ext^1_X(I, j_*F(s)) is equal to zero. However, the latter Ext is
  % just T(k(s))."
  \textit{Source: Altman--Kleiman, §1, (1.3) proof p. 56 (``the latter Ext is just
  $T(k(s))$'').}
  Let $A$ be a Noetherian local ring with closed point $s$ and residue field
  $k = k(s)$, $f : X \to \Spec A$, and $I, F$ $\Ox$-modules with $I$ $S$-flat and
  $\Ext^1_{X(s)}(I(s),F(s)) = 0$. Then $(\texttt{functorT}\,A\,f\,I\,F).\mathtt{obj}\,k = 0$,
  i.e.\ $T(k(s)) = 0$.
  \begin{proof}
    By \cref{lem:functorT_obj_isZero_iff} it suffices to show
    $\mathtt{extGroup}\,1\,I\,\big(F\tensor_S \widetilde k\big) = 0$. The fibre
    identification \cref{thm:ega_fibre_tensor_residue} gives
    $F\tensor_S \widetilde k \cong j_*j^*F$ for $j$ the closed-fibre inclusion at
    $s$, and transporting along this isomorphism in the second $\Ext$ argument
    reduces the claim to $\mathtt{extGroup}\,1\,I\,(j_*j^*F) = 0$, which is
    \cref{lem:ext1_X_pushforward_fibre_vanishes} (using $I$ $S$-flat and the fibre
    $\Ext$-vanishing hypothesis).
  \end{proof}
\end{lemma}

\begin{lemma}\leanok
[$T(M) = 0$ for every finitely generated $A$-module (homological heart at $\Spec A$)]
  \label{lem:functorT_eq_zero}
  \lean{MR0555258CompactifyingPicard.functorT_eq_zero}
  \uses{def:functorT, lem:functorT_residue_eq_zero, lem:functorT_eq_zero_of_residue, def:isSFlat}
  % SOURCE: [Altman--Kleiman], §1, (1.3) proof p. 56 (read from references/MR0555258-compactifying-picard.pdf)
  % SOURCE QUOTE: "Since T is half-exact and since T(k(s)) is equal to zero, T(M)
  % is equal to zero for every finitely generated A-module M [OB, 2.1] or [EGA
  % III_2, 7.5.3])."
  \textit{Source: Altman--Kleiman, §1, (1.3) proof p. 56.}
  Let $A$ be a Noetherian local ring with closed point $s$, $f : X \to \Spec A$,
  and $I, F$ $\Ox$-modules with $I, F$ $S$-flat and $\Ext^1_{X(s)}(I(s),F(s)) = 0$.
  Then $(\texttt{functorT}\,A\,f\,I\,F).\mathtt{obj}\,M = 0$ for every finitely
  generated $A$-module $M$.
  \begin{proof}
    Combine the residue vanishing $T(k(s)) = 0$ (\cref{lem:functorT_residue_eq_zero})
    with the half-exact upgrade \cref{lem:functorT_eq_zero_of_residue} (which
    invokes \cref{thm:ob_halfexact_free}): a half-exact functor over a Noetherian
    local ring with vanishing residue value vanishes on all finitely generated
    modules.
  \end{proof}
\end{lemma}

%==============================================================================
% Brick 6 of (1.3): from the homological heart T(M) = 0 (lem:functorT_eq_zero)
% to local freeness of H. The free conclusion at S = Spec A is BUILT (the
% commutative-algebra endpoint Module.free_of_flat_of_isLocalRing is in Mathlib);
% the EGA reduction from a general base S to Spec A (Noetherian local, closed
% point) is anchored as standard EGA (O_I 2.4.1, 5.4.1; IV_3 Sect. 8; 1.1
% base-change), since the base-change compatibility of H it requires is not built
% project-side (H_tensor / H_tensor_invertible are paused).
%==============================================================================

\begin{lemma}

\leanok
[Exactness of $M \mapsto \Hom_X(I, F \tensor_S \widetilde M)$ from $T = 0$]
  \label{lem:homTensorFunctor_tilde_exact}
  \lean{MR0555258CompactifyingPicard.homTensorFunctor_tilde_exact}
  \uses{def:functorT, lem:functorT_eq_zero, def:homTensorFunctor, thm:flat_tensor_exact, def:isSFlat}
  % SOURCE: [Altman--Kleiman], §1, (1.3) proof p. 56 (read from references/MR0555258-compactifying-picard.pdf)
  % SOURCE QUOTE: "Therefore the functor M |-> Hom_X(I, F (x)_S M~) is exact."
  \textit{Source: Altman--Kleiman, §1, (1.3) proof p. 56.}
  Under the hypotheses of \cref{lem:functorT_eq_zero}, the functor
  $M \mapsto \Hom_X(I, F \tensor_S \widetilde M)$ on finitely generated
  $A$-modules carries short exact sequences to short exact sequences (it is exact).
  \begin{proof}
    \uses{lem:functorT_eq_zero, thm:flat_tensor_exact}
    Let $0 \to M' \to M \to M'' \to 0$ be a short exact sequence of $A$-modules.
    Applying $\widetilde{(-)}$ (\cref{lem:tilde_exact}) and the flat base change
    $F \tensor_S (-)$ (\cref{thm:flat_tensor_exact}, $F$ being $S$-flat) yields a
    short exact sequence
    $0 \to F\tensor_S\widetilde{M'} \to F\tensor_S\widetilde M \to
    F\tensor_S\widetilde{M''} \to 0$ of $\Ox$-modules. The covariant $\Ext$ long
    exact sequence for $\Hom_X(I, -) = \Ext^0_X(I, -)$ reads
    \[
      0 \to \Hom_X(I, F\tensor_S\widetilde{M'}) \to
      \Hom_X(I, F\tensor_S\widetilde M) \to
      \Hom_X(I, F\tensor_S\widetilde{M''}) \to
      \Ext^1_X(I, F\tensor_S\widetilde{M'}).
    \]
    The last term is $(\texttt{functorT}\,A\,f\,I\,F).\mathtt{obj}\,M' = 0$ by
    \cref{lem:functorT_eq_zero}, so the functor is right exact; it is left exact as
    $\Hom_X(I,-)$ always is. Hence it is exact.
  \end{proof}
\end{lemma}

\begin{theorem}[Local finite presentation of $H$ (EGA finiteness)]
  \label{thm:ega_H_isLFP}
  \lean{MR0555258CompactifyingPicard.External.H_isLFP}
  \uses{def:H, def:isLFP, def:admissible}
  % NOTE (iter-030): the finiteness of the representing module H, used in (1.3)
  % both for the affine bridge (H = Gamma(H)~ over Spec A) of the free conclusion
  % and for the EGA O_I 5.4.1 free-at-point criterion. AK: "it suffices to show
  % H(I,F) is free at s, because it is locally finitely presented [EGA O_I,
  % 5.4.1]." Finiteness of H is [EGA III_2, 7.7.6] (coherence of the corepresenting
  % module over a Noetherian base); absent from Mathlib. Small EGA anchor.
  % SOURCE: [Altman--Kleiman], §1, (1.3) proof p. 56 (read from references/MR0555258-compactifying-picard.pdf)
  % SOURCE QUOTE: "Finally it suffices to show H(I, F) is free at s, because it is
  % locally finitely presented [EGA O_I, 5.4.1]."
  \textit{Source: Altman--Kleiman, §1, (1.3) proof p. 56; finiteness is [EGA
  III$_2$, 7.7.6].}
  For an admissible $f : X \to S$ over a Noetherian base and $\Ox$-modules $I, F$,
  the representing module $H(I,F)$ is locally finitely presented (\cref{def:isLFP}).
  \begin{proof}
    Standard EGA finiteness of the corepresenting module over a Noetherian base;
    used here as an external input.
  \end{proof}
\end{theorem}

\begin{lemma}\leanok
[Natural isomorphisms transport map-surjectivity]
  \label{lem:surjective_of_natIso_map}
  \lean{MR0555258CompactifyingPicard.surjective_of_natIso_map}
  Let $\alpha : F \cong G$ be a natural isomorphism of functors $C \to \mathbf{Type}$
  and $\varphi$ a morphism of $C$. If $G(\varphi)$ is surjective then so is
  $F(\varphi)$. (Generic helper, $\mathbf{Type}$-valued; isolates the
  concrete-category bookkeeping of the transport core
  \cref{lem:hom_postcomp_surjective_of_functorT_zero}.)
  \begin{proof}
    Naturality of $\alpha$ gives $F(\varphi) = \alpha_{Y}^{-1} \circ G(\varphi)
    \circ \alpha_{X}$; the outer maps are bijections, so $F(\varphi)$ is surjective
    iff $G(\varphi)$ is.
  \end{proof}
\end{lemma}

\begin{lemma}\leanok
[Postcomposition surjectivity from $T = 0$ (transport core)]
  \label{lem:hom_postcomp_surjective_of_functorT_zero}
  \lean{MR0555258CompactifyingPicard.hom_postcomp_surjective_of_functorT_zero}
  \uses{thm:H_represents, lem:homTensorFunctor_tilde_exact, lem:surjective_of_natIso_map, thm:stacks_affine_fp_tilde}
  In the situation of \cref{thm:H_free_of_functorT_zero}, let $N$ be a finitely
  presented $A$-module with $\widetilde N \cong \HIF(I,F)$ and let
  $0 \to \cdot \to \cdot \to N \to 0$ be a short exact sequence of $A$-modules with
  surjection $g$. Then postcomposition
  $(\varphi \mapsto \varphi \circ g) : \Hom_A(N, \cdot) \to \Hom_A(\cdot, \cdot)$ is
  surjective. (The representability + tilde-full-faithfulness transport core of
  \cref{thm:H_free_of_functorT_zero}.)
  \begin{proof}
    \uses{thm:H_represents, lem:homTensorFunctor_tilde_exact, lem:surjective_of_natIso_map}
    By \cref{lem:homTensorFunctor_tilde_exact} postcomposition is surjective on
    $\Hom_X(I, F\tensor_S\widetilde{(-)})$. Whiskering the representability iso
    \cref{thm:H_represents} by $\widetilde{(-)}$ and transporting along it
    (\cref{lem:surjective_of_natIso_map}) moves this to $\mathrm{coyoneda}$ of
    $\HIF(I,F) \cong \widetilde N$; full faithfulness of $\widetilde{(-)}$
    (\cref{thm:stacks_affine_fp_tilde}) descends it to $\Hom_A(N, -)$.
  \end{proof}
\end{lemma}

\begin{theorem}

\leanok
[Freeness of $H$ at $\Spec A$ from $T = 0$]
  \label{thm:H_free_of_functorT_zero}
  \lean{MR0555258CompactifyingPicard.H_free_of_functorT_zero}
  \uses{def:H, thm:H_represents, lem:homTensorFunctor_tilde_exact, lem:hom_postcomp_surjective_of_functorT_zero, thm:ega_H_isLFP, thm:stacks_affine_fp_tilde, def:isFreeMod}
  % SOURCE: [Altman--Kleiman], §1, (1.3) proof p. 56 (read from references/MR0555258-compactifying-picard.pdf)
  % SOURCE QUOTE: "Thus the functor M |-> Hom_S(H(I,F), M~) is exact. Hence H(I,F)
  % is free."
  \textit{Source: Altman--Kleiman, §1, (1.3) proof p. 56.}
  Let $A$ be a Noetherian local ring, $f : X \to \Spec A$ admissible, and $I, F$
  $\Ox$-modules with $I, F$ $S$-flat and $\Ext^1_{X(s)}(I(s),F(s)) = 0$ at the
  closed point $s$. Then $H(I,F)$ is a free $\mathcal{O}_{\Spec A}$-module
  (\cref{def:isFreeMod}).
  \begin{proof}
    \uses{lem:homTensorFunctor_tilde_exact, thm:H_represents}
    By \cref{lem:homTensorFunctor_tilde_exact} the functor
    $M \mapsto \Hom_X(I, F\tensor_S\widetilde M)$ is exact; by representability
    \cref{thm:H_represents} it is naturally isomorphic to
    $M \mapsto \Hom_S(H(I,F), \widetilde M)$, which is therefore exact. As $H(I,F)$
    is locally finitely presented (\cref{thm:ega_H_isLFP}) and $S = \Spec A$ is
    affine, the affine bridge \cref{thm:stacks_affine_fp_tilde} gives a finitely
    presented $A$-module $N$ with $H(I,F) \cong \widetilde N$, and full faithfulness
    of $\widetilde{(-)}$ identifies $\Hom_S(\widetilde N, \widetilde M)$ with
    $\Hom_A(N, M)$. Exactness of $\Hom_A(N, -)$ means $N$ is projective, hence flat
    (\texttt{Module.Flat.of\_projective}); $N$ is finite and $A$ is local, so $N$ is
    free (\texttt{Module.free\_of\_flat\_of\_isLocalRing}). Therefore
    $\widetilde N \cong H(I,F)$ is a free $\mathcal{O}_{\Spec A}$-module.
  \end{proof}
\end{theorem}

\begin{theorem}[Locally free of finite rank from free at a point (EGA O$_I$ 5.4.1, 2.4.1)]
  \label{thm:ega_lfp_free_at_point}
  \lean{MR0555258CompactifyingPicard.External.lfp_free_at_point_locallyFree}
  \uses{def:isLFP, def:isFreeMod, def:isLocallyFreeOfFiniteRankOn, def:isRetrocompact}
  % NOTE (iter-030): the EGA O_I 5.4.1 / 2.4.1 packaging used at the very end of
  % the (1.3) reduction: a locally finitely presented module that is free at a
  % point s is locally free of finite rank on an open retrocompact neighbourhood of
  % s. ("Retrocompact means the inclusion map is quasi-compact [EGA O_I, 2.4.1]";
  % "it suffices to show H is free at s, because it is locally finitely presented
  % [EGA O_I, 5.4.1]".) Standard EGA; absent from Mathlib. Sound general anchor (no
  % section-1 content).
  % SOURCE: [Altman--Kleiman], §1, (1.3) proof p. 56 (read from references/MR0555258-compactifying-picard.pdf)
  % SOURCE QUOTE: "Retrocompact means the inclusion map is quasi-compact [EGA O_I,
  % 2.4.1]. ... Finally it suffices to show H(I, F) is free at s, because it is
  % locally finitely presented [EGA O_I, 5.4.1]."
  \textit{Source: Altman--Kleiman, §1, (1.3) proof p. 56; [EGA O$_I$, 5.4.1, 2.4.1].}
  Let $M$ be a locally finitely presented $\mathcal{O}_S$-module (\cref{def:isLFP})
  that is free at a point $s \in S$ (its restriction to some affine open
  neighbourhood of $s$ is free of finite rank). Then there is an open, retrocompact
  neighbourhood $U$ of $s$ with $M|U$ locally free of finite rank
  (\cref{def:isLocallyFreeOfFiniteRankOn}, \cref{def:isRetrocompact}).
  \begin{proof}
    Standard EGA: a finitely presented module is free in a neighbourhood of any
    point where it is free (O$_I$ 5.4.1), and every open of $S$ has quasi-compact
    inclusion when restricted suitably (O$_I$ 2.4.1, retrocompactness). Used as an
    external input.
  \end{proof}
\end{theorem}

\begin{theorem}[Reduction of (1.3) to a Noetherian local base (EGA descent + base change)]
  \label{thm:ega_reduce_H_locallyFree}
  \lean{MR0555258CompactifyingPicard.External.reduce_H_locallyFree_to_spec_local}
  \uses{def:H, thm:H_represents, thm:ega_descent_noetherian, def:admissible, def:isSFlat, def:fiberExtVanishes, def:isLocallyFreeOfFiniteRankOn, def:isRetrocompact}
  % NOTE (iter-030): the EGA reduction step of (1.3). The assertion is local on S
  % (so assume S affine); descent to a Noetherian base (thm:ega_descent_noetherian,
  % [EGA IV_3 Sect. 8]) together with the base-change compatibility of H ((1.1))
  % reduces to S = Spec A with A Noetherian local and s its closed point. This
  % reduction transports H along base change, which requires the base-change
  % compatibility of H that is NOT built project-side (H_tensor / H_tensor_invertible
  % are paused), so it is anchored as standard EGA. It carries NO section-1 content:
  % the section-1 argument is the input "H is free at the closed point of Spec A"
  % (thm:H_free_of_functorT_zero, via the functor T), supplied as the premise.
  % SOURCE: [Altman--Kleiman], §1, (1.3) proof p. 56 (read from references/MR0555258-compactifying-picard.pdf)
  % SOURCE QUOTE: "The assertion is clearly local on S, so we may assume S is
  % affine. It then follows from [EGA IV_3, Sect. 8] and the compatibility of H(I,
  % F) with base-change (1.1) that we may assume S is Noetherian. ... So we may
  % assume S is the spectrum of a Noetherian local ring A and s is the closed point
  % of S."
  \textit{Source: Altman--Kleiman, §1, (1.3) proof p. 56; [EGA IV$_3$, Sect. 8].}
  Let $f : X \to S$ be admissible with $I, F$ $S$-flat and
  $\Ext^1_{X(s)}(I(s),F(s)) = 0$ at $s \in S$. Suppose that for every admissible
  $f' : X' \to \Spec A$ over a Noetherian local ring $A$ (closed point $s'$) with
  $S$-flat $I', F'$ and fibre $\Ext$ vanishing at $s'$, the representing module
  $H(I',F')$ is free. Then there is an open, retrocompact neighbourhood $U$ of $s$
  with $H(I,F)|U$ locally free of finite rank.
  \begin{proof}
    Standard EGA reduction: localness on $S$, descent to a Noetherian base
    (\cref{thm:ega_descent_noetherian}) with base-change compatibility of $H$
    (\cref{thm:H_represents}, (1.1)), localisation to a Noetherian local ring, and
    the free-at-a-point criterion \cref{thm:ega_lfp_free_at_point}. Used as an
    external input. (The premise — freeness of $H$ at the closed point of a
    Noetherian local base — is exactly \cref{thm:H_free_of_functorT_zero}, the
    section-1 output.)
  \end{proof}
\end{theorem}

\begin{theorem}\leanok
[Local freeness of $H(I,F)$]
  \label{thm:H_locallyFree}
  \lean{MR0555258CompactifyingPicard.H_locallyFree_of_ext_vanishing}
  \uses{def:H, def:isRetrocompact, def:isHalfExact, thm:H_represents, lem:acyclic_surjection, lem:ext_tensor_halfexact, lem:ext_tensor_halfexact_of_flat, def:functorT, lem:functorT_additive, lem:functorT_halfexact, lem:functorT_eq_zero_of_residue, lem:functorT_residue_eq_zero, lem:functorT_eq_zero, lem:homTensorFunctor_tilde_exact, thm:H_free_of_functorT_zero, thm:ega_H_isLFP, thm:ega_lfp_free_at_point, thm:ega_reduce_H_locallyFree, lem:ext1_vanishing_of_acyclic_chase, lem:ext1_X_pushforward_fibre_vanishes, thm:ega_flat_fibre_restriction, thm:ega_fibre_tensor_residue, thm:flat_tensor_exact, thm:ob_halfexact_free, thm:ega_descent_noetherian, thm:ega_adjunction}
  % SOURCE: [Altman--Kleiman], §1, (1.3), p. 55--56 (read from references/MR0555258-compactifying-picard.pdf)
  % SOURCE QUOTE: "(1.3) THEOREM. Let f : X -> S be a finitely presented, proper
  % morphism of schemes, and let I and F be locally finitely presented, S-flat
  % O_X-modules. Assume the relation, Ext^1_{X(s)}(I(s), F(s)) = 0, holds for
  % some point s in S. Then there exists an open, retrocompact neighborhood U of
  % s such that H(I, F) | U is locally free with a finite rank."
  \textit{Source: Altman--Kleiman, §1, (1.3), p. 55--56.}
  Let $f : X \to S$ be a finitely presented, proper morphism, and let $I$, $F$ be
  locally finitely presented, $S$-flat $\Ox$-modules. If
  \[
    \Ext^1_{X(s)}\!\big(I(s), F(s)\big) = 0
  \]
  for some point $s \in S$, then there is an open, retrocompact neighbourhood $U$
  of $s$ such that $\HIF(I,F)|U$ is locally free of finite rank.
  % SOURCE QUOTE PROOF: "Proof. Retrocompact means the inclusion map is
  % quasi-compact [EGA O_I, 2.4.1]. Obviously the notion is stable under
  % base-change and obviously every subset of a locally Noetherian space is
  % retrocompact. The assertion is clearly local on S, so we may assume S is
  % affine. It then follows from [EGA IV_3, Sect. 8] and the compatibility of
  % H(I, F) with base-change (1.1) that we may assume S is Noetherian. Finally it
  % suffices to show H(I, F) is free at s, because it is locally finitely
  % presented [EGA O_I, 5.4.1]. So we may assume S is the spectrum of a
  % Noetherian local ring A and s is the closed point of S. Consider the functor,
  % T(M) = Ext^1_X(I, F (x)_S M~). from the category of finitely generated
  % A-modules M to itself. (Note that T(M) is finitely generated because f is
  % proper and S is Noetherian [GD, IV, 3.2, p. 74].) We shall now show that
  % T(k(s)) is equal to zero. Consider an exact sequence, 0 -> K -> J -> I -> 0,
  % in which J is as specified in (1.2). Since I is S-flat, the sequence remains
  % exact when restricted to X(s). So it yields a commutative diagram with exact
  % rows, [...] where j is the inclusion map of the closed fiber X(s) into X. The
  % two vertical maps are the adjunction isomorphisms. Now, j*I is equal to I(s),
  % and Ext^1_{X(s)}(I(s), F(s)) is equal to zero. Hence Ext^1_X(I, j_*F(s)) is
  % equal to zero. However, the latter Ext is just T(k(s)). Since T is half-exact
  % and since T(k(s)) is equal to zero, T(M) is equal to zero for every finitely
  % generated A-module M [OB, 2.1] or [EGA III_2, 7.5.3]). Therefore the functor
  % M |-> Hom_X(I, F (x)_S M~) is exact. Thus the functor M |-> Hom_S(H(I,F), M~)
  % is exact. Hence H(I,F) is free."
  \begin{proof}
    \uses{thm:ega_reduce_H_locallyFree, thm:H_free_of_functorT_zero, lem:functorT_eq_zero, lem:functorT_residue_eq_zero, lem:homTensorFunctor_tilde_exact, thm:ega_H_isLFP, thm:ega_lfp_free_at_point, def:isSFlat}
    \emph{Reduction (\cref{thm:ega_reduce_H_locallyFree}).} The notion of
    retrocompact (the inclusion is quasi-compact) is stable under base change, and
    every subset of a locally Noetherian space is retrocompact; the assertion is
    local on $S$, so we may assume $S$ affine. Using \cref{thm:ega_descent_noetherian}
    together with the base-change compatibility of $\HIF(I,F)$
    (\cref{thm:H_represents}) we reduce to $S$ Noetherian, and since $\HIF(I,F)$ is
    locally finitely presented (\cref{thm:ega_H_isLFP}) the free-at-a-point criterion
    \cref{thm:ega_lfp_free_at_point} shows it suffices to prove $\HIF(I,F)$ free at
    $s$; so let $S = \Spec A$ with $A$ Noetherian local and $s$ its closed point.
    This whole reduction is the standard EGA step \cref{thm:ega_reduce_H_locallyFree}.

    \emph{Homological heart ($T(k(s)) = 0 \Rightarrow T \equiv 0$).} Consider
    $T(M) = \Ext^1_X(I, F \tensor_S \widetilde M)$ as a functor on finitely
    generated $A$-modules (\cref{def:functorT}). By \cref{lem:functorT_residue_eq_zero}
    (whose proof restricts the acyclic sequence $0 \to K \to J \to I \to 0$ of
    \cref{lem:acyclic_surjection} to the closed fibre using $S$-flatness of $I$,
    runs the contravariant $\Ext$ diagram chase, and identifies the result with the
    residue value) we have $T(k(s)) = 0$; as $T$ is half-exact,
    \cref{lem:functorT_eq_zero} (via \cref{thm:ob_halfexact_free}) gives $T(M) = 0$
    for every finitely generated $A$-module $M$.

    \emph{Free conclusion (\cref{thm:H_free_of_functorT_zero}).} By
    \cref{lem:homTensorFunctor_tilde_exact} the functor
    $M \mapsto \Hom_X(I, F \tensor_S \widetilde M)$ is exact, hence so is
    $M \mapsto \Hom_S(\HIF(I,F), \widetilde M)$ (\cref{thm:H_represents}); over the
    affine $\Spec A$ this forces the finitely presented $\HIF(I,F)$ to be free
    (\cref{thm:H_free_of_functorT_zero}). Feeding this back through
    \cref{thm:ega_reduce_H_locallyFree} produces the required open, retrocompact
    neighbourhood $U$ of $s$ on which $\HIF(I,F)$ is locally free of finite rank.
  \end{proof}
\end{theorem}

%==============================================================================
% Affine bridge (project infrastructure): the equivalence between
% quasi-coherent sheaves on an affine scheme and modules over its global
% sections, realized through Mathlib's tilde construction. These are the
% project-local Mathlib-gap lemmas that turn module-level finite-presentation
% facts into sheaf-level ones; they feed the free-resolution step (1.4) and,
% downstream, the relative-Ext layer. Standard (Hartshorne II.5); stated here
% as project infrastructure (no external paper source — they support 1.4).
%==============================================================================

\begin{definition}[Tilde / affine quasi-coherent sheaf]
  \label{def:tildeMod}
  \lean{AlgebraicGeometry.tilde}
  \mathlibok
  % NOTE: retargeted \lean{ModuleCat.tilde} -> \lean{AlgebraicGeometry.tilde} (iter-012);
  % ModuleCat.tilde is now a @[deprecated] alias of AlgebraicGeometry.tilde in current Mathlib.
  For a commutative ring $A$ and an $A$-module $M$, the associated
  quasi-coherent $\mathcal{O}_{\Spec A}$-module $\widetilde M$ on $\Spec A$.
  Realized by Mathlib's \texttt{ModuleCat.tilde}; its stalk at a prime
  $\mathfrak p$ is the localization $M_{\mathfrak p}$
  (\texttt{ModuleCat.Tilde.stalkIso}).
\end{definition}

\begin{lemma}\leanok
[Tilde of a finite free module is a finite free sheaf]
  \label{lem:tilde_free}
  \lean{MR0555258CompactifyingPicard.tilde_free}
  \uses{def:tildeMod, def:isFreeMod}
  For a finite index type $\Lambda$, the tilde $\widetilde{A^{(\Lambda)}}$ of the
  free $A$-module on $\Lambda$ is isomorphic, as an
  $\mathcal{O}_{\Spec A}$-module, to the finite free sheaf
  $\mathcal{O}_{\Spec A}^{(\Lambda)}$ (\texttt{SheafOfModules.free}). In
  particular it satisfies \cref{def:isFreeMod}.
  \begin{proof}
    \uses{def:tildeMod}
    Tilde is additive and sends the structure sheaf $A = \widetilde A$ to the
    unit $\mathcal{O}_{\Spec A}$ and preserves the finite coproduct indexed by
    $\Lambda$ (a finite biproduct, which tilde, being additive, carries to the
    corresponding biproduct of copies of the unit). Identifying the latter with
    \texttt{SheafOfModules.free}~$\Lambda$ gives the isomorphism.
  \end{proof}
\end{lemma}

\begin{lemma}\leanok
[Tilde is exact]
  \label{lem:tilde_exact}
  \lean{MR0555258CompactifyingPicard.tilde_exact}
  \uses{def:tildeMod}
  % NOTE: PROVEN axiom-clean (iter-013). The "absent stalk-exactness detection layer"
  % collapsed once Mathlib.Algebra.Homology.ShortComplex.ExactFunctor was found: tilde is
  % already right exact (left adjoint), so left-exactness reduces to PreservesMonomorphisms
  % (tilde.functor R), supplied by tilde_preservesMono (proved stalkwise via
  % stalkFunctor_map_injective_of_isBasis + IsLocalizedModule.map_injective). Then
  % ShortExact.map_of_exact closes it.
  A short exact sequence $0 \to K \to N \to M \to 0$ of $A$-modules induces a
  short exact sequence
  $0 \to \widetilde K \to \widetilde N \to \widetilde M \to 0$ of
  $\mathcal{O}_{\Spec A}$-modules.
  \begin{proof}\leanok
    \uses{def:tildeMod, lem:tilde_preservesFiniteLimits}
    It suffices to show the tilde functor preserves short exact sequences. The
    epimorphism $\widetilde N \to \widetilde M$ is automatic: tilde is a left
    adjoint, hence preserves all colimits, in particular cokernels and
    epimorphisms. The remaining content is left-exactness, i.e. that tilde
    preserves the kernel $\widetilde K = \ker(\widetilde N \to \widetilde M)$;
    equivalently that $\widetilde K \to \widetilde N$ is a monomorphism with the
    sequence exact in the middle.

    Both are local statements verified on stalks. The stalk of $\widetilde{(-)}$
    at a prime $\mathfrak p$ is localization $(-)_{\mathfrak p}$ at
    $\mathfrak p$ (the stalk isomorphism), and localization of $A$-modules is an
    exact functor: it preserves injectivity and exactness in the middle. Hence at
    every point the stalk sequence
    $0 \to K_{\mathfrak p} \to N_{\mathfrak p} \to M_{\mathfrak p} \to 0$ is
    short exact. A sequence of sheaves of modules is exact precisely when it is
    exact on all stalks, so the sheaf sequence is short exact, completing the
    short-exactness. (Monomorphism may also be obtained sectionwise: a map of
    (pre)sheaves of modules is a monomorphism iff it is injective on each
    section, and the sections of $\widetilde{(-)}$ over a distinguished open
    $D(f)$ are the localizations $(-)_f$, on which an injective module map stays
    injective.)
  \end{proof}
\end{lemma}

% ---- supporting lemmas for tilde_exact (coverage of prover helpers) ----

\begin{lemma}\leanok
[Tilde preserves monomorphisms]
  \label{lem:tilde_mono}
  \lean{MR0555258CompactifyingPicard.tilde_mono}
  \uses{def:tildeMod}
  For an injective $R$-linear map $f : M \to N$ the induced map of
  $\mathcal{O}_{\Spec R}$-modules $\widetilde f : \widetilde M \to \widetilde N$
  is a monomorphism.
  \begin{proof}
    \uses{def:tildeMod}
    Monomorphisms of sheaves of modules are detected on stalks, and a stalk map
    is detected on a basis of opens. Over a basic open $D(g)$ the section map of
    $\widetilde f$ is the localization-away-of-$g$ of $f$ (the canonical map
    $\widetilde{(-)} \to (-)_g$ is the universal localization at the powers of
    $g$), and localization preserves injectivity. Hence $\widetilde f$ is
    injective on each basic open, so injective on every stalk, so a monomorphism.
  \end{proof}
\end{lemma}

\begin{lemma}\leanok
[Tilde preserves monomorphisms (instance form)]
  \label{lem:tilde_preservesMono}
  \lean{MR0555258CompactifyingPicard.tilde_preservesMono}
  \uses{lem:tilde_mono}
  The functor $\widetilde{(-)} : \ModuleCat R \to (\Spec R)\text{-mod}$ has the
  \texttt{PreservesMonomorphisms} property; this is \cref{lem:tilde_mono}
  repackaged as a typeclass instance.
  \begin{proof}
    \uses{lem:tilde_mono}
    Immediate from \cref{lem:tilde_mono}: a functor preserves monomorphisms iff it
    carries each mono to a mono, which is exactly that statement.
  \end{proof}
\end{lemma}

\begin{lemma}\leanok
[Tilde preserves finite limits]
  \label{lem:tilde_preservesFiniteLimits}
  \lean{MR0555258CompactifyingPicard.tilde_preservesFiniteLimits}
  \uses{def:tildeMod, lem:tilde_mono}
  The functor $\widetilde{(-)} : \ModuleCat R \to (\Spec R)\text{-mod}$ preserves
  finite limits. (The instance \texttt{tilde\_preservesMono} packaging
  \cref{lem:tilde_mono} as a \texttt{PreservesMonomorphisms} typeclass is the
  load-bearing input.)
  \begin{proof}
    \uses{lem:tilde_mono}
    Tilde is additive and a left adjoint, hence preserves finite colimits and
    epimorphisms (it is right exact). For an additive functor, preserving finite
    limits is equivalent to sending every short exact sequence to a left-exact one
    together with preservation of monomorphisms; the monomorphism half is
    \cref{lem:tilde_mono} and the rest is the already-available right exactness.
  \end{proof}
\end{lemma}

\begin{lemma}\leanok
[Tilde--$\Gamma$ counit is an iso for a presented sheaf]
  \label{lem:tilde_counit_of_pres}
  \lean{MR0555258CompactifyingPicard.tilde_of_presentation}
  \uses{def:tildeMod}
  Let $M$ be an $\mathcal{O}_{\Spec R}$-module admitting a \emph{global}
  presentation $P$ (a free presentation by sheaves of modules on the trivial
  cover). Then the counit $\widetilde{\Gamma(\Spec R, M)} \to M$ of the
  tilde\,/\,global-sections adjunction is an isomorphism; equivalently
  $M \cong \widetilde{\Gamma(\Spec R, M)}$.
  \begin{proof}
    \uses{def:tildeMod}
    This is the essential-surjectivity half of the affine equivalence, supplied
    directly by Mathlib (\texttt{isIso\_fromTildeΓ\_of\_presentation}): a global
    presentation forces the counit to be an isomorphism, since both source and
    target then arise from the same generators-and-relations data.
  \end{proof}
\end{lemma}

\begin{theorem}[Affine quasi-coherence: a finitely presented sheaf on $\Spec R$ is a tilde of an fp module (Stacks)]
  \label{thm:stacks_affine_fp_tilde}
  \lean{MR0555258CompactifyingPicard.External.affine_fp_tilde}
  % NOTE (iter-015): External anchor. The essential-surjectivity half of the affine
  % equivalence QCoh(Spec R) ≃ Mod_R is ABSENT from Mathlib v4.30.0 (verified iter-014
  % four ways + iter-015 leansearch/loogle: tilde.functor has Full/Faithful/IsLeftAdjoint
  % but NO EssSurj/essImage, and SheafOfModules.IsFinitePresentation unfolds only to LOCAL
  % QuasicoherentData). It is elementary, buildable AG (Stacks 01IA + 01PC) — anchored here
  % to unblock the affine bridge after a 3-iter from-scratch-build attempt; flagged as a
  % Mathlib gap to build/upstream later (STRATEGY.md), and recommended protected (TO_USER).
  % SOURCE: [Stacks], Lemma 26.7.4 (Tag 01IA) + Lemma 28.17.2 (Tag 01PC) (read from references/stacks-qcoh-affine.md)
  % SOURCE QUOTE: "Let $(X, \mathcal{O}_X) = (\mathop{\mathrm{Spec}}(R), \mathcal{O}_{\mathop{\mathrm{Spec}}(R)})$
  % be an affine scheme. Let $\mathcal{F}$ be a quasi-coherent $\mathcal{O}_X$-module. Then
  % $\mathcal{F}$ is isomorphic to the sheaf associated to the $R$-module $\Gamma(X, \mathcal{F})$."
  % SOURCE QUOTE: "The quasi-coherent sheaf of $\mathcal{O}_X$-modules $\widetilde M$ is an
  % $\mathcal{O}_X$-module of finite presentation if and only if $M$ is an $R$-module of finite presentation."
  \textit{Source: Stacks Project, Lemma 26.7.4 (Tag 01IA) and Lemma 28.17.2 (Tag 01PC).}
  Let $X = \Spec R$ and let $I$ be a finitely presented (hence quasi-coherent)
  $\mathcal{O}_X$-module. Then there exists a finitely presented $R$-module $M$
  with $I \cong \widetilde M$.
  \begin{proof}
    By Stacks 01IA the quasi-coherent $I$ is isomorphic to
    $\widetilde{\Gamma(X, I)}$; take $M = \Gamma(X, I)$. By Stacks 01PC, since
    $\widetilde M \cong I$ is of finite presentation as an $\mathcal{O}_X$-module,
    $M$ is of finite presentation as an $R$-module. Used here as an external
    input (affine quasi-coherence equivalence).
  \end{proof}
\end{theorem}

\begin{lemma}\leanok
[A finitely presented sheaf on an affine scheme is a tilde]
  \label{lem:fp_sheaf_affine_tilde}
  \lean{MR0555258CompactifyingPicard.fp_sheaf_affine_tilde}
  \uses{def:tildeMod, def:isLFP, thm:stacks_affine_fp_tilde}
  % NOTE (iter-015): now a thin project-facing restatement of the external anchor
  % thm:stacks_affine_fp_tilde. The earlier sheaf-presentation route (lem:affine_global_presentation,
  % removed) was a detour: the 1.4 free-resolution step consumes the (module M fp, I ≅ ~M) form
  % directly, tilde-ing M's MODULE-level free presentation — not a sheaf-level I.Presentation.
  Let $X = \Spec R$ and $I$ a finitely presented (in particular quasi-coherent)
  $\mathcal{O}_X$-module. Then there is a finitely presented module $M$ over
  $\Gamma(X, \mathcal{O}_X) = R$ with $I \cong \widetilde M$.
  \begin{proof}
    \uses{thm:stacks_affine_fp_tilde}
    Immediate from the affine quasi-coherence anchor
    \cref{thm:stacks_affine_fp_tilde}: it provides a finitely presented
    $R$-module $M = \Gamma(X, I)$ together with an isomorphism
    $I \cong \widetilde M$.
  \end{proof}
\end{lemma}

% ---- coverage of prover-created helpers (iter-014 lean_aux) ----

\begin{lemma}

\leanok
[Mathlib's free sheaf of modules is free in the project sense]
  \label{lem:free_isFreeMod}
  \lean{MR0555258CompactifyingPicard.free_isFreeMod}
  \uses{def:isFreeMod}
  For an index type $\Lambda$, Mathlib's free sheaf of modules
  \texttt{SheafOfModules.free}~$\Lambda$ on a scheme $X$ satisfies the project
  predicate \cref{def:isFreeMod}.
  \begin{proof}
    \uses{def:isFreeMod}
    Immediate: \texttt{SheafOfModules.free}~$\Lambda$ is, by definition, the free
    $\mathcal{O}_X$-module on $\Lambda$, so the identity isomorphism witnesses
    \cref{def:isFreeMod} with index set $\Lambda$.
  \end{proof}
\end{lemma}

\begin{lemma}

\leanok
[A global presentation yields a free presentation in project terms]
  \label{lem:exists_free_presentation_of_presentation}
  \lean{MR0555258CompactifyingPicard.exists_free_presentation_of_presentation}
  \uses{def:isFreeMod, lem:free_isFreeMod}
  Let $M$ be an $\mathcal{O}_{\Spec R}$-module admitting a global presentation
  $P$ (a \texttt{SheafOfModules.Presentation}: generators and relations, both
  free sheaves). Then $M$ has a \emph{free presentation} in project terms: a
  free module $G$ with an epimorphism $g : G \twoheadrightarrow M$, together with
  a free module $K$ and an epimorphism $K \twoheadrightarrow \ker g$.
  \begin{proof}
    \uses{lem:free_isFreeMod}
    The presentation $P$ provides its generators as a free sheaf with an
    epimorphism $G = \mathrm{free}\,\Lambda \twoheadrightarrow M$ and its
    relations as a free sheaf with an epimorphism
    $K = \mathrm{free}\,\Lambda' \twoheadrightarrow \ker g$. Both covers are free
    by \cref{lem:free_isFreeMod}; this is exactly the asserted free presentation.
  \end{proof}
\end{lemma}

\begin{lemma}\leanok
[Tilde of a finitely presented module is a finitely presented sheaf (Stacks 01PC, forward)]
  \label{lem:tilde_isFP}
  \lean{MR0555258CompactifyingPicard.tilde_isFP}
  \uses{def:tildeMod, def:isLFP, lem:tilde_free, lem:tilde_exact}
  % NOTE (iter-016): the Lean proof does NOT follow the constructive
  % tilde_free+tilde_exact+IsFinitePresentation.mk route written below — it routes
  % through the new boundary axiom External.tilde_isFP (Stacks 01PC forward). The
  % constructive route needs (a) building SheafOfModules.Presentation(~M) from a
  % Module.FinitePresentation and (b) a Presentation.IsFinite ⇒ IsFinitePresentation
  % brick, both absent from Mathlib v4.30.0. Anchored per the iter-016 plan fallback.
  % SOURCE: [Stacks], Lemma 28.17.2 (Tag 01PC), forward direction (read from references/stacks-qcoh-affine.md)
  % SOURCE QUOTE: "The quasi-coherent sheaf of $\mathcal{O}_X$-modules $\widetilde M$ is an
  % $\mathcal{O}_X$-module of finite presentation if and only if $M$ is an $R$-module of finite presentation."
  \textit{Source: Stacks Project, Lemma 28.17.2 (Tag 01PC), forward direction.}
  For a finitely presented module $M$ over a commutative ring $R$, the tilde
  $\widetilde M$ is a finitely presented $\mathcal{O}_{\Spec R}$-module, i.e. it
  satisfies \cref{def:isLFP}.
  \begin{proof}
    \uses{lem:tilde_free, lem:tilde_exact}
    Since $M$ is finitely presented it admits a global free presentation
    $R^{(\Lambda')} \to R^{(\Lambda)} \to M \to 0$ with $\Lambda, \Lambda'$
    finite. Tilde is exact (\cref{lem:tilde_exact}) and carries finite free
    modules to finite free sheaves (\cref{lem:tilde_free}), so this tildes to a
    presentation $\mathcal{O}^{(\Lambda')} \to \mathcal{O}^{(\Lambda)} \to
    \widetilde M \to 0$ of $\widetilde M$ by finite free sheaves. Such a global
    presentation by finite free sheaves is exactly a finite-presentation datum
    for the sheaf of modules (the constructor
    \texttt{SheafOfModules.IsFinitePresentation.mk} from a quasi-coherent
    presentation on the trivial cover); hence $\widetilde M$ is finitely
    presented. This is the forward direction of Stacks 01PC (the affine anchor
    \cref{thm:stacks_affine_fp_tilde} supplies the reverse direction).
  \end{proof}
\end{lemma}

\begin{lemma}\leanok
[Finite free presentation of a finite module over a Noetherian ring]
  \label{lem:finiteFreePresentation_noetherian}
  \lean{MR0555258CompactifyingPicard.exists_finiteFreePresentation_of_noetherian}
  \uses{def:isFreeMod}
  Let $R$ be a Noetherian commutative ring and $M$ a finite $R$-module. Then
  there is a finite free module $R^n$ with a surjection $p : R^n \to M$ whose
  kernel $\ker p$ is again finitely presented (indeed finite, hence finitely
  presented because $R$ is Noetherian). This is the module-theoretic core of the
  free-resolution step \cref{lem:free_resolution_step}.
  \begin{proof}
    \uses{def:isFreeMod}
    A finite generating set of $M$ gives a surjection $p : R^n \to M$ from a
    finite free module. Its kernel is a submodule of $R^n$, finitely generated
    because $R$ is Noetherian; a finitely generated module over a Noetherian ring
    is finitely presented. Transporting through the categorical kernel
    isomorphism gives the stated finite presentation of $\ker p$.
  \end{proof}
\end{lemma}

\begin{lemma}\leanok
[Finite free sheaf of rank $n$ on $\Spec R$]
  \label{lem:tilde_free_fin}
  \lean{MR0555258CompactifyingPicard.tilde_free_fin}
  \uses{def:tildeMod, def:isFreeMod, lem:tilde_free}
  For a natural number $n$, the tilde $\widetilde{R^n} = \widetilde{(\Fin n \to_0
  R)}$ of the rank-$n$ finite free module is a free $\mathcal{O}_{\Spec R}$-module
  (\cref{def:isFreeMod}). This is the rank-$n$ ($\Fin n$-indexed) variant of
  \cref{lem:tilde_free}, used for the middle term $J_0 = \widetilde{A_0^n}$ of the
  free-resolution step.
  \begin{proof}
    \uses{lem:tilde_free}
    Reindex $\Fin n \to_0 R$ along the equivalence $\mathrm{ULift}(\Fin n) \simeq
    \Fin n$ to land the index in the universe used by \cref{lem:tilde_free}, apply
    tilde, and use that tilde carries a finite-rank free module to a free sheaf.
  \end{proof}
\end{lemma}

\begin{lemma}\leanok
[Pullback of a free sheaf is free]
  \label{lem:pullback_isFreeMod}
  \lean{MR0555258CompactifyingPicard.pullback_isFreeMod}
  \uses{def:isFreeMod, thm:pullback_preserves_free}
  For a scheme morphism $p : X \to Y$ and a free $\mathcal{O}_Y$-module $M$, the
  pullback $p^* M$ is a free $\mathcal{O}_X$-module (\cref{def:isFreeMod}). A thin
  restatement of the standard input \cref{thm:pullback_preserves_free}; used for the
  middle term $p^* J_0$ of the free-resolution step.
  \begin{proof}
    \uses{thm:pullback_preserves_free}
    Immediate from \cref{thm:pullback_preserves_free}.
  \end{proof}
\end{lemma}

\begin{lemma}\leanok
[Pullback of a finitely presented sheaf is finitely presented]
  \label{lem:pullback_isLFP}
  \lean{MR0555258CompactifyingPicard.pullback_isLFP}
  \uses{def:isLFP, thm:pullback_preserves_fp}
  For a scheme morphism $p : X \to Y$ and a finitely presented $\mathcal{O}_Y$-module
  $M$, the pullback $p^* M$ is finitely presented (\cref{def:isLFP}). A thin
  restatement of the standard input \cref{thm:pullback_preserves_fp}; used for the
  terms $p^* \widetilde{K_0}$ and $p^* J_0$ of the free-resolution step.
  \begin{proof}
    \uses{thm:pullback_preserves_fp}
    Immediate from \cref{thm:pullback_preserves_fp}.
  \end{proof}
\end{lemma}

%==============================================================================
% Strand (b): the relative local Ext sheaves, the base-change map, exchange,
% and the finiteness/flatness theorem.
%==============================================================================

\begin{lemma}\leanok
[Free resolution step over an affine base]
  \label{lem:free_resolution_step}
  \lean{MR0555258CompactifyingPicard.exists_free_resolution_step}
  \uses{thm:ega_descent_noetherian, def:isFreeMod, lem:tilde_free, lem:tilde_free_fin, lem:tilde_exact, lem:tilde_isFP, lem:fp_sheaf_affine_tilde, lem:finiteFreePresentation_noetherian, thm:flat_pullback_exact, lem:pullback_isFreeMod, lem:pullback_isLFP}
  % SOURCE: [Altman--Kleiman], §1, (1.4), p. 56--57 (read from references/MR0555258-compactifying-picard.pdf)
  % SOURCE QUOTE: "(1.4) LEMMA. Let f : X -> S be a finitely presented morphism
  % of affine schemes, and let I be an S-flat, finitely presented O_X-Module.
  % Then there exists an exact sequence 0 -> K -> J -> I -> 0, (1.4.1) with K and
  % J finitely presented and with J free."
  \textit{Source: Altman--Kleiman, §1, (1.4), p. 56--57.}
  Let $f : X \to S$ be a finitely presented morphism of affine schemes and let
  $I$ be an $S$-flat, finitely presented $\Ox$-module. Then there is an exact
  sequence
  \[
    0 \to K \to J \to I \to 0
    \tag{1.4.1}
  \]
  with $K$ and $J$ finitely presented and $J$ free.
  % SOURCE QUOTE PROOF: "Proof. By [EGA IV_3, Sect. 8], there exists a Noetherian
  % affine scheme such that all the data descend to S_0. Since X_0 is Noetherian,
  % there exists a sequence like (1.4.1) on X_0. (On X, we can construct such a
  % sequence with K finitely generated [CA, I, Sect. 2.8, Lemma 9, p. 21]; on X_0
  % any finitely generated Module is finitely presented.) Since I is flat, the
  % pullback of the sequence on X_0 is the desired sequence on X."
  \begin{proof}
    \uses{thm:ega_descent_noetherian, lem:tilde_free, lem:tilde_free_fin, lem:tilde_exact, lem:tilde_isFP, lem:fp_sheaf_affine_tilde, lem:finiteFreePresentation_noetherian, thm:flat_pullback_exact, lem:pullback_isFreeMod, lem:pullback_isLFP}
    \emph{Step 1 (descend to a Noetherian base).} By
    \cref{thm:ega_descent_noetherian} the data $(X, I)$ descend: there are a
    Noetherian ring $A_0$, a morphism $p : X \to X_0 = \Spec A_0$, and a finitely
    presented $\mathcal{O}_{X_0}$-module $I_0$ with $I \cong p^* I_0$.

    \emph{Step 2 (resolve over the Noetherian affine $X_0$).} Since $X_0 = \Spec
    A_0$ is affine, the affine bridge \cref{lem:fp_sheaf_affine_tilde} writes
    $I_0 \cong \widetilde{M_0}$ for a finitely presented $A_0$-module $M_0$. As
    $A_0$ is Noetherian, \cref{lem:finiteFreePresentation_noetherian} gives a
    finite free module $A_0^n$ with a surjection $A_0^n \twoheadrightarrow M_0$
    whose kernel $K_0$ is finitely presented, i.e. a short exact sequence
    $0 \to K_0 \to A_0^n \to M_0 \to 0$ of $A_0$-modules. Applying the exact
    functor tilde (\cref{lem:tilde_exact}) yields a short exact sequence
    $0 \to \widetilde{K_0} \to \widetilde{A_0^n} \to I_0 \to 0$ of
    $\mathcal{O}_{X_0}$-modules in which, by \cref{lem:tilde_free},
    $J_0 = \widetilde{A_0^n}$ is finite free, and by \cref{lem:tilde_isFP} both
    $\widetilde{K_0}$ and $J_0$ are finitely presented (the forward direction of
    Stacks 01PC, since $K_0$ and $A_0^n$ are finitely presented $A_0$-modules).

    \emph{Step 3 (pull back to $X$).} Apply the pullback functor $p^*$ to the
    sequence over $X_0$. Pullback sends the finite free sheaf $J_0$ to a finite
    free sheaf $p^* J_0$ (\cref{lem:pullback_isFreeMod}) and preserves finite
    presentation (\cref{lem:pullback_isLFP}), so $p^* \widetilde{K_0}$
    and $p^* J_0$ are finitely presented with $p^* J_0$ free. Because $I \cong
    p^* I_0$ is flat over $S$, \cref{thm:flat_pullback_exact} ensures the
    pulled-back sequence
    $0 \to p^* \widetilde{K_0} \to p^* J_0 \to p^* I_0 \to 0$ stays short exact.
    Transporting along $I \cong p^* I_0$ gives the desired sequence
    $0 \to K \to J \to I \to 0$ over $X$ with $J = p^* J_0$ free and $K, J$
    finitely presented.
  \end{proof}
\end{lemma}

\begin{lemma}[Local Ext commutes with limits]
  \label{lem:ext_limit}
  \lean{MR0555258CompactifyingPicard.ext_limit}
  % NOTE: deferred — Lean decl not yet scaffolded (TODO block in Basic.lean); needs projective limits of schemes and modules, absent from Mathlib.
  \uses{lem:free_resolution_step, thm:ega_basechange_q0}
  % SOURCE: [Altman--Kleiman], §1, (1.5), p. 57 (read from references/MR0555258-compactifying-picard.pdf)
  % SOURCE QUOTE: "(1.5) LEMMA. Let X_0 be an affine scheme, and let
  % S = lim S_lambda be a projective limit of S_0-schemes S. Let f_0 : X_0 -> S_0
  % be a finitely presented morphism, let I_0 and F_0 be locally finitely
  % presented O_{X_0}-Modules, and let X = lim(X_lambda), I = lim(I_lambda), and
  % F = lim(F_lambda) be the natural limits induced. Fix an integer q. Assume
  % that I_0 is S_0-flat if q >= 1 and that X_0 is S_0-flat if q >= 2. Then there
  % is a canonical isomorphism, lim Ext^q_{X_lambda}(I_lambda, F_lambda) =
  % Ext^q_X(I, F)."
  \textit{Source: Altman--Kleiman, §1, (1.5), p. 57.}
  Let $X_0$ be affine and $S = \varprojlim S_\lambda$ a projective limit of
  $S_0$-schemes. Let $f_0 : X_0 \to S_0$ be finitely presented, let $I_0, F_0$ be
  locally finitely presented $\mathcal{O}_{X_0}$-modules, and form the induced
  limits $X = \varprojlim X_\lambda$, $I = \varprojlim I_\lambda$,
  $F = \varprojlim F_\lambda$. Fix $q$, and assume $I_0$ is $S_0$-flat if
  $q \geq 1$ and $X_0$ is $S_0$-flat if $q \geq 2$. Then there is a canonical
  isomorphism
  \[
    \varprojlim \sExt^q_{X_\lambda}(I_\lambda, F_\lambda)
      \;=\; \sExt^q_X(I,F).
  \]
  % SOURCE QUOTE PROOF: "Proof. The assertion is local, so we may assume X_0 and
  % the S_lambda are affine. The proof now proceeds by induction on q >= 0. For
  % q = 0, the assertion results from [EGA IV_3, 8.5.2 (i)]. Consider the case
  % q = 1. By (1.4) there exists on X_0 an exact sequence, 0 -> K_0 -> J_0 -> I_0
  % -> 0, (1.5.1) with J_0 and K_0 finitely presented and with J_0 free. [...]
  % Consider the case q >= 2. The sequence (1.5.1) yields diagrams, [...] Since
  % X_0 is S_0-flat, the free O_{X_0}-Module J_0 is also S_0-flat. Hence since I_0
  % is S_0-flat, K_0 is also. Therefore, by induction on q, the left-hand
  % vertical maps in (1.5.3) induce an isomorphism in the limit, Hence, so do the
  % right-hand vertical maps. The dotted arrows in (1.5.2) and (1.5.3) do not
  % depend on the choice of the exact sequence (1.5.1) because any two such
  % sequences are homotopic since J_0 is free."
  \begin{proof}
    \uses{lem:free_resolution_step, thm:ega_basechange_q0}
    The assertion is local, so we may assume the $S_\lambda$ affine, and we
    induct on $q$. For $q = 0$ it is \cref{thm:ega_basechange_q0}. For $q=1$,
    \cref{lem:free_resolution_step} gives an exact sequence
    $0 \to K_0 \to J_0 \to I_0 \to 0$ on $X_0$ with $J_0$ free; since $I_0$ is
    $S_0$-flat this stays exact after base change, producing compatible exact
    $\Hom$-to-$\sExt^1$ sequences whose $q=0$ terms agree in the limit by the
    previous case, forcing the $\sExt^1$ terms to agree. For $q \geq 2$, freeness
    of $J_0$ makes $\sExt^{\geq 1}(J_0, -)$ vanish, so the connecting maps give
    $\sExt^{q-1}(K_\lambda, -) \cong \sExt^q(I_\lambda, -)$; flatness of $X_0$
    makes $J_0$, hence $K_0$, flat, and the inductive hypothesis applied to $K$
    in degree $q-1$ gives the isomorphism in degree $q$. The induced maps are
    independent of the chosen sequence since any two free presentations are
    homotopic.
  \end{proof}
\end{lemma}

\begin{lemma}[Affine base change for Ext-modules]
  \label{lem:ext_algebra_basechange}
  \lean{MR0555258CompactifyingPicard.ext_basechange_algebra}
  % NOTE: deferred — Lean decl not yet scaffolded (TODO block in Basic.lean); the affine/algebraic incarnation, needs the comparison over Spec B, absent from Mathlib.
  \uses{lem:ext_limit}
  % SOURCE: [Altman--Kleiman], §1, (1.6), p. 58 (read from references/MR0555258-compactifying-picard.pdf)
  % SOURCE QUOTE: "(1.6) LEMMA. Let A be a ring, let B a finitely presented
  % A-algebra, and let M and N be finitely presented B-modules. Fix an integer q.
  % Assume that M is A-flat if q >= 1 and that B is A-flat if q >= 2. Then there
  % is a canonical isomorphism, Ext^q_A(M, N)~ = Ext^q_B(M~, N~)."
  \textit{Source: Altman--Kleiman, §1, (1.6), p. 58.}
  Let $A$ be a ring, $B$ a finitely presented $A$-algebra, and $M, N$ finitely
  presented $B$-modules. Fix $q$, and assume $M$ is $A$-flat if $q \geq 1$ and $B$
  is $A$-flat if $q \geq 2$. Then there is a canonical isomorphism of sheaves on
  $\Spec B$,
  \[
    \widetilde{\Ext^q_A(M, N)} \;=\; \sExt^q_B(\widetilde M, \widetilde N).
  \]
  % SOURCE QUOTE PROOF: "Proof. The case q = 0 is proved in [EGA I, 1.3.12,
  % (ii)]. The rest of the proof is straightforward and similar to that of the
  % preceding lemma."
  \begin{proof}
    \uses{lem:ext_limit}
    The case $q = 0$ is [EGA I, 1.3.12(ii)]. The general case is proved exactly
    as \cref{lem:ext_limit}: take a finite free presentation of $M$ as a
    $B$-module and induct on $q$, using the flatness hypotheses to keep the
    syzygies flat.
  \end{proof}
\end{lemma}

\begin{lemma}[Adjunction for local Ext under base change]
  \label{lem:ext_adjunction}
  \lean{MR0555258CompactifyingPicard.ext_adjunction}
  % NOTE: deferred — Lean decl not yet scaffolded (TODO block in Basic.lean); needs fibre products X_T and pushforward along 1×g, absent from Mathlib.
  \uses{thm:ega_adjunction, lem:ext_limit}
  % SOURCE: [Altman--Kleiman], §1, (1.7), p. 58 (read from references/MR0555258-compactifying-picard.pdf)
  % SOURCE QUOTE: "(1.7) LEMMA. Let f : X -> S be a finitely presented morphism
  % of schemes, and let I and F be locally finitely presented O_X-Modules. Fix an
  % integer q. Assume that I is S-flat if q >= 1 and that f is S-flat if q >= 2.
  % Then there exists, for each base-change morphism g : T -> S and each
  % quasi-coherent O_T-Module M, a canonical "adjunction" isomorphism,
  % Ext^q_X(I, (1 x g)_*(F (x)_S M)) ~-> (1 x g)_* Ext^q_{X_T}(I_T, F (x)_S M).
  % (1.7.1)"
  \textit{Source: Altman--Kleiman, §1, (1.7), p. 58.}
  Let $f : X \to S$ be finitely presented and $I, F$ locally finitely presented
  $\Ox$-modules. Fix $q$, and assume $I$ is $S$-flat if $q \geq 1$ and $f$ is flat
  if $q \geq 2$. Then for each base change $g : T \to S$ and each quasi-coherent
  $\Ot$-module $M$ there is a canonical adjunction isomorphism
  \[
    \sExt^q_X\!\big(I, (1\times g)_*(F \tensor_S M)\big)
      \;\xrightarrow{\;\sim\;}\;
      (1\times g)_* \sExt^q_{X_T}\!\big(I_T, F \tensor_S M\big),
      \tag{1.7.1}
  \]
  compatible with further base change and with passage to limits as in
  \cref{lem:ext_limit}.
  % SOURCE QUOTE PROOF: "Proof. For q = 0, the isomorphism (1.7.1) comes from the
  % usual adjunction isomorphism [EGA O_I, 4.4.3.1]. The compatibilities are
  % straightforward. For general q, the construction is straightforward,
  % following the line of reasoning of (1.5). The compatibilities follow,
  % similarly, from those for q = 0."
  \begin{proof}
    \uses{thm:ega_adjunction, lem:ext_limit}
    For $q = 0$ this is the usual adjunction isomorphism
    (\cref{thm:ega_adjunction}). For general $q$ one builds the map by the method
    of \cref{lem:ext_limit} (a free presentation of $I$ and induction on $q$),
    and the compatibilities reduce to those for $q = 0$.
  \end{proof}
\end{lemma}

\begin{definition}[The base-change map for local Ext's]
  \label{def:extBaseChangeMap}
  \lean{MR0555258CompactifyingPicard.extBaseChangeMap}
  % NOTE: deferred — Lean decl not yet scaffolded (TODO block in Basic.lean); needs X_T, the canonical unit/counit maps and the adjoint of (1.7.1), absent from Mathlib.
  \uses{lem:ext_adjunction, lem:ext_algebra_basechange}
  % SOURCE: [Altman--Kleiman], §1, (1.8), p. 58--59 (read from references/MR0555258-compactifying-picard.pdf)
  % SOURCE QUOTE: "(1.8) (The base-change map for local Ext's). Let f : X -> S be
  % a finitely presented morphism of schemes, and let I and F be locally finitely
  % presented O_X-Modules. Let g : T -> S be a morphism, and let M be a
  % quasi-coherent O_T-Module. The canonical map, F -> (1 x g)_*(1 x g)*F,
  % induces a map, Ext^q_X(I, F) (x)_S M -> Ext^q_X(I, (1 x g)_*(1 x g)*F) (x)_S
  % M. (1.8.1) [...] Assume that I is S-flat if q >= 1 and that f is flat if
  % q >= 2. Then composing (1.8.1) and (1.8.2) with the adjoint of (1.7.1) yields
  % a canonical base-change map, b^q(M) : Ext^q_X(I, F) (x)_S M ->
  % Ext^q_{X_T}(I_T, F (x)_S M). [...] If the base-change g : T -> S is flat, then
  % the base-change map b^q(O_T) is an isomorphism."
  \textit{Source: Altman--Kleiman, §1, (1.8), p. 58--59.}
  Let $f : X \to S$ be finitely presented, $I, F$ locally finitely presented
  $\Ox$-modules, $g : T \to S$ a morphism, and $M$ a quasi-coherent $\Ot$-module.
  The canonical map $F \to (1\times g)_*(1\times g)^* F$ induces
  \[
    \sExt^q_X(I,F) \tensor_S M \to
      \sExt^q_X\!\big(I, (1\times g)_*(1\times g)^* F\big) \tensor_S M,
      \tag{1.8.1}
  \]
  and writing out the canonical map $R(\Ot) \tensor_S M \to R(M)$ with
  $R(M) = (1\times g)^* \sExt^q_X(I, (1\times g)_*(F \tensor_S M))$ gives
  \[
    \sExt^q_X\!\big(I, (1\times g)_*(1\times g)^* F\big) \tensor_S M \to
      (1\times g)^* \sExt^q_X\!\big(I, (1\times g)_*(F \tensor_S M)\big).
      \tag{1.8.2}
  \]
  Assuming $I$ is $S$-flat for $q \geq 1$ and $f$ flat for $q \geq 2$, composing
  \textnormal{(1.8.1)} and \textnormal{(1.8.2)} with the adjoint of
  \textnormal{(1.7.1)} (\cref{lem:ext_adjunction}) defines the canonical
  base-change map
  \[
    b^q(M) : \sExt^q_X(I,F) \tensor_S M \to
      \sExt^q_{X_T}\!\big(I_T, F \tensor_S M\big).
  \]
  It commutes with restriction to open subschemes, to $\Spec(\Os)$, and to other
  localizations of $S$, is compatible with further base change and limits as in
  \cref{lem:ext_limit}, and when $g$ is flat, $b^q(\Ot)$ is an isomorphism.
  \begin{proof}
    \uses{lem:ext_adjunction, lem:ext_algebra_basechange}
    The map is the composite \textnormal{(1.8.1)}, \textnormal{(1.8.2)}, adjoint
    of \textnormal{(1.7.1)}, well defined under the stated flatness hypotheses.
    For flat $g$ the isomorphism $b^q(\Ot)$ is local on $S, X, T$ and holds in
    the affine case by \cref{lem:ext_algebra_basechange} and flat base change for
    modules.
  \end{proof}
\end{definition}

\begin{theorem}[Property of exchange for local Ext's]
  \label{thm:exchange}
  \lean{MR0555258CompactifyingPicard.exchange}
  % NOTE: deferred — Lean decl not yet scaffolded (TODO block in Basic.lean); built on extBaseChangeMap + the Nakayama iso, both deferred (infra absent from Mathlib).
  \uses{def:extBaseChangeMap, lem:ext_limit, thm:ob_nakayama_iso, thm:ob_halfexact_free}
  % SOURCE: [Altman--Kleiman], §1, (1.9), p. 59 (read from references/MR0555258-compactifying-picard.pdf)
  % SOURCE QUOTE: "(1.9) THEOREM (property of exchange for local Ext's). Let
  % f : X -> S be a finitely presented morphism of schemes, and let I and F be
  % locally finitely presented O_T-Modules. Assume F is S-flat. Fix an integer q.
  % Assume (a) I is S-flat if q >= 1 and (b) f is flat if q >= 2. Fix a point s in
  % S and a point x in X(s). Assume that the base-change map to the fiber,
  % b^q(k(s)) : Ext^q_X(I, F) (x)_S k(s) -> Ext^q_{X(s)}(I(s), F(s)), is
  % surjective at x. Then, (i) For every map g : T -> S and every quasi-coherent
  % O_T-Module M, the base-change map b^q(M) is an isomorphism at every point of
  % (1 x g)^{-1}(x). (ii) The following three statements are equivalent: (1)
  % b^{q-1}(k(s)) is surjective at x. (2) b^{q-1}(M) is an isomorphism at every
  % point of (1 x g)^{-1}(x) for every g and every M. (3) Ext^q_X(I, F) is S-flat
  % at x."
  \textit{Source: Altman--Kleiman, §1, (1.9), p. 59--61.}
  Let $f : X \to S$ be finitely presented and $I, F$ locally finitely presented
  $\Ox$-modules with $F$ $S$-flat. Fix $q$ and assume $I$ is $S$-flat if
  $q \geq 1$ and $f$ flat if $q \geq 2$. Fix $s \in S$ and $x \in X(s)$, and
  assume the base-change map to the fibre
  \[
    b^q(k(s)) : \sExt^q_X(I,F) \tensor_S k(s) \to
      \sExt^q_{X(s)}\!\big(I(s), F(s)\big)
  \]
  is surjective at $x$. Then:
  \begin{enumerate}
    \item[(i)] for every $g : T \to S$ and every quasi-coherent $\Ot$-module $M$,
      $b^q(M)$ is an isomorphism at every point of $(1\times g)^{-1}(x)$;
    \item[(ii)] the following are equivalent: \textnormal{(1)} $b^{q-1}(k(s))$ is
      surjective at $x$; \textnormal{(2)} $b^{q-1}(M)$ is an isomorphism at every
      point of $(1\times g)^{-1}(x)$ for all $g, M$; \textnormal{(3)}
      $\sExt^q_X(I,F)$ is $S$-flat at $x$.
  \end{enumerate}
  % SOURCE QUOTE PROOF: "Proof. (i) Clearly we may assume S and X are affine.
  % Write S as a limit, S = lim S_lambda, where each S_lambda is the spectrum of
  % a finitely generated Z-algebra. [...] Define a functor from the category of
  % O_s-modules N to the category of O_s-modules, R(N) = Ext^q_{O_x}(I_x, F_x
  % (x)_{O_x} N). It is easy to see that R commutes with direct limits. Moreover,
  % if N is finitely generated, then R(N) is also finitely generated [GD IV, 3.2
  % (i), p. 74]. Since b^q(k(s))_x is surjective, the natural map, R(O_s) (x)_{O_s}
  % k(s) -> R(k(s)), is surjective. Moreover, the (unique) maximal ideal of O_x
  % contracts to the (unique) maximal ideal of O_s. Therefore, by [OB, 4.1], the
  % map, R(O_s) (x)_{O_s} N -> R(N), (1.9.1) [...] is an isomorphism. [...]
  % (ii) The implication (1) => (2) holds by (i). For the implications (2) => (3)
  % and (3) => (1), clearly we may assume S = Spec(O_s). Let 0 -> M' -> M -> M'' ->
  % 0 be an arbitrary exact sequence of quasi-coherent O_S-Modules, and consider
  % the following diagram, with two commutative squares and exact lower sequence:
  % [...] The maps b^q(M') and b^q(M) are isomorphisms at x by (i)."
  \begin{proof}
    \uses{def:extBaseChangeMap, lem:ext_limit, thm:ob_nakayama_iso, thm:ob_halfexact_free, lem:ext_adjunction}
    \emph{(i)} We may assume $S, X$ affine and, writing $S = \varprojlim
    S_\lambda$ over finitely generated $\mathbf{Z}$-algebras, descend the data; by
    \cref{lem:ext_limit} the fibre maps $b^q(k(s_\lambda))$ have limit
    $b^q(k(s))$, so the hypotheses descend and we may take $S$ Noetherian. Set
    $R(N) = \Ext^q_{\mathcal{O}_x}(I_x, F_x \tensor_{\mathcal{O}_x} N)$ on
    $\mathcal{O}_s$-modules; $R$ commutes with direct limits and preserves finite
    generation. Surjectivity of $b^q(k(s))$ at $x$ makes
    $R(\mathcal{O}_s) \tensor k(s) \to R(k(s))$ surjective, and since the maximal
    ideal of $\mathcal{O}_x$ contracts to that of $\mathcal{O}_s$,
    \cref{thm:ob_nakayama_iso} gives that $R(\mathcal{O}_s)\tensor N \to R(N)$ is
    an isomorphism for all $N$. Taking $N = M_t$ and comparing with the stalk of
    the adjunction isomorphism \textnormal{(1.7.1)} shows $b^q(M)$ is an
    isomorphism at $x$, hence at every point of $(1\times g)^{-1}(x)$.

    \emph{(ii)} The implication $(1)\Rightarrow(2)$ is part (i) applied in degree
    $q-1$. For the rest assume $S = \Spec(\mathcal{O}_s)$ and apply $b^\bullet$ to
    an exact sequence $0 \to M' \to M \to M'' \to 0$: by (i) the maps $b^q(M')$
    and $b^q(M)$ are isomorphisms at $x$, and a diagram chase using
    right-exactness of $\tensor$ yields $(2)\Rightarrow(3)$ (flatness of
    $\sExt^q_X(I,F)$ at $x$) and $(3)\Rightarrow(1)$ (taking $M = \mathcal{O}_s$,
    $M'' = k(s)$). The half-exactness input of \cref{thm:ob_halfexact_free}
    underlies the identification of the relevant functors.
  \end{proof}
\end{theorem}

\begin{theorem}\leanok
[Finiteness and flatness of local Ext]
  \label{thm:ext_finite_flat}
  \lean{MR0555258CompactifyingPicard.ext_finite_flat}
  \uses{thm:exchange, thm:ega_ext_vanishing_open, thm:ega_ext_coherent_fp}
  % SOURCE: [Altman--Kleiman], §1, (1.10), p. 61 (read from references/MR0555258-compactifying-picard.pdf)
  % SOURCE QUOTE: "(1.10) THEOREM. Let f : X -> S be a finitely presented, proper
  % morphism of schemes, and let I and F be locally finitely presented
  % O_X-Modules. Assume F is S-flat. Fix an integer q. Assume that I is S-flat if
  % q >= 1 and that f is flat if q >= 2. Then, (i) Let V denote the set of s in S
  % where we have Ext^q_{X(s)}(I(s), F(s)) = 0. Then V is open and retrocompact,
  % and for each base-change g : T -> S factoring through V and for each
  % quasi-coherent O_T-Module M, we have Ext^q_{X_T}(I_T, F (x)_S M) = 0. (ii) Fix
  % an integer c and let V denote the set of s in S, where we have
  % Ext^p_{X(s)}(I(s), F(s)) = 0 for p = c + 1, c - 1. Then V is open and
  % retrocompact, the restriction, Ext^c_X(I, F) | f^{-1}(V), is locally finitely
  % presented and flat over V, and the map b^p(M) is an isomorphism for every
  % base-change g : T -> S and for every quasi-coherent O_T-Module M."
  \textit{Source: Altman--Kleiman, §1, (1.10), p. 61.}
  Let $f : X \to S$ be finitely presented and proper, and $I, F$ locally finitely
  presented $\Ox$-modules with $F$ $S$-flat. Fix $q$ and assume $I$ is $S$-flat if
  $q \geq 1$ and $f$ flat if $q \geq 2$. Then:
  \begin{enumerate}
    \item[(i)] the set $V = \{ s \in S : \sExt^q_{X(s)}(I(s),F(s)) = 0 \}$ is open
      and retrocompact, and for each base change $g : T \to S$ factoring through
      $V$ and each quasi-coherent $\Ot$-module $M$ one has
      $\sExt^q_{X_T}(I_T, F \tensor_S M) = 0$;
    \item[(ii)] fixing $c$, the set $V$ of $s$ with
      $\sExt^p_{X(s)}(I(s),F(s)) = 0$ for $p = c+1, c-1$ is open and retrocompact;
      on it the restriction $\sExt^c_X(I,F)|f^{-1}(V)$ is locally finitely
      presented and flat over $V$, and $b^c(M)$ is an isomorphism for every base
      change $g : T \to S$ and every quasi-coherent $\Ot$-module $M$.
  \end{enumerate}
  % SOURCE QUOTE PROOF: "Proof. (i) Let U denote the set of x in X, where we have
  % Ext^q_{X(f(x))}(I(f(x)), F(f(x))) = 0. Then the proof of [EGA IV_3, 12.3.4]
  % shows that U is open and retrocompact, although this is not fully stated. [...]
  % It is easy to prove that, because f is proper and finitely presented, the set
  % V of points s such that f^{-1}(x) lies in U is open and retrocompact. The
  % assertion now follows from (1.9(i)) or from (1.10.1). (ii) By (i) the set V is
  % open and retrocompact. By (1.9(ii)) applied twice, first with q = c + 1 and
  % then with q = c, the restricted Ext is flat over V and the map b^q(M) is an
  % isomorphism for every g factoring through V and for every M. Finally, the
  % assertion of local finite presentation is local on S and is compatible with
  % base-change. So we may assume S is affine and by [EGA IV_3, Sects. 8, 11]
  % Noetherian. Then the assertion holds by [EGA III_2, 12.3.3]."
  \begin{proof}
    \uses{thm:exchange, thm:ega_ext_vanishing_open, thm:ega_ext_coherent_fp}
    \emph{(i)} Let $U = \{ x \in X : \sExt^q_{X(f(x))}(I(f(x)),F(f(x))) = 0 \}$;
    the argument of [EGA IV$_3$, 12.3.4] shows $U$ is open and retrocompact. As
    $f$ is proper and finitely presented, the set $V$ of $s$ with $f^{-1}(s)
    \subseteq U$ is open and retrocompact, and the vanishing statement follows
    from \cref{thm:exchange}(i).

    \emph{(ii)} By (i), $V$ is open and retrocompact. Applying
    \cref{thm:exchange}(ii) twice — first in degree $c+1$, then in degree $c$ —
    shows $\sExt^c_X(I,F)$ is flat over $V$ and $b^c(M)$ is an isomorphism for all
    $g$ factoring through $V$ and all $M$. Local finite presentation is local on
    $S$ and base-change compatible, so one reduces to $S$ affine Noetherian, where
    it is [EGA III$_2$, 12.3.3].
  \end{proof}
\end{theorem}
